\documentclass[12pt]{article}
% %%%%%%%%%%%%%%%%%%%%%%%%%%%%%%%%%%%%%%%%%%%%%%%%%%%%%%%%%%%%%%%%%%%%%%%%%%%%%%%%%%%%%%%%%%%%%%%%%%%%%%%%%%%%%%%%%%%%%%%%%%%%GGsection
%
\usepackage{amsmath,amssymb,amsthm,enumerate,graphicx}
\usepackage{ifthen,latexsym,syntonly}
\usepackage{setspace}
\usepackage{subcaption,caption}
\usepackage{mathabx}
%\usepackage[showrefs]{refcheck}  %use this to show equation and section labels       
\usepackage{color}
\usepackage[round,comma,authoryear]{natbib}   % for natbib
%\usepackage{subfigure}
 %\usepackage{graphics}
\usepackage{hyperref}
\usepackage{float}
\usepackage{epstopdf}
\usepackage{mathabx}
%\bibliographystyle{mynat}
%\bibliographystyle{ormsv080}

\usepackage{psfrag}
\usepackage{tikz}
\usetikzlibrary{decorations.pathreplacing}


%\usepackage{mathtools}
\usepackage{multirow}
%\usepackage{soul}


\onehalfspacing

\DeclareMathOperator*{\argmax}{arg\,max}
\DeclareMathOperator*{\argmin}{arg\,min}
\setlength{\textwidth}{6.5in}
\setlength{\textheight}{8.5in}
\usepackage{color}
\allowdisplaybreaks[3]
\usepackage{calc,xspace}
\definecolor{wacky}{rgb}{1,0,1}
\onehalfspacing

\newcommand{\wacky}[1]{{\large\textcolor{wacky}{#1}} }

\definecolor{ForestGreen}{RGB}{34, 139, 34}


\newtheorem{theorem}{Theorem}%[section]
\newtheorem{remark}[theorem]{Remark}
\newtheorem{assumption}{Assumption}
\newtheorem{case}[theorem]{Case}
\newtheorem{claim}[theorem]{Claim}
\newtheorem{corollary}{Corollary}
\newtheorem{condition}{Condition}
\newtheorem{criterion}[theorem]{Criterion}
\newtheorem{definition}{Definition}
\newtheorem{example}{Example}
\newtheorem{lemma}[theorem]{Lemma}
\newtheorem{problem}[theorem]{Problem}
\newtheorem{bound}[theorem]{Bound}
\newtheorem{proposition}{Proposition}
%\newtheorem{solution}[theorem]{Solution}
\newtheorem{thm}[theorem]{Theorem}
\newtheorem{restructure}{Debt Restructure}

\providecommand{\keywords}[1]{\textbf{\textit{Keywords---}} #1}


\setlength{\oddsidemargin}{.05in} \setlength{\topmargin}{-.45in}
\setlength{\textwidth}{6.4in} \setlength{\textheight}{8.5in}



\newcommand{\NW}[1]{{\textcolor{blue}{#1}}}
\newcommand{\WJ}[1]{{\textcolor{red}{#1}}}
\newcommand{\TS}[1]{{\textcolor{magenta}{#1}}}
%\newcommand{\JQ}[1]{{\textcolor{cyan}{#1}}}



\begin{document}

\date{%December 2009, revised
%TCIMACRO{\TeXButton{Today}{\today}}%
%BeginExpansion
\today%
%EndExpansion
}
\title{Implementing a Ramsey Tax Plan in a Stochastic Continuous Time Economy\thanks{We thank XXX and YYY for helpful comments.}}
\date{\today }
\author{	 Wei Jiang\thanks{%
Department of Industrial Engineering and Decision Analytics, Hong Kong University of Science and Technology. E-mail: weijiang@ust.hk } \and Thomas J. Sargent\thanks{New York University. E-mail: thomas.sargent@nyu.edu} \and Neng Wang\thanks{Cheung Kong Graduate School of Business (CKGSB). E-mail: {newang@gmail.com}} }

\maketitle


\begin{abstract}
\noindent
We extend a deterministic continuous-time formulation of a \cite{LucasStokey1983} optimal tax problem to a stochastic setting driven by a Markov chain. To render markets dynamically complete, we equip the government with locally risk-free instantaneous debt and  state-contingent hedging contracts. The government finances cumulative liabilities by rolling over instantaneous debt and managing a hedging portfolio, without rescheduling inherited term debts. A local commitment condition that preserves the purchasing power of instantaneous debt  induces continuation governments to confirm the Ramsey plan.  The value function is additively separable in a continuous state variable  that equals the product of cumulative liabilities and the household's marginal utility of consumption. This state variable  evolves continuously even when Markov transitions cause jumps in bond prices and marginal utilities.

\medskip

\thispagestyle{empty}

\noindent  \textbf{Keywords}: Ramsey planner, continuation Ramsey planner, implementability, instantaneous debt, local commitment, dynamic programming, hedging, complete markets.
\newline
\newline
JEL Classification: E62, H21, H63
\end{abstract} 
%
%
%\WJ{$\heartsuit$}
%$\spadesuit$ 
%
%\WJ{$\diamondsuit$}
%$\clubsuit$

\clearpage
\pagenumbering{arabic}
\setcounter{page}{1}
\newpage

\section{Introduction}


\citet{LucasStokey1983} studied optimal fiscal policy in a dynamic, stochastic, complete-markets economy. A  government finances an exogenous stochastic process for government expenditures $\{G_t\}$ and services  an inherited term debt coupon process $\{b_{0-,t}\}$ by choosing a  process of distorting flat rate taxes. \citeauthor{LucasStokey1983} characterized a Ramsey plan that maximizes the representative household's welfare and showed how to implement it by appropriately rescheduling the  term structure of government debt, thereby inducing  a sequence of continuation governments  to confirm the Ramsey tax plan. \citet*{debortoli2021optimal} showed that \citeauthor{LucasStokey1983}'s implementation can fail when initial debt is so high that the Ramsey planner sets the tax rate above the peak of the Laffer curve.

\NW{\cite{JiangSargentWang2024}, a non-stochastic predecessor of this paper,}  showed  that this difficulty can be resolved by expanding the government's contractible subspace to include  ``instantaneous debt'' --- a locally risk-free security of infinitesimal maturity --- together with a local commitment condition that requires each continuation government to preserve the ``purchasing power'' of outstanding instantaneous debt, measured as the product of the debt balance  and the representative household's marginal utility of consumption.  The deterministic setting of \cite{JiangSargentWang2024} featured  government expenditure and initial term debt service processes that were known functions of calendar time, so  that  bond prices were known at time~$0$ and  no hedging was needed. 

This paper extends   \cite{JiangSargentWang2024} to a stochastic setting  while  remaining faithful to   \citeauthor{LucasStokey1983}'s (1983) specification of complete markets. \NW{In continuous time, this requires some care.} \citeauthor{LucasStokey1983}'s (1983) original discrete-time  complete markets assumption, namely  that Arrow securities for every   state and date  are traded at time~$0$, has \NW{a  well-known continuous-time counterpart in  finance:} the dynamic   spanning constructions  of \cite{BlackScholes1973}, \cite{merton1973theory}, and \cite{HarrisonandKreps1979} that use  instantaneous debt  together with  competitively priced hedging contracts that pay off at state transitions. Our stochastic environment features a time-homogeneous Markov chain $\{s_t\}$ with $n$ states; starting from state~$s_{t-}$, there are $n-1$ possible transitions,  each associated with a hedging contract. Together, the instantaneous (locally risk-free)  debt and these $n-1$~hedging contracts  span all contingencies  over each infinitesimal interval, rendering  markets dynamically complete.   In contrast with the deterministic predecessor, the government must  actively manage both instantaneous  debt and a portfolio of hedging contracts  in order to  finance its  cumulative liabilities under the Ramsey plan. \NW{Our  main  results carry over from the deterministic setting, but the proofs are   more involved because  state transitions induce jumps in bond prices, marginal utilities, and the cumulative liability process.
} 



%We now briefly summarize   sections~\ref{sec:environ}--\ref{sec:indeterminacy}.

Section~\ref{sec:environ} describes the environment: a continuous-time economy with a representative household, a benevolent government, and Markovian exogenous processes $\{G_t,\, b_{0-,t}\}$.  The household's first-order conditions pin down the stochastic discount factor and establish a key  implementability constraint that encodes the government's present-value budget constraint purely in terms of allocations.  


Section~\ref{sec:Lagrangian1} formulates a Lagrangian for the Ramsey problem, attaching a multiplier~$\Lambda_0$ to the implementability constraint. Under the Markov assumption, optimal consumption $C(s_t;\Lambda_0)$ depends only on the current state  and  the  multiplier. This section also derives the equilibrium risk-free interest rate, which equals the household's subjective discount rate adjusted for risk premia associated with state transitions, and the market price of  state-transition risk.

Section~\ref{sec:bellmanramsey} reformulates the Ramsey problem as a dynamic program. We define the cumulative government liability $\Pi_t = \int_0^t \frac{q_{0,u}}{q_{0,t}} (b_{0-,u}+G_u-\tau_{0,u}N_{0,u})\, du$, which tracks what the government owes absent  any rescheduling of its  initial term debt.  While~$\Pi_t$ can  jump when the Markov state transitions, the product $X_t = \Pi_t\, U_{C,t}$ --- the ``purchasing power'' of cumulative liabilities --- evolves continuously in time. This continuity makes $X_t$ a natural state variable. We show that the value function takes the additively separable form  $V(X_t, s_t) = -\Phi_0\, X_t + H(s_t;\, \Phi_0)$, where the scalar~$\Phi_0$ is chosen by the Ramsey planner to maximize $H(s_0;\, \Phi_0)$ subject to the implementability constraint. The   Lagrangian multiplier $\Lambda_0^\ast$ from section~\ref{sec:Lagrangian1} equals $\Phi_0^\ast$, establishing  that the Lagrangian and Bellman equation approaches   yield the same Ramsey plan.


Section~\ref{sec:financing} shows how  the government finances its  cumulative liabilities. We extend \citeauthor{LucasStokey1983}'s complete markets assumption to our continuous-time stochastic setting by allowing the government to trade locally risk-free instantaneous debt and $n-1$~competitively priced hedging contracts. The government finances the Ramsey plan  by leaving the initial term debt structure $\{b_{0-,t}\}$ untouched and   adjusting  only its  instantaneous debt balance and hedging positions. The   instantaneous debt balance $B_t$ then  equals the cumulative liability $\Pi_t^\ast$, while the hedging   demand $\Delta_t^{(s_{t-},j)} = (1-e^{-\eta^\ast(s_{t-},j)})\Pi_{t-}^\ast$ offsets    state-transition risk.

Section~\ref{implementation} addresses  the  question that motivated \cite{LucasStokey1983} and \cite*{debortoli2021optimal}: how
to implement the Ramsey plan sequentially under a timing protocol that grants ``partial commitment'' to continuation governments. We adopt  a  local commitment condition:  each government's current-period tax rate is set by its immediate predecessor. Under this timing protocol, the continuation government at time~$t$  inherits an instantaneous debt balance $B_t = \Pi_t^\ast$ and the tail of the original term debt structure, and faces the   same  purchasing-power-adjusted state variable $X_t^\ast = B_t\, U_{C,t}^\ast$ as the initial Ramsey planner. We show that the continuation government then voluntarily confirms  the Ramsey plan.

Section~7 presents quantitative illustrations in a two-state regime-switching environment. When the initial term debt process $\{b_{0-,t}\}$ switches between  high and low levels while government spending is  constant, the Ramsey planner sets taxes low and interest rates low in the high-debt state, manipulating bond prices  to reduce the market value of its debt. The cumulative liability $\Pi_t^\ast$ jumps at state transitions, but $X_t = \Pi_t\, U_{C,t}$ remains continuous.  When instead  government spending $\{G_t\}$  switches while the  initial  debt service is constant, the Ramsey planner smooths  taxes  completely, holding consumption, the tax rate, and the interest rate all  constant --- an echo of \citet{barro1979determination}'s finding that intertemporal properties of an optimal tax rate can be disconnected from those of the government expenditure process.  

Section~\ref{sec:indeterminacy} constructs an alternative recursive formulation based on a  forward-looking   promised liability  $\Omega_t$ that represents   the present discounted value  of all future government liabilities. The associated state variable $Y_t = \Omega_t\, U_{C,t}$ has a law of motion whose coefficient on the  martingale increment is indeterminate, reflecting the government's freedom in choosing Arrow securities. Despite this indeterminacy, the value function retains the  same additively separable structure and the Ramsey plan is  unique: the optimal~$\Psi_0^\ast$ equals~$\Phi_0^\ast$, so consumption and tax  paths are identical to those derived in section~\ref{sec:bellmanramsey}.
Section \ref{sec:conclusion} offers concluding remarks.

\section{Environment \label{sec:environ}}
%of mass one 


Time $t \in [0, \infty)$ is continuous.  
A representative household and  benevolent government  participate  in  a complete set of  perfectly competitive markets.  The government finances a stochastic exogenous stream of expenditures with  a stream of distorting  flat tax rates. %There is no uncertainty. 
The exogenous  shock in the economy is driven by a time-homogeneous (continuous-time) Markov chain $\{s_t; t\geq 0\}$ where $s_t$ takes values in the set ${\bf N} = \{1, 2, \cdots, n\}$. The probability of $s_t$ transitioning from  $j\in {\bf N} $ to $k \in {\bf N}$ over a small time interval $\Delta t$ is  $\lambda_{j, k}\Delta t$ where $\lambda_{j, k} \geq 0$  for   $k\neq j$. The probability of remaining in  $j$ over $\Delta t$ is then $\big(1-\lambda_{j, j} \Delta t\big) $,  where $\lambda_{j, j} =  -\sum_{k\neq j} \lambda_{j, k}\leq 0$. %\nw{find the natural lan
%There are many competitive equilibrium with distorting taxes.
 At time $0$, a Ramsey planner    selects the best competitive equilibrium with distorting taxes.


%follows a $n-$state time-homogeneous Markov chain and 



% Variables describing government actions, quantities, and prices are continuous functions of time.  A continuous flow of government expenditures $\overrightarrow {G_0} = \{G_t;t\geq 0\}$ and a  continuous flow of  initial debt-payouts $\overrightarrow {b_0} = \{b_{0,\,t}; t\geq 0 \}$ are exogenous.

%\footnote{Let $\beta_{0,t}$ be the accumulative amount of term debt matured up to time $t$. We focus on the smooth accumulative term debt policy: $d \beta_{0,t}=b_{0,t} dt$ as in DNY. Note that $b_{0,t}$ can jump but $\beta_{0,t}$ cannot. In the basedline model, we assume $b_{0,t}$ is continuous in time. In an extended model, we relax this assumption to allow jumps in $b_{0,t}$.} 


The representative  household supplies  a labor process $%\overrightarrow {N_0}= 
\{N_{0,\,t}; t\geq 0 \}$ that produces a flow  of a single nonstorable good that can be divided between a consumption flow $%\overrightarrow {C_0} =
 \{C_{0,\,t}; t\geq 0 \}$  and a government
expenditure flow  $%\overrightarrow {C_0} =
 \{G_{0,\,t}; t\geq 0 \}$: %$\overrightarrow {G_0}$: %Feasible allocations require  
\begin{equation}\label{resourceconstraint}
	%C_{0, t}+G_t = N_{0, t} \, , 
	C_{0, t}+G_t = N_{0, t} \, ,\; \text{for\; all}\; t\geq 0\, . 
\end{equation}


%Technically, we  express this requirement as $ {\{C_{0, t}; t \geq 0\}}  \in {\mathbb S}^2$, where ${\mathcal S}^2$ denotes a (Sobelev) space of twice continuously differentiable functions.  
%The economy is populated by a continuum of infinitely-lived identical households. 
At time $0$, the representative household orders   consumption  and labor supply streams $\{C_{0, t}; t\geq 0\}_{}$ and
% labor supply streams 
$\{N_{0, t}; t\geq 0\}_{}$ according to
\begin{equation}\label{utility}
	\mathbb{E}_0 \left[\int_{0}^{\infty}  e^{-\rho t} {U(C_{0, t}, N_{0, t}) }dt\right]\, , 
\end{equation}
% \begin{equation}\label{utility}
%	\int_{0}^{\infty} e^{-\rho t} U(C_{0, t}, N_{0, t}) dt, 
%	%\int_{0}^{\infty} e^{-\rho t} U(C_{0,t}, N_{0,t}) dt, 
%\end{equation}
where $\rho>0$ and $U(\cdot, \cdot)$ is strictly increasing in consumption $C$, strictly decreasing in labor supply $N$, globally concave, and continuously differentiable.  Along with \citet*{debortoli2021optimal}, we assume that labor supply $N_{0, t}$ has  no upper bound. 




Let $\tau_{0, t}$  denote a tax rate at $t$ chosen at time 0 and   $q_{0, t}$ be the time-0 value of a zero-coupon bond with a unit payoff at $t$. %Going forward, as for $\tau_{0, t}$ and   $q_{0, t}$, 
A first  subscript  denotes the  time that a  variable is chosen and that a  second subscript denotes a time that the variable is realized. 



The representative household faces a single intertemporal  budget constraint: 
\begin{equation}\label{budgetc}
	\mathbb{E}_0 \left[ \int_{0}^{\infty}  C_{0, t} q_{0, t}  \, dt \right]\leq \mathbb{E}_0 \left[ \int_{0}^{\infty} {b_{0-, t} q_{0, t}} dt \right] + \mathbb{E}_0\left[ \int_{0}^{\infty}  (1-\tau_{0, t})N_{0, t} q_{0, t} dt \right] \, . 
\end{equation}
%\begin{equation}\label{budgetc}
%	\int_{0}^{\infty}  C_{0, t} 	q_{0,t} \, dt \leq   \int_{0}^\infty b_{0, t} 	q_{0,t}dt + \int_{0}^{\infty} (1-\tau_{0, t})N_{0, t}  	q_{0,t}dt \, . 
%\end{equation}
%\begin{equation}\label{budgetc}
% \int_{0}^{\infty}  C_{0, t} \frac{M_{0, t}}{M_{0, 0}}  \, dt \leq B_{0-} +\int_{0}^\infty b_{0-, t}\frac{M_{0, t}}{M_{0, 0}}dt + \int_{0}^{\infty} \left[ (1-\tau_{0, t})N_{0, t} \frac{M_{0, t}}{M_{0, 0}}  \right]dt \, . 
%\end{equation}
%\begin{equation}\label{budgetc}
% \int_{0}^{\infty} C_{0, t} \frac{M_{0, t}}{M_{0, u}}  \, dt \leq \frac{M_{0, 0}}{M_{0, u}} B_{0-} +\int_{0}^\infty b_{0-, t}\frac{M_{0, t}}{M_{0, u}}dt + \int_{0}^{\infty} \left[ (1-\tau_{0, t})N_{0, t} \frac{M_{0, t}}{M_{0, u}}  \right]dt \, . 
%\end{equation}
%\begin{eqnarray}\label{budgetc}
% \int_{s}^{\infty} C_{0, t} \frac{M_{0, t}}{M_{0, u}}  \, dt &\leq & \frac{M_{0, 0}}{M_{0, u}}  B_{0-} +\int_{s}^\infty b_{0-, t}\frac{M_{0, t}}{M_{0, u}}dt + \int_{s}^{\infty} \left[ (1-\tau_{0, t})N_{0, t} \frac{M_{0, t}}{M_{0, u}}  \right]dt \notag \\
% &+&\, \frac{M_{0, 0}}{M_{0, u}}  \int_0^s \left[ b_{0-, t} + (1-\tau_{0, t})N_{0, t}  - C_{0, t}\right]  d t . 
%\end{eqnarray}
The left side of (\ref{budgetc}) is the present value of  the household's consumption stream $\{ C_{0,t};  t\geq 0\}$ and the right side is the  sum of the household's financial wealth and its human wealth in the form of the present value of  after-tax labor income.
Given  tax rate process   $\{\tau_{0, t}; t\geq 0\}_{}$ and  bond price process $\{ q_{0, t}; t\geq 0\}_{}$, the household chooses $\{C_{0, t}, N_{0, t}; t\geq 0\}_{}$ to maximize \eqref{utility} subject to constraint (\ref{budgetc}).
That optimum problem yields
the following 
first-order necessary  conditions: 
\begin{eqnarray}
	1-\tau_{0, t}& = & -\frac{U_{N, t}}{U_{C, t}}  \, ,  \label{foch1} \\ && \notag \\
	q_{0,t} & = & %\frac{ e^{-\rho t} U_C(C_t, L_t) }{ \Phi} =
	e^{-\rho t}\frac{U_{C, t}}{U_{C, 0}}  \label{foch2}\, . 
\end{eqnarray}
The tax rate $\tau_{0, t}$ and  the stochastic discount factor (SDF)  $q_{0,t}$ are both stochastic processes.  
%Define a time-discount-factor process $\{M_{0,t}\}_{t \geq 0}$ by
%\begin{equation}\label{sdf0}
%\frac{M_{0, t}}{M_{0, 0}} \equiv  q_{0, t} = e^{-\rho t} \frac{U_{C, t}}{U_{C, 0}}\, .
%\end{equation}

%Equation (\ref{foch1}) is the standard FOCs for the tradeoff between consumption $C$ and labor $N$ for a given  tax rate $\tau_{0, t}$.  Equation (\ref{foch2}) equates the time-0 price of the equilibrium zero-coupon bond maturing at $t$, $q_{0,t}$, with the household's marginal rate of substitution (MRS) between time $t$ and  $0$ on the right side of (\ref{foch2}).
%

%\WJ{At every $t$, the household buys / sells government's short-term debt $B_t$ to satisfy a dynamic budget constraint. That is for a short period $[t, t+dt]$, $B_t$, the household's wealth at time $t$, will increase by discounting factor $\rho B_t$, time $t$ due debt $b_{0-, t}$, household's after tax income $(1-\tau_{0, t})N_{0, t}$, and decrease by consumption $C_{0, t}$. The increased quantity of household's wealth at time $t$ will be passed to time $t+dt$ at government's announced DF. Therefore, we have
%	\begin{equation*}
%	U_C(C_{0, t}, N_{0, t})  \left[ B_t +\left(\rho  B_t-C_{0, t}+(1-\tau_{0, t})N_{0, t}+b_{0-, t} \right) dt \right]  = B_{t+dt} U_C(C_{0, t+dt}, N_{0, t+dt})
%	\end{equation*} 
%	Taking the limit, we obtain: %\begin{equation}\label{dB}
%	%d (B_t U_C)  =  \left[\rho (B_t U_C)-C_{0, t}U_{C}-N_{0, t}U_{N} +b_{0-, t}U_{C}\right]dt. 
%	%\end{equation}
%	\begin{equation}\label{dB}
%	d (B_t M_{0,t})  =  \left[\rho (B_t M_{0,t})-C_{0, t}M_{0,t}-N_{0, t}U_{N}(C_{0, t}, N_{0, t}) +b_{0-, t}M_{0,t}\right]dt. 
%	\end{equation}
%	where $M_{0,t}$ is the marginal utility of consumption determined at time $0$: 
%	\begin{equation}
%	M_{0,t} = U_C(C_{0, t}, N_{0, t})\, . 
%	\end{equation}	
%}

%\WJ{We propose to use short-term debt rollover for the government  to satisfy their budget constraints. %period-by-period cash flow balance. and household
%	That is government purchases / sells only the short-term debt to fill the payment gap of each period without touching the long-term debt $\{b_{0-, t}\}$. In Lucas-Stocky, the government needs to buy long-term debt with different maturity to maintain its period-by-period balance. However, it's not clear what quantity for different maturity debt the government needs to buy at every $t$. }


A time$-0$ government finances  an  initial debt and its exogenous spending stream by taxing labor income:
% This implies the following time$-0$ government budget constraint: 
\begin{equation}\label{gbudgetc} 
\mathbb{E}_0 \left[ \int_{0}^{\infty} b_{0-, t} q_{0, t} dt \right] \leq  {\mathbb E}_0 \left[\int_0^{\infty}  q_{0, t} \left( \tau_{0, t} N_{0, t} -G_{t} \right) d t\right] \, ,
\end{equation} 
%\begin{equation}\label{gbudgetc}
%	\int_{0}^{\infty}  (\tau_{0, t}N_{0, t}-G_{t})	q_{0,t} \, dt \geq   \int_{0}^\infty b_{0, t} 	q_{0,t}dt  \, . 
%\end{equation}
The left  side is   time 0  value of payouts on government debt: $\{b_{0, t}; t\geq 0\}$ and the right  side is time 0 value of $\{S_{0, t}; t\geq 0\}$. We let $S_{0, t}=\tau_{0, t}N_{0, t}-G_t$ denote the primary surplus. 

\begin{definition}\label{defn1}
	Given the government's initial debt term structure $\{b_{0-,t};\, t \ge 0\}$, spending flow process $\{G_t;\, t \ge 0\}$, and a flat rate tax process $\{\tau_{0,t};\, t \ge 0\}$, a \textbf{competitive equilibrium} is a feasible allocation, $\{C_{0,t}, N_{0,t};\, t \ge 0\}$, and a bond price process $\{q_{0,t};\, t \ge 0\}$ for which
	\begin{itemize}
		\item the government's budget constraint~\eqref{gbudgetc} is satisfied and
		\item the allocation $\{C_{0,t}, N_{0,t};\, t \ge 0\}$ solves the household's optimization problem.
	\end{itemize}
\end{definition}
%\begin{definition} \label{defn1} 
%	Given the government's   initial  debt term structure  $\{b_{0-, t}; t\geq 0\}$ and  spending flow process $\{G_t;t\geq 0\}$,
%	a competitive equilibrium is a feasible  allocation $\{C_{0, t}, N_{0, t}; t\geq 0\}$, % satisfying (\ref{resourceconstraint}),
%	a flat rate tax process $\{\tau_{0, t}; t\geq 0 \}$,  and a bond price process $\{q_{0,t}; t\geq 0\}$ for which	\begin{itemize}
%	\item  the government's budget constraint \eqref{gbudgetc} is satisfied, and 
%	\item  given the government spending,  tax, and bond price process, the allocation solves the household's optimization problem
%		\end{itemize}
%\end{definition}


We follow \cite{LucasStokey1983} and use first-order conditions  \eqref{foch1} and  \eqref{foch2} together with  feasibility constraint \eqref{resourceconstraint} to eliminate tax rates and bond prices from the
government budget constraint. We thereby obtain the following implementability constraint on competitive equilibrium allocations $\{C_{0,t},  N_{0,t}\}_{t=0}^\infty$: % with distorting labor taxes : 
%\begin{equation}
%	%\int_{0}^{\infty} \left(C_{0, t} U_C + (C_{0, t}+G_t)U_N \right) dt = \int_{0}^{\infty} b_{0, t}U_C dt \,.
%	\int_{0}^{\infty} e^{-\rho t} ( C_{0, t}U_{C, t}+(C_{0, t}+G_t)U_{N, t} )  dt  \geq   \int_{0}^{\infty} e^{-\rho t}b_{0, t} U_{C, t} dt \, .  \label{impletconstraint}
%\end{equation}
\begin{equation} %\label{imc2}
	\mathbb{E}_0  \left[\int_{0}^{\infty} e^{-\rho t} ( C_{0, t}U_{C, t}+(C_{0, t}+G_t)U_{N, t} )  dt\right]  \geq \mathbb{E}_0 \left[ \int_{0}^{\infty}  e^{-\rho t} b_{0-, t} U_{C, t} dt \right]  \, .  \label{impletconstraint}
\end{equation}
%This is the implementability constraint that enters into the government's optimization problem.







\begin{proposition}\label{propo::1}
	Given an initial  term-debt  structure  $\{b_{0-, t}; t\geq 0\}$ and government spending path $\{G_t;t\geq 0\}$, an allocation $\{C_{0, t}, N_{0, t}; t\geq 0\}$ is a competitive equilibrium allocation if and only if it satisfies (\ref{resourceconstraint}) for all $t\geq 0$ and (\ref{impletconstraint}) at $t=0$. 
\end{proposition} % \TS{Tom: Maybe we don't need feasibility -- it might be build into implementability already.}
%
%A government shall commit to its outstanding short-term debt such that the immediately matured short-term debt's purchasing power will not be manipulated. To gurantee such commitmence, we focus on the following smooth equilibrium.
%
%
%
%\begin{definition}
%	A smooth equilibrium is a competitive equilibrium with allocation $\{C_{0, t}, B_{0, t}; t\geq 0\}$ such that
%	\begin{equation} \label{smooth}
%	\lim_{s\downarrow t} C_{0, s} = C_{0, t}\,.
%	\end{equation}
%\end{definition}
%
%\WJ{The smooth equilibrium in our continuous-time model is different from that converged from the discrete-time model. For the later, there is a discrete jump from time-0 consumption to time-1 consumption. When the length of time interval shrinks to zero, i.e., $\Delta t\to 0$, the commitment policy converges to a consumption series where consumption at time zero is infinity, such that the inherited short-term debt, i.e., $B_0$, can be inflated away. In our model, we assume the government must commit to all debts such that the value of short-term debt can not be inflated away immediately, which can be ensured by the smooth equilibrium condition (\ref{smooth}).
%}

%\begin{eqnarray} 
%&\lim_{s\downarrow t} B_sU_{C, s} = B_t U_{C, t};\, \lim_{s\downarrow t} B_s = B_t
%\end{eqnarray}



\paragraph{Ramsey problem.}
A Ramsey planner chooses a process   $ {\{C_{0, t}; t \geq 0\}}$  that maximizes 
%\eqref{utility}.
\begin{align} \label{obj}
	\mathbb{E}_0 \bigg[ \int_{0}^{\infty}  e^{-\rho t} U(C_{0, t}, C_{0, t}+G_t) dt  \bigg]%\\
	%s.t. &  %\\
\end{align} %+B_{0}U_{C}
%\begin{align} \label{obj}
%	&  \int_{0}^{\infty}  e^{-\rho t} U(C_{0, t}, C_{0, t}+G_t) dt  %\\
%	%s.t. &  %\\
%\end{align} %+B_{0}U_{C}
subject to 
%the smoothness condition (\ref{smooth}) and the 
implementability constraint \eqref{impletconstraint}.  By using first-order conditions
\eqref{foch1} and \eqref{foch2}, we can 
compute  an associated tax rate process and bond price system. 
After the Ramsey planner chooses   a tax  plan and associated price system,  
the representative agent's best  responses confirm the associated  competitive equilibrium allocation.

{To ease exposition, for the remainder of the paper we follow  \cite{LucasStokey1983} by focusing on a Markovian setting:} 
\begin{assumption}\label{assump1}
%Let $\{G_t; t\geq 0\}$  denote an exogenously specified %government spending 
%government spending process: 
% which  depends on the state variable $s_t$ via a time-invariant function: %$G(\,\cdot\,)$.  
%\begin{equation} \label{G_t}
%G_t= G(s_t),
%\end{equation}
%where %depends on the evolution of a state variable $s_t$.  We assume that 
%$s_t$ follows a $n-$state time-homogeneous Markov chain and takes values in the set $\{1, 2, \cdots, n\}$. The probability of $s_t$ transitioning from state $j$ to $k\neq j$ within a small time interval $\Delta t$ is  $\lambda_{j, k}\Delta t$ where $\lambda_{j, k} \geq 0$  for   $k\neq j$ and the probability of staying in state $j$ over $\Delta t$ is then $\big(1-\lambda_{j, j} \Delta t\big) $,  where $\lambda_{j, j} =  -\sum_{k\neq j} \lambda_{j, k}\leq 0$. %The transition dynamics induced by the Markov chain is fully captured by the generator matrix   which we denote by ${{\Lambda}}=[ \lambda_{j, k}]$ for $j, k \in\{1, 2, \cdots, n\}$.  %Simply put, t
	Exogenous flows of government expenditures  $ \{G_t;t\geq 0\}$ and 
	debt-service payouts $\{b_{0,\,t}; t\geq 0 \}$ are described by the following (time-homogenous) functions of the state variable $s_t$:  
\begin{equation} \label{G_t}
G_t= G(s_t) \quad \text{and} \quad b_{0-, t} = b(s_t)\, . 
\end{equation}
%Markov chain processes
	% {C\`adl\`ag processes, i.e., $\lim_{s\downarrow t} b_{0, s} = b_{0, t}$ and $\lim_{s\downarrow t} G_{s} = G_{t}$ for all $t\geq 0$;  and $\lim_{s\uparrow t} b_{0, s}$ and $\lim_{s\uparrow t} G_{s}$ exist for all $t>0$.}
\end{assumption}   

%\NW{We may not need the following part of this section?} 
%
%We next turn to how   \citeauthor{LucasStokey1983}'s Ramsey planner    reschedules an  initial {debt  structure} $\{ b_{0,t} \}_{t=0}^\infty$ to another  debt  structure $\{\hat b_{0,t}\}_{t=0}^\infty \neq  \{ b_{0,t}\}_{t=0}^\infty$, hoping to induce  future planners to     confirm the Ramsey  plan it has designed. %in order to motivate  continuation planners 
%
%
%%
%%\NW{Next, we introduce Lucas-Stokey's notion of implementation of the Ramsey plan. \citeauthor{LucasStokey1983}'s  Ramsey planner    chooses a particular  debt term structure $\{\hat b_{0,t}\}_{t=0}^\infty \neq  \{ b_{0,t}\}_{t=0}^\infty$.
%%%We use   the follow language.
%%}
%\subsection*{\normalsize \citeauthor{LucasStokey1983}'s continuation {debt  structure}}
%%\paragraph{Continuation debt term structure}
%To begin, we note that
%at  a Ramsey  allocation $\{C_{0,t}^\ast, N_{0,t}^\ast; t\geq 0\}$, flat tax rate process $\{\tau_{0,t}^\ast;  t\geq 0\}$, and price system $\{q_{0,t}^\ast;  t\geq 0\}$, many possible {debt   structures} $\{\hat b_{0,t} \}_{t=0}^\infty$  satisfy   time $0$ and time $t >0 $ continuation government budget constraints: 
%\begin{eqnarray}%\scriptsize{
%\mathbb{E}_0 \bigg[	\int_{0}^\infty \hat b_{0, s} 	q_{0,s}^*ds  &\leq &	\int_{0}^{\infty}  (\tau_{0, s}^\ast N_{0, s}^*-G_{s})	q_{0,s}^* \, ds \,  \bigg] \label{eqhatb1} \\ %}
%%		\end{equation*} 	\begin{equation*}%\scriptsize{
%\mathbb{E}_t \bigg[	\int_{t}^\infty \hat b_{0, s}  	q_{0,s}^*ds  &\leq & 	\int_{t}^{\infty}  (\tau_{0, s}^\ast N_{0, s}^*-G_{s})	q_{0,s}^* \, ds \, \bigg] \label{eqhatb2}.  %}
%	\end{eqnarray} 
%
%
%
%\begin{definition}
%A continuation of a Ramsey plan at $t >0$ is the tail $\{C_{0,s}^\ast, N_{0,s}^\ast, \tau_{0,s}^\ast, q_{0,s}^\ast; s\geq t\}$  
% of a Ramsey plan for $s \geq t$. A continuation  planner at  $t> 0$ must honor the (inherited) continuation {debt  structure}  $\{\hat b_{0,s}\}_{s=t}^\infty$  and chooses a  continuation plan    $\{\tau_{t, s}\}_{s=t}^\infty$ for tax rates    that need  not equal the continuation of the original Ramsey plan, $\{\tau_{0,s}^\ast\}_{s=t}^\infty$, and a new {debt  structure} $\{\hat b_{t, s}\}_{s=t}^\infty$.
% A Ramsey plan is said to be	 implemented (or ``time consistent'')	 when all  continuation Ramsey planners  (i.e., for all $t>0$) choose to confirm it.
%\end{definition}
%
% %  \WJ{Tom: time $0$ Ramsey  
%% \noindent Thus, a  continuation Ramsey planner at $s>0$ that must honor 
%%  debt term structure $\{\hat b_{0,t}\}_{t=s}^\infty$ constructs a %that was chosen by Ramsey a 
%% continuation  plan.
%
%	
%%	\TS{\tiny Do not mention  Ricardian equivalence and MM here as they refer to the first-best world allocation where taxes are not distortionary?} 
%
%%	\TS{Neng and Wei: in the above limits, I want $0$ in the lower limit in all 3 places}
%
%%\TS{I changed the lower bound to 1 from 0 for the integrals} 
%	
%
%\NW{Figure \ref{fig:timing}  illustrates the   timing protocol used by  \citet*{debortoli2021optimal} in their model with no shocks.} At time $0$, a Ramsey planner confronts a
%government purchase stream $\{G_t\}_{t=0}^\infty$ and an initial {debt  structure} $\{b_{0,t}\}_{t=0}^\infty$ that it must finance with a sequence  of flat rate taxes  $\{\tau_{0,t}\}_{t=0}^\infty$  and possibly restructured {debt} $\{\hat b_{0,t}\}_{t=0}^\infty$.  Knowing the tax and debt structure   sequences,  at  time $0$ the representative agent chooses
%$\{C_{0,t}, N_{0,t}\}_{t=0}^\infty$.  In  \citeauthor{debortoli2021optimal}'s laboratory, a single \NW{continuation Ramsey planner} at time $t=1$ confronts    $\{G_{t}\}_{t=1}^\infty$ and $\{\hat b_{0,t}\}_{t=1}^\infty$ that it must finance
%with a time $t=1$ continuation tax plan $\{\tau_{1,t}\}_{t=1}^\infty$.
%
%
%\begin{figure}[t]
%	\caption{Timing protocol for Ramsey planner and \NW{time $1$ continuation Ramsey planner.}}
%	\centering
%	\begin{tikzpicture}[scale=0.85]
%	
%	\draw [thick,->] (-1.2,0) -- (9.7,0);
%	%\draw [thick,-] (-0.19,0.1) -- (-0.19,0);
%	
%	%\node[align=center] at (4.1,1.3) {\bf  \ZL{ \scriptsize Step 2}};
%	\node[align=center] at (4.1, 1.3) {\bf  \NW{\scriptsize Household}};
%	\node[align=center] at (4.1, 0.9) {\bf  \NW{\scriptsize chooses $\{C_{0, t}, N_{0, t}\}_{t=0}^{\infty}$}};
%	
%	\draw [thick,->] (-1.2,0.1) -- (-1.2,0.5);
%	\node[align=center] at (-0.8, 1.9) {\bf \WJ{\scriptsize $\{b_{0, t}\}_{t=0}^{\infty},\; \{G_t\}_{t=0}^{\infty}$  }};
%	\node[align=center] at (-0.8, 1.3) {\bf \TS{\scriptsize Ramsey planner  }};
%	\node[align=center] at (-0.8, 0.9) {\bf \TS{\scriptsize chooses $\{\tau_{0, t}\}_{t=0}^{\infty},\; \{\hat{b}_{0, t}\}_{t=0}^{\infty}$  }};
%	%\node[align=center] at (-1, 1.3) {\bf  \NW{ \scriptsize  Step 3}};
%	
%	%\draw [thick,-] (5.20,0.23) -- (5.20,0);
%	\draw [thick,decorate,decoration={brace,amplitude=12pt,raise=4pt}] (-1.15,0) -- (9.5,0);
%	
%	
%	%\node[align=center] at (8.9, 1.3) {\bf   \YC{\scriptsize  Step 1}};
%	%	\node[align=center] at (8.9, 0.9) {\bf   \TS{\scriptsize Agent pays $M_{\tau}$  }};
%	%	\draw [thick,-] (9.4,0.23) -- (9.4,0);
%	%\draw [thick,decorate,decoration={brace,amplitude=12pt,raise=4pt}] (5.25,0) -- (9.35,0);
%	
%	
%	%\draw [thick,decorate,decoration={brace,amplitude=12pt,raise=4pt}] (9.45,0) -- (10.75,0);
%	
%	\draw[fill,color=blue] (-1.2,0.005) circle (.05);	
%	%\draw[fill,color=black] (-0.2,0.005) circle (.05);
%	
%	%\node[align=center] at (5.2,-0.3) {\footnotesize $ \tau^d$};
%	%\node[align=center] at (4.2,-0.3) {$\tau_F$};
%	\node[align=center] at (-1.2,-0.3) {\footnotesize $0 $};
%	\node[align=center] at (10.1,0) {\;\footnotesize  \; time};
%	\node[align=center] at (9.6,-0.3) {\footnotesize $\infty $};
%	
%	\end{tikzpicture}
%	
%	\bigskip
%	\begin{tikzpicture}[scale=0.85]
%	
%	
%	\draw [thick,->] (1.5,0) -- (9.7,0);
%	%\draw [thick,--] (-1.2,0) -- (9.7,0);
%	%\draw [thick,-] (-0.19,0.1) -- (-0.19,0);
%	
%	%\node[align=center] at (4.1,1.3) {\bf  \ZL{ \scriptsize Step 2}};
%	\node[align=center] at (5.6, 1.3) {\bf  \NW{\scriptsize Household}};
%	\node[align=center] at (5.6, 0.9) {\bf  \NW{\scriptsize chooses $\{C_{1, t}, N_{1, t}\}_{t=1}^{\infty}$}};
%	
%	\draw [thick,->] (1.5,0.1) -- (1.5,0.5);
%	\node[align=center] at (1.3, 1.9) {\bf \WJ{\scriptsize $\{\hat{b}_{0, t}\}_{t=1}^{\infty},\; \{G_t\}_{t=1}^{\infty}$  }};
%	\node[align=center] at (-0.8, 0.9) {\bf \color{white}{\scriptsize Ramsey planner commits }};
%	\node[align=center] at (1.1, 1.3) {\bf \TS{\scriptsize Continuation Ramsey planner  }};
%	\node[align=center] at (1.1, 0.9) {\bf \TS{\scriptsize chooses $\{\tau_{1, t}\}_{t=1}^{\infty}, \{\hat{b}_{1, t}\}_{t=1}^{\infty}$ }};
%	
%	%\node[align=center] at (-1, 1.3) {\bf  \NW{ \scriptsize  Step 3}};
%	
%	%\draw [thick,-] (5.20,0.23) -- (5.20,0);
%	\draw [thick,decorate,decoration={brace,amplitude=12pt,raise=4pt}] (1.6,0) -- (9.5,0);
%	
%	
%	%\node[align=center] at (8.9, 1.3) {\bf   \YC{\scriptsize  Step 1}};
%	%	\node[align=center] at (8.9, 0.9) {\bf   \TS{\scriptsize Agent pays $M_{\tau}$  }};
%	%	\draw [thick,-] (9.4,0.23) -- (9.4,0);
%	%\draw [thick,decorate,decoration={brace,amplitude=12pt,raise=4pt}] (5.25,0) -- (9.35,0);
%	
%	
%	%\draw [thick,decorate,decoration={brace,amplitude=12pt,raise=4pt}] (9.45,0) -- (10.75,0);
%	
%	%\draw[fill,color=white] (-1.2,0.005) circle (.05);	
%	\draw[fill,color=blue] (1.5,0.005) circle (.05);
%	%\draw[fill,color=black] (-0.2,0.005) circle (.05);
%	
%	%\node[align=center] at (5.2,-0.3) {\footnotesize $ \tau^d$};
%	%\node[align=center] at (4.2,-0.3) {$\tau_F$};
%	%\node[align=center] at (-1.2,-0.3) {\footnotesize $0 $};
%	\node[align=center] at (1.5,-0.3) {\footnotesize $1 $};
%	\node[align=center] at (10.1,0) {\;\footnotesize \; time};
%	\node[align=center] at (9.6,-0.3) {\footnotesize $\infty $};
%	
%	\end{tikzpicture}
%	
%\label{fig:timing}	
%\end{figure}
%
%
%
% \citet{LucasStokey1983}  constructed examples in  which     budget-feasible {debt  structures} $\{\hat b_{0,t} \}_{t=1}^\infty$ induce a continuation Ramsey planner to confirm a continuation of  a Ramsey tax plan $\{\tau_{0,t}^\ast;  t\geq 1\}$ and its associated allocation and price system: $\{C_{0,t}^\ast, N_{0,t}^\ast, q_{0,t}^\ast; t\geq 1\}$. 
%However,  \citet*{debortoli2021optimal} constructed examples  for which  debt  term structure  $\{\hat b_{0,t} \}_{t=1}^\infty$  along  lines recommended by     \citet{LucasStokey1983} fails to  induce a continuation Ramsey planner to confirm a continuation {of} the Ramsey  plan.
%To set the stage for a continuous time version of   \citeauthor{debortoli2021optimal}'s counterexample, the next section presents  the  Ramsey planner's  Lagrangian.
%






\section{A  Lagrangian}\label{sec:Lagrangian1}

In this section, we construct a  Lagrangian for the Ramsey problem, use it to characterize the equilibrium allocation associated with the Ramsey plan, then derive equilibrium asset-pricing implications. 
%In this section, we use  
\subsection{Posing the Ramsey Problem} 
%a Lagrangian approach to solve the Ramsey problem. 
Attach a  multiplier  $\Lambda_0$  to implementability constraint  (\ref{impletconstraint}) and form  a Lagrangian: 
%\begin{align}\label{L0}
%	{\mathcal L}^{}_0 & =  \int_{0}^{\infty}  e^{-\rho t}\left[ U(C_{0, t}, C_{0, t}+G_t)+ \Lambda_0^{}\left( C_{0, t}U_{C, t}+(C_{0, t}+G_t)U_{N, t}-b_{0, t}^{} U_{C, t}\right)  \right] dt \, . %\nonumber \\%- \Psi_0(b_{0, t}^{} U_C- b_{0-, t} U_C)	\right]dt \\
%	%&  \quad \quad \;  \, .
%\end{align} %  to  minimize   \eqref{L0}
\begin{align}\label{L0}
	{\mathcal L}^{}_0 & = \mathbb{E}_0 \bigg[ \int_{0}^{\infty}  e^{-\rho t}\left[ U(C_{0, t}, C_{0, t}+G(s_t))+ \Lambda_0^{}\left( C_{0, t}U_{C, t}+(C_{0, t}+G(s_t))U_{N, t}-b(s_t)^{} U_{C, t}\right)  \right] dt \bigg] \, . %\nonumber \\%- \Psi_0(b_{0, t}^{} U_C- b_{0-, t} U_C)	\right]dt \\
	%&  \quad \quad \;  \, .
\end{align} %  to  minimize   \eqref{L0}
The Ramsey planner  maximizes the right side of  \eqref{L0} over consumption plans $\{C_{0,t}\}_{t=0}^\infty$   and minimizes over the nonnegative multiplier $\Lambda_0$.  We denote the extremizing values $\{C_{0,t}^*\}_{t=0}^\infty$ and $\Lambda_0^*$.
% Maximize under the integral sign on the right side of \eqref{L0} to get a continuum of static problems: 
%% We characterize the optimal policy for time $0+$ onward as follows.  For a given time interval $(t, t+dt), t>0$, the planner chooses consumption to solve
%\begin{equation} \label{optp-L}
%	\max_{C_{0, t}} \;\; \left[ U(C_{0, t}, C_{0, t}+G_t)+ \Lambda_0^{}\left( C_{0, t}U_C+(C_{0, t}+G_t)U_N-b_{0, t} U_C \right) \right]dt\,, \;\; t \geq 0
%\end{equation} %If there are no constraints on $C_{0, t}$, then 
For a given $\Lambda_0$, the optimal consumption plan satisfies
\footnote{{We assume  a regularity condition under which 
	the optimal consumption plan for the problem defined in   \eqref{Csoln}   is unique for a  given $\Lambda_0$.}}
\begin{equation}\label{Csoln}
	C(s_t; \Lambda_0):=\arg\max_{C_{0, t}}  \left[ U(C_{0, t}, C_{0, t}+G(s_t))+ \Lambda_0\left( C_{0, t}U_{C, t}+(C_{0, t}+G(s_t))U_{N, t}-b(s_t) U_{C, t} \right) \right] .
\end{equation}
%\footnote{We assume  a regularity condition under which 
%the optimal consumption plan for the problem defined in   \eqref{Csoln}   is unique for a  given $\Lambda_0$.} 
%\begin{equation}\label{Csoln}
%	C(\Lambda_0; b_{0, t}, G_t):=\argmax_{C_{0, t}}  \left[ U(C_{0, t}, C_{0, t}+G_t)+ \Lambda_0\left( C_{0, t}U_C+(C_{0, t}+G_t)U_N-b_{0, t} U_C \right) \right] .
%\end{equation}
%solve problem (\ref{optp-L}) 
%, and $\tau(\Lambda_0^{}, b_{0, t}, G_t)$, respectively.
Substituting  \eqref{Csoln} into  implementability constraint (\ref{impletconstraint}),  we deduce that   % Lagrange multiplier 
$\Lambda_0$ must satisfy: 
\begin{eqnarray}\label{eqPhi}
	&& \mathbb{E}_0 \, \bigg[ \int_{0}^{\infty} e^{-\rho t} \big[ \left( C(s_t; \Lambda_0)-b(s_t)\right)U_{C, t}+(C(s_t; \Lambda_0)+G(s_t))U_{N, t} \big]dt  \bigg] = 0\,. %\mathbb{E}_0\, \int_{0}^{\infty} e^{-\rho t}b(s_t) U_{C, t} dt\,.
	%&&  +B_{0}U_{C}( C(\Lambda_0^{}, b_{0, 0}, G_{0}) ,C(\Lambda_0^{}, b_{0, 0}, G_{0})+G_{0}) .
\end{eqnarray}
%\begin{eqnarray}\label{eqPhi}
%	&& \int_{0}^{\infty} e^{-\rho t} \big[ C(\Lambda_0; b_{0, t}, G_t)U_C+(C(\Lambda_0; b_{0, t}, G_t)+G_t)U_N \big]dt  =  \int_{0}^{\infty} e^{-\rho t}b_{0, t} U_C dt\,.
%	%&&  +B_{0}U_{C}( C(\Lambda_0^{}, b_{0, 0}, G_{0}) ,C(\Lambda_0^{}, b_{0, 0}, G_{0})+G_{0}) .
%\end{eqnarray}
%\begin{eqnarray}\label{eqPhi}
%	&& \int_{0}^{\infty} e^{-\rho t} ( C(\Lambda_0^{*}, b_{0, t}, G_t)U_C+(C(\Lambda_0^{*}, b_{0, t}, G_t)+G_t)U_N )  dt  =  \int_{0}^{\infty} e^{-\rho t}b_{0, t} U_C dt\,.
%	%&&  +B_{0}U_{C}( C(\Lambda_0^{}, b_{0, 0}, G_{0}) ,C(\Lambda_0^{}, b_{0, 0}, G_{0})+G_{0}) .
%\end{eqnarray}
%where we impose the smooth consumption condition
%\begin{equation}
%	b_{0-, t} \text{is cadlag process.}
%\end{equation}
%\NW{Wei: maybe this is left continuous with right limit? cadlag process? } 
%\WJ{Need to justify this condition.  }
Substituting  $C_{0,t} = C(s_t; \Lambda_0)$ into the right side of   (\ref{L0}) lets us  write  $\mathcal L_0$ as a function %of $\Lambda_0^*$:
 $\mathcal{L}_0(\Lambda_0)$. 
An optimal  $\Lambda_0^{*}$ is  the non-negative root of equation (\ref{eqPhi}) that maximizes $\mathcal{L}_0(\Lambda_0)$. 
The Ramsey allocation satisfies $C_{0,t}^*=C(s_t; \Lambda_0^\ast)$ and $N_{0,t}^* = C_{0,t}^*+G_t = N(s_t; \Lambda_0^\ast)$. 
%\NW{	$C(\Lambda_0, b_{0, t}, G_t)  =  \argmax_{C_{0, t}}  \left[ U(C_{0, t}, C_{0, t}+G_t)+ \Lambda_0^{}\left( C_{0, t}U_C+(C_{0, t}+G_t)U_N-b_{0, t} U_C \right) \right]$.} 

%In summary, we conclude with

%\NW{make this prettier and maybe add an eqn as you noted.} 


%_{\{\Lambda_0 =  \textrm{a positive root of}  (\ref{eqPhi})\}}
\begin{proposition} \label{proplag} 
	%When %Assuming  , 
	Under  conditions provided in  Appendix \ref{appendix:A}, there  exists  a non-negative root  $\Lambda_0^*$ of equation (\ref{eqPhi}) that solves
	\begin{equation*}
		\Lambda_0^*  = \argmin_{\Lambda_0}  \,\,  \mathcal{L}_0(\Lambda_0)\, . %(\{C_{0,t}^*\}_{t=0}^\infty,  
	\end{equation*}	
	 The    Ramsey allocation,  tax rate, and bond price  process satisfy
	\begin{eqnarray*}
		&	 {C}^\ast_{0, t} =     C(s_t; \Lambda_0^*); \, \quad 
		{N}^*_{0, t}  =    {C}_{0, t}^\ast +G(s_t) ;  \\
		&	  {\tau}^*_{0, t} =  1+\frac{U_{N, t}({C}^*_{0, t}, {N}^*_{0, t})}{U_{C, t}({C}^*_{0, t}, {N}^*_{0, t})};\,   \quad
		{q}^*_{0,t}  = e^{-\rho t}\frac{U_{C, t}({C}^*_{0, t}, {N}^*_{0, t})}{U_{C, 0}({C}^*_{0, 0}, {N}^*_{0, 0})}\, .
	\end{eqnarray*} 
	%	        \begin{eqnarray*}
		%	        	 C(\Phi_0^{\mathcal L}, b_{0-, 0}, G_0)= \lim _{s\downarrow 0} C(\Phi_0^{\mathcal L}, b_{0-, s}, G_s)
		%	        \end{eqnarray*}	
	%	where 
	%and the marginal utilities $U_C$ and $U_N$ are evaluated at $\widetilde{C}_{0, t},\widetilde{N}_{0, t}$.
	The value ${\mathcal L}^{*}_0 $ of a Ramsey plan is
	%	\begin{equation*}
		$	{\mathbb E}_0 \big[ \int_{0}^{\infty}  e^{-\rho t} U(	{C}^*_{0, t}, {C}^*_{0, t}+G(s_t))  dt \big] \, .
		$	%\end{equation*}
	\end{proposition}
%	\NW{report $r_{0,t}^\ast$?} 
%	 	\begin{eqnarray*}
%		&	 {C}^\ast_{0, t} =     C(\Lambda_0^{*}; b_{0, t}, G_t); \, \quad 
%		{N}^*_{0, t}  =    {C}_{0, t}^\ast +G_t ;  \\
%		&	  {\tau}^*_{0, t} =  1+\frac{U_{N, t}({C}^*_{0, t}, {N}^*_{0, t})}{U_{C, t}({C}^*_{0, t}, {N}^*_{0, t})};\,   \quad
%		{q}^*_{0,t}  = e^{-\rho t}\frac{U_{C, t}({C}^*_{0, t}, {N}^*_{0, t})}{U_{C, 0}({C}^*_{0, 0}, {N}^*_{0, 0})}\, .
%	\end{eqnarray*}
%	%	        \begin{eqnarray*}
%	%	        	 C(\Phi_0^{\mathcal L}, b_{0-, 0}, G_0)= \lim _{s\downarrow 0} C(\Phi_0^{\mathcal L}, b_{0-, s}, G_s)
%	%	        \end{eqnarray*}	
%	%	where 
%	%and the marginal utilities $U_C$ and $U_N$ are evaluated at $\widetilde{C}_{0, t},\widetilde{N}_{0, t}$.
%	The value ${\mathcal L}^{*}_0 $ of a Ramsey plan is
%	%	\begin{equation*}
%	$	 \int_{0}^{\infty}  e^{-\rho t} U(	{C}^*_{0, t}, {C}^*_{0, t}+G_t)  dt.
%	$	%\end{equation*}
%\end{proposition}
%
%The primary surplus under the Ramsey  plan is 
%\begin{equation}\label{SfunRam}
%S_{0, t}^\ast = S(C_{0, t}^\ast ) =\tau_{0, t}^\ast N_{0, t}^\ast- G_t
%\,.
%\end{equation} %	(\ref{Sfun}). 
%The  time-$t$ value of cumulative deficits in the absent of any rescheduling  the initial  term debt  process $\{b_{0,t}\}$ is 
%\begin{eqnarray} \label{Pistar} 
%			\Pi_{t}^*  = \frac{1}{q_{0, t}^*}\int_{0}^t q_{0, u}^*\left(b_{0, u} - S(C_{0, u}^*)\right) du \, , % \\
%	%		& 	= & \int_{0}^t e^{\int_u^t r_{0,s}^\ast d s } \left(b_{0, u} - S(C_{0, u}^*)\right) du 
%		\end{eqnarray}
%		where the  increment to  the government's  IOU's over a small $du$ interval is 
% $\left(b_{0, u} - S(C_{0, u}^*)\right) du$ and ${q_{0, u}^*}/{q_{0, t}^*}$ converts a time $u$ deficit to its time $t$ value. 
%The stock  $	\Pi_{t}^* $  plays a crucial role in our way of  implementing the Ramsey plan. 
%% 
%		%	\NW{Add  expression with interest rate.} %	is time $t$ value of  %under Ramsey plan 
%%Recognizing that   $r_{0,t}^*$ is a Dirac delta at $t=1$ under the Ramsey allocation,  define 
%%		\begin{equation*}
%%		\Pi_{1-}^* =  \frac{1}{\TS{q_{0, 1-}^*}}\int_{0}^{1-} q_{0, s}^*\left(b_{0, s} - S(C_{0, s}^*)\right) ds. 
%%		\end{equation*}
%%\NW{
%%		\begin{equation*}
%%		\Pi_{1-}^* = C_0^* \frac{\left(e^{\rho}-1\right)}{\rho} \left( \frac{b_{0}}{C_0^*} -1 + (C_0^*+G)   \right)
%%		%\frac{1}{\TS{q_{0, 1-}^*}}\int_{0}^{1-} q_{0, s}^*\left(b_{0, s} - S(C_{0, s}^*)\right) ds. 
%%		\end{equation*}
%%}

%Absent term debt adjustment,   where $ S(C_{0, u}^*)$ is the primary surplus given in 
%the primary surplus $S_{0, t}$ can be written as % into a function of consumption, i.e., 
%}	\NW{

%\section{The Counterexample}\label{sec:counterexample}
%To construct a continuous-time version  of 
%\citeauthor{debortoli2021optimal}'s counterexample to  \citeauthor{LucasStokey1983}'s implementation of a Ramsey plan, we adopt 
%the special instantaneous utility function: 
%\begin{equation}\label{DNYUtility}
%U(C, N) = \log C - \eta \frac{N^\gamma}{\gamma}  \, , 
%\end{equation}
%where $\eta>0$ and $\gamma\geq 1$.
%The Ramsey planner must finance the time-invariant government expenditure process 
%%\begin{equation*}
%$G_t =G$ for all $t \geq 0$ 
%%\end{equation*}
%and an  initial {debt  structure}:
%\begin{equation}\label{DNYdebt}
%b_{0, t}  =
%\begin{cases} b_0 > 0  & t\in[0, T), \\ % for a fixed $T$; in DNY, $T=1$\medskip
%			0  &  t\geq T,
%			\end{cases}
%\end{equation} 
%for a given $T > 0$.  To ease exposition,  for the remainder of the paper, we'll work with the $\gamma=1$ case where $U_C=1/C$ and $U_N=-1$. The  tax rate  then satisfies {$\tau_{0,t}  = 1- \eta C_{0,t} $ for all $t\geq 0$} and the primary surplus $S_{0, t}$ given in (\ref{SfunRam}) takes the form of: \begin{equation}\label{Sfun}
%S_{0, t} = S(C_{0, t}) = C_{0, t}\left(1-\eta (C_{0, t} + G_t) \right)\,.
%\end{equation}
%
% 
%{Next, we state a  continuous-time version of Lemma 3 in 	\cite*{debortoli2021optimal} for the $\gamma=1$ case.} 
%\begin{lemma}
%	When $b_0$ is not too high,	a Ramsey plan exists and  is	\begin{eqnarray} 
%	& C_{0, t}^{*}  =  C_0^* 1_{\{t\in[0, T)\}} +  C_1^*1_{\{t\in[T, \infty)\}}; \quad
%	N_{0, t}^{*}  =  N_0^* 1_{\{t\in[0, T)\}} +  N_1^*1_{\{t\in[T, \infty)\}} ,  \label{RamseyCN} \\
%	& 	\tau_{0, t}^{*}   =   \tau_0^* 1_{\{t\in[0, T)\}} +  \tau_1^*1_{\{t\in[T, \infty)\}}; \quad 
%	q^*_{0,t}  = 
%	e^{-\rho t}\left( 1_{\{t\in[0, T)\}} +  C_0^*/C_1^*1_{\{t\in[T, \infty)\}} \right), \label{Ramseytauq} 
%	\end{eqnarray}
%	where $1_{\{A\}}$ is an indicator function that equals one if  the event $A$ occurs and zero otherwise, 
%	\begin{equation}
%	C_0^*=C(\Lambda_{0}^*; b_{0}, G) = \frac{2b_0\Lambda_0^*}{\sqrt{1+4\eta b_0\Lambda_0^*(1+\Lambda_0^*)}-1} ,  \quad C_1^*=C(\Lambda_{0}^*; 0, G) = \frac{1}{\eta(1+\Lambda_0^*)},  \label{C0C1DNY}
%	\end{equation}
%	$N_0^*=C_0^*+G$, and $N_1^*=C_1^*+G$. The Ramsey tax plan is given by $\tau_0^*=1-\eta C_0^*$ and $\tau_1^*=1-\eta C_1^*$. 
%	Finally, 	$\Lambda_{0}^*$ is the unique non-negative root of 
%	\begin{align}
%	\frac{1-e^{-\rho T}}{\rho}\left( 1-\eta(C(\Lambda_{0}; b_0, G)+G)-\frac{b_0}{C(\Lambda_{0}; b_0, G)} \right)+\frac{e^{-\rho T}}{\rho}\left(1-\eta(C(\Lambda_{0}; 0, G)+G)\right) = 0\, . \label{C0C1imc} 
%	\end{align}
%	%	It is also true that  $\Lambda_{0}^* ={\argmin}_{} \,\mathcal L_0(\Lambda_0)$. % $= \argmax_{\{\Lambda_{0}\geq 0: \Psi_0(\Phi_{0})=0\} } \;  \int_{0}^{\infty} e^{-\rho t}U(C(\Phi_0, b_{0, t}, G), C(\Phi_0, b_{0, t}, G)+G)dt .$
%	
%\end{lemma}
%
%
%As $C_{0,t}^\ast=C_0^\ast$ for all $t<T$ and $C_{0,t}^\ast=C_1^\ast$ for all $t\geq T$, the  bond price satisfies
%\begin{equation} \label{bondpricdny}
%q^*_{0,t}  = e^{-\rho t}, \; \text{for}\; t< T; \; \,\,\, q^*_{0,t}  = e^{-\rho t}\frac{C_0^*}{C_1^*}, \; \text{for}\; \; t\geq T. 
%\end{equation}  
%Consequently, the competitive equilibrium interest rate is  $r_{0, t}^* = \rho$ for both $t<T$ and $t> T$. 
%Immediately before  $t=T$,  the  interest  rate approaches $\infty$, so  {$\{r_{0, t}^*-\rho\}_{t=0}^\infty$ is a Dirac delta function}.
%We shall soon see how the  Dirac delta interest rate at $t=T$ plays a key role. 
%
%In the context of this example, we now describe how \citeauthor{debortoli2021optimal} use   \citeauthor{LucasStokey1983}'s way of structuring a continuation debt  term structure   in order to  motivate the continuation Ramsey planner to implement the continuation  of the   Ramsey plan. The Ramsey allocation  (\ref{RamseyCN}) is piece-wise linear, so the  Ramsey planner  repurchases $\{b_{0,t}\}_{t=0}^\infty$ at time $0$ and sells   debt with the following coupon payment schedule: $\hat b_{0,t} = \hat b_0 $ for all $t<T$ and $\hat{b}_{0,t}=\hat b_1$ for all $t\geq T$, making sure that $\hat b_0$ satisfies both  the budget constraint (\ref{eqhatb1}) at time $0$ and the budget constraint  (\ref{eqhatb2}) at time $t=T$. We  compute  $\hat b_0 $ and   $\hat b_1$ as follows. %to be determined. %This is  
%%$C_{0,t}^\ast=C_0^\ast$ for all $t<1$ and $C_{0,t}^\ast=C_1^\ast$ for all $t\geq 1$ 
%
%
%
%% Because 
%%\NW{In , to induce a {continuation Ramsey planner}  to continue   the Ramsey plan, the %set  a  continuation debt service flow given by: % equal to %\NW{$\hat{b}_1$} that satisfies			% to time−T gov’t%Within the set of feasible restructuring policies using term debt, 
%%%\TS{Term-debt restructuring policy}: \NW{gov't at time $T-$} leaves a unique term debt $\hat{b} _T$ to gov't at \NW{time $T$}\,%\hyperlink{explain1}{\beamerskipbutton{Details}}
%%}
%
%
%%\WJ{How to obtain $\hat{b}_1$?}
%
%\begin{itemize}
%\item 
%To compute $\hat b_1$, 
%%note that $q_{0, t}^* = e^{-\rho t}$ for $t<T$; $q_{0,t}^\ast =e^{-\rho t}C_0^*/C_1^*$  for $t\geq T$; 
%we use the bond price given in (\ref{bondpricdny}) and the property that the primary surplus under the Ramsey plan is constant for $t\geq 1$: $S(C_{1}^*) = \tau_1^\ast N_1^\ast- G = C_{1}^*\left( 1- \eta(C_1^*+G)\right)$ 
%%To induce time $1$ gov't to choose allocation $C_1^*$, leave a constant term-debt flow $\hat{b}$ to time $1$ government. In particular, 
%to simplify the time-$T$ budget equation (\ref{eqhatb2}) as follows: %dynamic budget constraint for time$-1$ gov't under the Ramsey plan is given by
%	\begin{equation*}
%		\int_{T}^{\infty} \hat{b}_1 e^{-\rho t} \frac{C_0^*}{C_1^*}dt = \int_{T}^{\infty}  S(C_{1}^*)e^{-\rho t} \frac{C_0^*}{C_1^*}dt \, .  %\Rightarrow \, \hat{b}_1 = S(C_{1}^*) \,.
%	\end{equation*} 
%This yields: $\hat{b}_1 = S(C_{1}^*)$. 
%Simplifying the time $0$ budget constraint (\ref{gbudgetc}) for the initial debt term structure $\{b_{0, t}\}$ we obtain
%\begin{equation*}
%\frac{1-e^{-\rho }}{\rho}\left( {S(C_{0}^*)}-{b_{0}} \right)+\frac{e^{-\rho }}{\rho}S(C_{1}^*)\frac{C_0^*}{C_1^*} = 0\, . 
%\end{equation*}	
%Combining the above two results gives $\hat b_{1} =  \left(e^{\rho T}-1\right) \left( \frac{{b_0}}{{C_0^*}} -1 + \eta({C_0^*}+G)   \right) {C_1^*}$.
%%	Knowing that 
%\item 
%Using the preceding expression for  $\hat b_1$ to rewrite the  time-0 budget constraint (\ref{eqhatb1}) as follows: %for time$-0$ under the Ramsey plan satisfies
%	\begin{equation*}
%		\int_{0}^{T} e^{-\rho t} \hat{b}_0 dt + \int_{T}^{\infty} e^{-\rho t} \hat{b}_1\frac{C_0^*}{C_1^*}dt = \int_{0}^{T} e^{-\rho t} S(C_0^*) dt + \int_{T}^{\infty} e^{-\rho t} S(C_1^*)\frac{C_0^*}{C_1^*}dt\, , 
%	\end{equation*} 
%we obtain  $\hat b_{0}  = {C_0^*}\left(1-\eta({C_0^*}+G)\right) $. 
%\end{itemize}
%In summary, we obtain the following Lucas-Stokey debt restructuring policy: 
%\begin{equation} \label{eqn:hatdebt}
%\hat b_{0, t} = \begin{cases}
%\hat b_{0}  = {C_0^*}\left(1-\eta({C_0^*}+G)\right),  & t \in [0, T) \, , \\
%  \hat b_{1} =  \left(e^{\rho T}-1\right) \left( \frac{{b_0}}{{C_0^*}} -1 + \eta({C_0^*}+G)   \right) {C_1^*}, 
%    & t \geq T \, . 
%\end{cases}
%\end{equation}
%
%
%\medskip
%\noindent We can state   key findings of \citet*{debortoli2021optimal} in the context of our continuous-time version of their model.
%\begin{lemma}\label{ref:lemmab^*}
% When $b_0$ is not too high,  {debt  structure} \eqref{eqn:hatdebt}  induces a time $T$ continuation Ramsey planner to confirm the Ramsey plan.  But for $b_0$ above a threshold $b^*$, {debt  structure} \eqref{eqn:hatdebt} is unable to induce the continuation planner   to confirm  the Ramsey plan.   
%\end{lemma}
%%\begin{lemma}
%%	When the initial debt is higher than a threshold $b^*$, i.e., $b_0>b^*$, the time$-T$ government that faces the unique $b_1$ optimally renege the Ramsey plan chosen by the time$-0$ government. 
%%\end{lemma}
%
%\noindent To verify Lemma \ref{ref:lemmab^*}, we would first note that a negative Lagrange multiplier $\Lambda_{T}^*$ on a   Lagrangian for a candidate continuation Ramsey plan indicates that the candidate is not optimal. Then we would execute the  following three steps.
%		\begin{itemize}
%		\item[1.] $C_1^*$ is a decreasing function of $b_0$, and $b_1/C_1^*$ is an increasing function of $b_0$.
%		\item[2.]  Because $C_1^*$ is too low,  when evaluated along  the continuation Ramsey plan, the multiplier $\Lambda_{T}^*$ on the implementability constraint  for the time $-T$ continuation Ramsey planner  is negative:
%		\begin{eqnarray*}
%			\Lambda_{T}^* & = & -\frac{1-\eta C_1^*\left(C_1^*+G \right)^{\gamma-1}}{b_1/C_1^* - \eta \gamma C_1^* \left( C_1^*+G \right)^{\gamma-1}}   \\
%			& = &  -\frac{1-\eta C_1^*\left(C_1^*+G \right)^{\gamma-1}}{1-\eta\left(C_1^*+G \right)^{\gamma} - \eta \gamma C_1^* \left( C_1^*+G \right)^{\gamma-1}}<0
%		\end{eqnarray*}
%		\item[3.]  When $b_0$ exceeds a threshold $b^*$, $\Lambda_{T}^*<0$ and the time $-T$ continuation planner chooses not to confirm the continuation Ramsey plan. 
%	\end{itemize}
	
%\noindent We shall soon apply  these three steps to  our continuous time version of \citeauthor{debortoli2021optimal}'s counterexample. 
%\end{proof}
%
%The intuition that the Ramsey plan is time inconsistent when $b_0$ is too high is as following. A future government that inherits $\widehat{b}_1$ from the past government has strong incentive to inflate the term debt because the inherited debt coupon flow is $dt$ level. By doing so, the future government may optimally choose a higher consumption level.
%
%{\cite*{debortoli2021optimal}'s implementation resembles %a debt restructuring that seems like a
%		 Alexander Hamilton's rescheduling of United States %Continental and state 
%		 government debts at the beginning of the first Washington administration in 1790.  Hamilton exchanged finite maturity Continental and state government bonds in exchange for   bundles of US federal  consols.\footnote{See \citet{hall2014fiscal}.} Notice how in  Figure \ref{fig:DNYexample} a Ramsey planner reschedules initial finite maturity debt streams with consols. %\TS{Neng and Wei: here is my second draft of a remark on Alexander Hamilton.} \NW{Tom: I like this reference to Hamilton's rescheduling of debt.}  and \ref{fig:DNYcounterexample}
%	}
%
%\paragraph{Numerical illustrations.} 
%
%To stay close to  \citeauthor{debortoli2021optimal}'s  discrete time example, we set    $T=1$ and $(\gamma,\eta) = (1,1) $, which makes  $ U(C, N) = \log C - N$.
% We use the \citeauthor{debortoli2021optimal}'s  one-period discount factor of $0.5$ by settomg $e^{-\rho}=0.5$, so that   $T=1$ corresponds to 13.9 years with an annual discount rate of $5\%$. 
% For these parameter values, it turns out that the threshold in Lemma \ref{ref:lemmab^*} is $b^* = 0.35$. By setting $b_0 = .3$, we recover an example in which committing the government to finance  Lucas and Stokey's recommended restructured debt term structure works.  By setting $b_0 = .5$, we'll construct a version of \citeauthor{debortoli2021optimal}'s counterexample in which it doesn't work. 
% 
% 
%\noindent\textbf{Lucas-Stokey Implementation works when $b_0=0.3$}.  When  $b_{0,t}= b_0 =0.3$ for $t\in [0, 1)$ and $b_{0,t}= 0$ for $t\geq 1$, %\TS{Have to say what $b_{0,t}$ also is for $t \geq 1$} 
%%\TS{Wei -- please fill in numbers here.}
%the Ramsey planner sets
%	 $C_{0,t}^*= C_0^\ast = 0.6987$  for  $t\in [0, 1), $  	 $C_{0,t}^* = C_1^\ast = 0.4721$   for  $t\geq 1,$  and associated tax rates: $\tau_{0}^\ast=0.3013$ and $\tau_{1}^\ast=0.5273$.
%To induce the continuation Ramsey planner  to confirm the  Ramsey plan, the Ramsey planner repurchases $\{b_{0,t}\}_{t=0}^\infty$ and sells  $\{\hat b_{0,t}\}_{t=0}^\infty$: 
%		$\hat{b} _{0, t} =	{\hat{b} _{0}}  =0.0708 $  for $t \in [0, 1)$ %\; = \; 
% and $\hat{b} _{0, t} = \hat{b} _{1} =  0.1548$  for $t \geq 1$. {This induces  the continuation Ramsey planner at $t=1$  %equal 
% 	 to  confirm the Ramsey plan:  $C_{1, t}^* = C_1^*=0.4721$ for all $t \geq 1$.}  
% 
%	 
% 
%	\begin{figure}[h!]
%		
%		\begin{subfigure}{\textwidth}
%			\caption{Initial debt term structure $b_{0, t}$ and restructured debt $\widehat{b}_{0, t}$}
%	\centering
%	\includegraphics[width=\linewidth]{figure_202403/ts_low_34}	
%		\end{subfigure}
%
%\begin{subfigure}{\textwidth}
%		\caption{Ramsey plan and continuation Ramsey plan. }
%	\centering
%	\includegraphics[width=\linewidth]{figure_202403/ts_low_44}
%\end{subfigure}
%\caption{Implementation works when $b_0=0.3$. Ramsey plan: $ C_0^\ast = 0.6987,  C_1^\ast = 0.4721$,  $\tau_{0}^\ast=0.3013,$ and $\tau_{1}^\ast=0.5273$. Parameter values are: $G_t=0.2, e^{-\rho}=0.5,$ and $\eta=1$. \label{fig:DNYexample}}
%	\end{figure}


%
%\paragraph{DNY (2021) Counterexample} % (1983)} %in Continuous Time} %Specification} %DNY}
% When  $\WJ{b_{0,t}= b_0 =0.5}$ for $t\in [0, 1)$ and $\WJ{b_{0,t}= 0}$ for $t\geq 1$, 			Ramsey planner sets: 	 $C_{0,t}^*= \TS{C_0^\ast = 0.7374}$ \quad for  $t\in [0, 1) ,$  
%			$C_{0,t}^* = \TS{C_1^\ast = 0.1846}$ \quad  for  $t\geq 1$.  %\rightarrow \hat{b} _T  = 0.0787$  %$, which  %
%			 To induce \WJ{continuation Ramsey planner}  to confirm \TS{continuation of Ramsey plan}, Ramsey planner repurchases $\{b_{0,t}\}_{t=0}^\infty$ and sells  $\{\hat b_{0,t}\}_{t=0}^\infty$:
%							$\NW{\hat{b} _{0, t} =	{\hat{b} _{0}} } = \TS{C_0^*}\left(1-(\TS{C_0^*}+G)\right) =0.0462$ 
%				\quad for $t \in [0, 1),$ \bigskip %&& \\ 
%					 $\NW{\hat{b} _{0, t} = \hat{b} _{1}} \; = \;  \left(e^{\rho}-1\right) \left( \frac{\WJ{b_0}}{\TS{C_0^*}} -1 + (\TS{C_0^*}+G)   \right) \TS{C_1^*} =  0.1136$ \quad for $t \geq 1$ 
%%		%	\end{eqnarray*}
%%		
%
%	
%	It is now our job to construct continuation the  Ramsey plan associated with $b_{1,t} = \hat{b}_1 = 0.1136$ for all $t\geq 1$. To do this, we proceed as follows.
%			 Substitute FOC ($1-\tau_{1, t} = C_{1, t}$) for labor into $S(C_{1, t}) = \tau_{1, t} N_{1, t}-G $ to get										 $$\TS{S(C_{1, t}) = C_{1, t}( 1- (C_{1, t}+G) )} 
%				$$
%The  \WJ{continuation IC} is:  $$\int_{1}^\infty e^{-\rho(t-1)} \frac{\hat{b}_1}{C_{1, t}} dt \leq \int_{1}^{\infty} e^{-\rho(t-1)} \frac{S(C_{1, t})}{C_{1, t}}dt . $$
%					From  continuation Lagrangian conclude that   \NW{$C_{1, t}=C_{1, 1}$} and  $\TS{S(C_{1, t})= \hat{b}_1}$ for all $t\geq 1.$
%	Now solve  $ \NW{S(C_{1,1}) = \hat{b}_{1} }$ for $C_{1,1}$. This equation		
%		\textbf{repeat}  $ \TS{S(C_{1,1}) = \hat{b}_1 }$ has two roots:
%		One root   $C_{1,1}=C_1^*=0.1846$ is indeed associated with the continuation of the Ramsey plan.
%	However it is the other root  $C_{1,1}=\widehat{C}_1=0.6154$ that is affiliated with the continuation Ramsey plan
%	The continuation Ramsey planner sets  \NW{$C_{1,1}=\widehat{C}_1=0.6154$} for 
%	\textbf{the following reason.}
%				
%	%	\NW{Make a hyperlink to Laffer curve so it's easy to go back and forth.} 
%
%
%
%
%%	\frame{\frametitle{Case where it works?}
%%	\TS{Neng and Wei: might it be worthwhile to add a numerical example with a lower $b_0$ in which Lucas Stokey's recommendation {\em does} work, and talk readers through exactly what the government is issuing -- linking it to how the government is issuing a consol? Maybe the associated "pictures" could also be enlightening to readers/watchers?}
%%	
%% %   \WJ{Added to the appendix, last four pages. $b_0=0.1$.}
%%
%%%	\NW{How about we  plot for the $b_0=0.3$ case which is just below $b^\ast=0.35$? This shows a bigger effect of tax smoothing motive.} 
%%}
%%
%%	\frame{\frametitle{Laffer curve?}
%%	\TS{Neng and Wei: might it be worthwhile at this point to indicate whether and now the "two roots" in the preceding slide are linked to DNY's "Laffer curve" discussion? Free disposal of my question, as usual.}
%%	\NW{Yes Tom: this is all in our discussion of $b^\ast$. see the hyperlinked slide.} 
%%	\WJ{Tom and Neng: see the appendix slides Laffer curve.}
%%	}
%


%
%
%\noindent\textbf{Lucas-Stokey Implementation doesn't work when $b_0=0.5$}. When  $b_{0,t}= b_0 =0.5$ for $t\in [0, 1)$ and $b_{0,t}= 0$ for $t\geq 1$, the  Ramsey planner sets %(with $\eta=1$, and $T=1$)		
% 	 $C_{0,t}^*= C_0^\ast = 0.7374$  for  $t\in [0, 1) $   and
%	 	 $C_{0,t}^* = C_1^\ast = 0.1846$   for  $t\geq 1$.  In order to devalue  the high initial debt 
%		 $b_0$ by  manipulating the bond price process,    the  Ramsey planner sets the  tax rate   $\tau_{0,t}^* = \tau^*_1= 1-C_1^\ast = 0.8154$  for  $t\geq 1$, which is  above the peak of the Laffer curve.   
%		 
%		 To induce the continuation Ramsey planner  to confirm the 
%		 continuation of the Ramsey plan, Lucas and Stokey advise the Ramsey planner to repurchase $\{b_{0,t}\}_{t=0}^\infty$ and  to construct a restructured debt term structure   $\{\hat b_{0,t}\}_{t=0}^\infty$ in which   $\hat{b}_{0, t} =	\hat{b} _{0}  =0.0462$ for $t \in [0, 1)$ and  $\hat{b}_{0, t} = \hat{b} _{1} =  0.1136$  for $t \geq 1$.
%		 
%		 
%To construct the associated continuation Ramsey plan, first % associated with $b_{1,t} = \hat{b}_1 = 0.1136$ for all $t\geq 1$.
%substitute the first order necessary condition  $1-\tau_{1, t} = C_{1, t}$ for labor into the  primary  surplus formula to get				%	 \begin{eqnarray*} 
%%\text{Surplus:} \, S(C_{T, t})&=&\tau_{T, t} N_{T, t}-G \\ 
%$S(C_{1, t}) = C_{1, t}( 1- (C_{1, t}+G) )\, . 
%$%			\end{eqnarray*} %\bigskip 
%The  implementation constraint for the continuation Ramsey planner is: %(which we also call the IC constraint) is
%  $$\int_{1}^\infty e^{-\rho(t-1)} \frac{\hat{b}_1}{C_{1, t}} dt \leq \int_{1}^{\infty} e^{-\rho(t-1)} \frac{S(C_{1, t})}{C_{1, t}}dt \, .  $$
%Form a    Lagrangian  for the continuation Ramsey planner to obtain   $C_{1, t}=C_{1, 1}$  and  $S(C_{1, t})= \hat{b}_1$ for all $t\geq 1.$ Solve  $ S(C_{1,1}) = \hat{b}_{1} $ for $C_{1,1}$ to  obtain two roots: $C_{1,1}=C_1^*=0.1846$, which describes the continuation of  the Ramsey plan,
%and {$C_{1,1}=\widehat{C}_1=0.6154$, which describes the  continuation Ramsey plan.   To verify that  the continuation Ramsey planner sets {$C_{1,1}=\widehat{C}_1=0.6154$} in order to maximize its objective function,
%form the  Lagrangian: 
%			\begin{equation*}
%				\mathcal{L}_1 = \int_{1}^{\infty} e^{-\rho(t-1)} \bigg[ \log C_{1, t}-\left(C_{1, t}+G\right) +\Lambda_1 \bigg( 1-(C_{1, t}+G) -\frac{\hat{b}_1}{C_{1, t}}\bigg) \bigg]dt .
%			\end{equation*}
%We can show that 		\begin{equation*}
%		U(\widehat{C}_1, \widehat{C}_1+G) > U(C_1^*, C_1^*+G)
%		\end{equation*}
%by using $\widehat{C}_1=0.6154$, $C_1^*=0.1846$, 
%and 		\begin{eqnarray*}
%			-1.8769	 = \int_{1}^{\infty} e^{-\rho(t-1)}U(\widehat{C}_1, \widehat{C}_1+G)dt > 
%		\int_{1}^{\infty} e^{-\rho(t-1)}U(C_1^*, C_1^*+G)dt  =  -2.9926\, . 
%		\end{eqnarray*}
%This completes our verification of \citeauthor{debortoli2021optimal}'s counterexample in which  a Ramsey plan can't be implemented by restructuring  a C\`adl\`ag  $\{b_{0,\,t}; t\geq 0 \}$ {debt  structure}. 		
%
%
%	\begin{figure}[h!]
%	
%	\begin{subfigure}{\textwidth}
%		\caption{Initial debt term structure $b_{0, t}$ and  restructured debt  $\widehat{b}_{0, t}$}
%		\centering
%		\includegraphics[width=\linewidth]{figure_202403/ts_high_34}	
%	\end{subfigure}
%	
%	\begin{subfigure}{\textwidth}
%		\caption{Ramsey plan and continuation Ramsey plan. }
%		\centering
%		\includegraphics[width=\linewidth]{figure_202403/ts_high_44}
%	\end{subfigure}
%	\caption{Implementation doesn't work when $b_0=0.5$. Ramsey plan: $C_0^\ast = 0.7374, C_1^\ast = 0.1846,$ $\tau_{0}^\ast=0.2626,$ and $\tau_{1}^\ast=0.8154$. Continuation Ramsey plan: $\widehat{C}_1=0.6154$ and $\hat{\tau}_{1}=0.3846$. Parameter values are: $G_t=0.2, e^{-\rho}=0.5,$ and $\eta=1$. \label{fig:DNYcounterexample}}
%\end{figure}
%
%%NW{Put a $\Pi^\ast$ and an interest rate graph over time.} 
%
%
%%\NW{Side by side: why our implementation strategy works when L-S doesn't -- and how local commitment kicks in... discuss annuity calculation and all that... why it's feasible under our strategy with Local commitment.}
%
%%\NW{Lucas, Merton  mid 1980s -- Tom will get the precise reference.} 
%
%%\NW{Tom: Yared's graph looks like Hamilton... pls add.} 
%
%  It is not possible to find a counterexample to a counterexample, so we won't try. 
%%  Instead, we'll seek a  minimal extension of  commitments to service continuation government debts that continuation Ramsey planners are bound to honor, an extension that,  in our context,  is faithful to    \citeauthor{LucasStokey1983}'s intention to give the government access to complete markets.    
%%
%%\WJ{Tom: how about the alternative in blue? Feel free to revert to your version. We suggest this because we feel it's more about  commitment mechanism than an extension.} 
%%
%Instead, we'll seek an arrangement with  minimal additional obligations   that continuation Ramsey planners are bound to honor, ones  that,  in our context,  are faithful to    \citeauthor{LucasStokey1983}'s intention to give the government access to complete markets.    
%%\NW{Wei: pls process this.} 
%%
%%one obtains $\hat{b} = \left(e^{\rho }-1\right) \left( \frac{b_{0}}{C_0^*} -1 + (C_0^*+G)   \right) C_1^*\,.$ For $b_0=0.5, C_0^*=0.7374, C_1^*=0.1846, e^{-\rho}=0.5$, we have $\hat{b}_1=0.1136$. Substituting $\hat{b}_1=0.1136$ into the following implementability constraint
%%	one obtains $\hat{b}_0=0.0462$. 
%
%
%%


%
%\TS{Tom Rewrite or drop following couple of paragraphs.}
%
%Our interpretation to this counterexample is that
%Lucas and Stokey restriction on the {\em contractible subspace}: $\{\hat b_{0,t}, t \geq 0\}$ 
%	is \TS{C\`adl\`ag} (right continuous with left limits) process 
%DNY finds that under Lucas and Stokey's assumption about {\em contractible subspace}, \WJ{continuation Ramsey planner} may not choose  \TS{continuation of Ramsey plan}. 
% Restructuring the debt term structure doesn't induce continuation Ramsey planner to implement  Ramsey plan
%

%\newpage 

%\newpage

\subsection{Asset Prices} 
%\NW{Move this part to Lag section?} 
%
%Below we derive asset-pricing implications 
%%report the equilibrium SDF $q_{0,t}$ given in (\ref{foch2}) %= e^{-\rho t} U(C_t)/U(C_0)$, 
%%risk-free rate, and market price of risk 
%associated with the Ramsey plan. 
%\paragraph{SDF, risk-free rate, and  market price of risk.}
%Let $r_{0, t}$ denote the equilibrium interest rate. %Conditional on $s_{t-} =s$,  
Applying Ito's Lemma to $q_{0,t}^*$  given in (\ref{foch2}), we obtain 
\begin{equation} \label{dq} 
d q_{0, t}^* = - \rho q_{0, t-}^* d t + q_{0, t-}^* \frac{dU_{C, t}}{U_{C, t-}} \,,  %= - \rho q_{0, t-} d t + \sum U'(C_t) d C_t  
\end{equation}
where the last term captures the effect of state transitions on the $\{q_{0, t}\}$ process.  
% - r^{*}_{0, t-} M_{t-} d t  - M_{t-}  \sum_{k\neq s_{t-}} \eta(s_{t-}, k)  \left(d {\mathcal J}_t^{(s_{t-}, k)} - \lambda_{s_{t-}, k} d t \right)    
%	d M_t & = & d(-\rho t) + \left[ U_C(C_t, L_t) - U_C(C_{t-}, L_{t-})  \right]d {\mathcal J}_t 
%We also know that 
The equilibrium SDF in a market economy where shocks are driven by jumps can be written  as  (see, e.g., \cite{duffie2010dynamic}): %requires that  
\begin{eqnarray}\label{sdfq}
d q_{0, t}^* = - r^{*}_{0, t-} q_{0, t-}^*dt  -  q_{0, t-}  \sum_{j\neq s_{t-}} \zeta^*(s_{t-}, j)\big(d {\mathcal J}_t^{(s_{t-}, j)} - \lambda_{s_{t-}, j} d t \big)\,,
\end{eqnarray}
%\begin{eqnarray} 
%d M_t = - r^{*}_{0, t-} M_{t-} d t  - M_{t-}  \sum_{k\neq s_{t-}} \eta(s_{t-}, k)  \left(d {\mathcal J}_t^{(s_{t-}, k)} - \lambda_{s_{t-}, k} d t \right)    
%%	d M_t & = & d(-\rho t) + \left[ U_C(C_t, L_t) - U_C(C_{t-}, L_{t-})  \right]d {\mathcal J}_t
%\end{eqnarray}
where $r^{*}_{0, t-}$ is the equilibrium interest rate  and  $\zeta^*(s_{t-}, j)  $ is the equilibrium market price of risk associated with a transition to $j\neq s_{t-}$ from $s_{t-}$.  Equilibrium requires that the expected rate of  change of $q_{0,t}^\ast$ equals  
$-r_{0,t-}^\ast$. Using the martingale representation theorem for a jump process, we thus can write the $q_{0, t}$ process as \eqref{sdfq}, where the conditional expectation of the second term at $t-$ equals zero.  %conditional on the event that a jump  arrives at $t$, bringing the 

%\NW{Wei: please take care of the details and check the signs...} 
Both \eqref{dq} and \eqref{sdfq} must hold, so using restrictions that equilibrium  puts on  the drift of the SDF  and the jump size from state $s_{t-}=i$ to each state $j\neq i$, we obtain: % the following expression for the market price of risk  associated with  transition from state $s_{t-}$ to $k$: %and is 
%\begin{equation}
%\eta(s_{t-}, k) = \left(e^{\kappa^*(s_{t-}, k)}-1 \right)\,,
%\end{equation}
%\NW{We need a different notation from $\kappa$ to avoid confusion with $k$.} 
\begin{equation}
\zeta^*(i, j) = -\big(e^{\eta^*(i, j)}-1 \big)\,,
\end{equation}
where {$\eta^*(i,j)$ for  $i\neq j$} is  %equals logarithmic %given by %the relative jump size which satisfies 
\begin{equation}\label{kappa0} 
\eta^*(i, j)= \ln\frac{U_C^*(C(j), N(j))}{U_C^*(C(i), N(i))} \end{equation}
%\left(e^{\kappa(s_{t-}, s_t)}-1 \right)
and the equilibrium risk-free rate is  
%Using \eqref{foc3}  and \eqref{qXi}, we obtain the risk-free rate
\begin{equation}\label{rstar}
r^{*}_{0, t-} =r^{*}(s_{t-}) = \rho-\sum_{j\neq s_{t-}}\big(e^{\eta^{*}(s_{t-}, j)}-1\big)\lambda_{s_{t-}, j} \,. %+ \frac{d U_{C, t}/dt}{U_{C, t-}}\,.
\end{equation}
Equation (\ref{kappa0}) implies that the market price of risk associated with  transition from state $i$ to $j$ is positive, i.e., $\zeta^*(i,j)>0$, if $\eta^*(i, j)<0$. This requires that the marginal utility of consumption decreases when the economy transitions to state $j$ from state $i$. %as  
Thus, if consumption increases 
when  the state transitions from  $s_{t-}=i$ to $s_t=j$, the zero-coupon bond price   decreases from  $q_{0,t-}$ to $q_{0,t}=q_{0,t-}e^{\eta^*(i, j)} $ where $j\neq i$,  %then 
%%\NW{fix a sign problem here?} 
%\[ q_{0,t} - q_{0,t-}  =  q_{0,t-}  \left(e^{\eta^*(i, j)}-1 \right) \, , 
%\] 
as we can see from the jump shock piece in \eqref{sdfq} and by noting $d {\mathcal J}_t^{(s_{t-}, j)} =1$. 
%\NW{FIX.
Equation (\ref{rstar}) characterizes the equilibrium risk-free rate: the second term captures both expected growth and 
precautionary savings effects that are associated with a jump in  the equilibrium interest rate. %and the second term captures the intertemporal smoothing effect.}} %growth effect. The first term captures the 
% price instantly decreases
%\NW{check the expectation operator?} 

%\[ \left(e^{\kappa^{*}(s_{t-}, k)}-1\right) \] 

%	where $\kappa^*(j, k) (k\neq j)$ is the relative jump size which satisfies 
%	\[ \kappa^*(j, k)= \ln\frac{U_C(C^*(k), L^*(k))}{U_C(C^*(j), L^*(j))}. \]

%	On the other hand, according to equation (\ref{foc3}) on SDF and equilibrium consumption and leisure in Proposition 3.1, the dynamics of equilibrium SDF satisfies
%	\begin{equation}
%	dM_t = -\rho M_{t-} dt + \sum_{s_t\neq s_{t-}}\left(e^{\kappa^*(s_{t-}, s_t)}-1\right)M_{t-}  d {\mathcal J}_t^{(s_{t-}, s_t)}
%	\end{equation}
%	
%	Therefore, the risk free rate $r_t$ is given as

%Let $\Xi_t $ denote the Radon-Nikodym derivative, which is a martingale under the physical measure with mean one: %	If we further introduce 
%	\begin{equation}
%	d\Xi_t  =  - \Xi_{t-} \sum_{k\neq s_{t-}}\left(e^{\kappa^*(s_{t-}, s_t)}-1\right) \left(d {\mathcal J}_t^{(s_{t-}, s_t)} - \lambda_{s_{t-}, s_t} d t \right)
%	\end{equation}
%	We may write the SDF as: %can be written as
%	\begin{equation}
%	\frac{M_t}{M_0} =  \Xi_t e^{-\int_0^{t} r_{0, u} du}.
%	\end{equation}

%\subsection{Dynamics of $\Pi_t$ and Risk Hedging}


%
%From \eqref{martingaleX}, we also know that
%\begin{equation*}
%\mathbb{E}_0 \left[ 	-\int_{0}^t q_{0, u} \left(b_{0, u-} +G_{u-}-\tau_{0, u-} N_{0, u-} \right)du + q_{0, t} \Pi_t  \right] \\
%= - \mathbb{E}_0 q_{0, t}\widehat \Pi_t + \mathbb{E}_0 q_{0, t}\Pi_t =0 \,,
%\end{equation*}
%where 
%\begin{equation}
%\widehat \Pi_t = \int_{0}^t \frac{q_{0, u-}}{q_{0, t}} \left(b_{0, u-} +G_{u-}-\tau_{0, u-} N_{0, u-} \right)du
%\end{equation}
%is the cumulative liabilities up to time $t$. 
%Therefore, we have
%\begin{equation}
%\mathbb{E}_0 q_{0, t}\widehat \Pi_t =\mathbb{E}_0 q_{0, t}\Pi_t\,. 
%\end{equation}
%which suggests that the forward-looking Cumulative Surplus $ \Pi_t$ not necessary equals to the backward-looking Cumulative Liabilities $\widehat\Pi_t$. 

%
%Let $\Xi_t $ denote the Radon-Nikodym derivative, which is a martingale under the physical measure with mean one: %	If we further introduce 
%\begin{equation}
%d\Xi_t  =  - \Xi_{t-} \sum_{k\neq s_{t-}}\eta_{t-}(k) \left(d {\mathcal J}_t^{(s_{t-}, k)} - \lambda_{s_{t-}, k} d t \right)
%\end{equation}
%We may write the debt price as: %can be written as
%\begin{equation}\label{qXi}
%q_{0, t} =  \Xi_t e^{-\int_0^{t} r_{0, u} du},
%\end{equation}
%where $r_{0, t}$ is risk-free rate over interval $[t, t+dt]$.
%Under the Ramsey plan, using first-order condition \eqref{foch2}, we obtain


%Nevertheless, in a scenario that the government chooses exposure to jump risks properly, the forward-looking cumulative surplus perfectly tracks the backward-looking Cumulative Liabilities. Specifically, when $\alpha_t=0$, we have $\Pi_t = \widehat \Pi_t$ for all $t\geq 0$.

%
%Let $\widetilde \lambda_{j, k} = {\lambda}_{j, k}e^{\kappa^*(j, k)}$ denote the risk-adjusted jump intensity from state $j$ to state $k$, $k\neq j$.
%Using the Martingale representation theory, we have
%\begin{align}
%d\Pi_t =   r_{0, t-}\Pi_{t-} dt +\left(b_{0, t-} +G_{t-}-\tau_{0, t-} N_{0, t-} \right)dt % \notag\\
%%	& \quad\quad  + 
%+ \sum_{s_t\neq s_{t-}} \beta_{t-}\Pi_{t-} \left( d {\mathcal J}_t^{(s_{t-}, s_t)} -\widetilde{\lambda}_{s_{t-}, s_t} dt\right)\,.
%\end{align}
%where $\beta_{t-}=\left(e^{-\kappa^*(s_{t-}, s_t)}-1\right)$ is unique for the $\Pi_t$ defined in \eqref{Pi}.
%
%
%
%
%\newpage
%
%\paragraph{Risk-neutral probability measure.} 	
%The SDF defines a risk-neutral probability measure $\mathbb{Q}$. If a change of state in the Markov chain corresponds to a jump in the SDF, then the corresponding large shock carries a risk premium, which leads to an adjustment of the transition intensity under $\mathbb{Q}$ as follows:
%\begin{equation}
%{\lambda}_{j, k}^{\mathbb Q}= e^{\kappa^*(j, k)}\lambda_{j, k}, \, j\neq k.
%\end{equation}
%
%
%
%\paragraph{Term structure.} Let $q_{0,t}$ denote the time-$0$ price of a zero-coupon bond (ZCB) maturing at $t$.  Using the SDF given in (\ref{foc3}), we have 
%\begin{eqnarray} 
%q_{0,t}^{} &=& {\mathbb E}_0 \left(\frac{M_t}{M_0}  \right)  = {\mathbb E}_0 ^{\mathbb Q} \left(  e^{-\int_0^{t} r^{}_{0, u} du}\right) \, . 
%%	\\	y_{0,t} &=& - \frac{\ln q_{0,t}}{ t}  
%\end{eqnarray}
%Let $y_{0,t} $ denote the time-0 yield for  a ZCB with maturity $t$:	\begin{equation}
%y^{}_{0,t} = - \frac{\ln q^{}_{0,t}}{ t} \, . 
%\end{equation}
%%We will plot the yield curve $\{y_{0,t}; t\geq 0\} $. %Then, we can  as a function of $t$. %$y_{0,t}$ generate term structure of interest rate and 
%
%%	M_t & = & e^{-\rho t} U_C(C_t, L_t) / \Phi = e^{-\rho t}\frac{U_C(C_t, L_t)}{U_{C}(C_0, L_0)}   \\
%%The time-0 price of a perpetual risk-free  debt with unit coupon rate is %fixed cash flow at the rate of one?
%%\begin{eqnarray} 
%%\int_0^{\infty}  q_{0,t}\,  d t  \, . 
%%\end{eqnarray}
%
%





%
%From \eqref{martingaleX}, we also know that
%\begin{equation*}
%\mathbb{E}_0 \left[ 	-\int_{0}^t q_{0, u} \left(b_{0, u-} +G_{u-}-\tau_{0, u-} N_{0, u-} \right)du + q_{0, t} \Pi_t  \right] \\
%= - \mathbb{E}_0 q_{0, t}\widehat \Pi_t + \mathbb{E}_0 q_{0, t}\Pi_t =0 \,,
%\end{equation*}
%where 
%\begin{equation}
%\widehat \Pi_t = \int_{0}^t \frac{q_{0, u-}}{q_{0, t}} \left(b_{0, u-} +G_{u-}-\tau_{0, u-} N_{0, u-} \right)du
%\end{equation}
%is the cumulative liabilities up to time $t$. 
%Therefore, we have
%\begin{equation}
%\mathbb{E}_0 q_{0, t}\widehat \Pi_t =\mathbb{E}_0 q_{0, t}\Pi_t\,. 
%\end{equation}
%which suggests that the forward-looking Cumulative Surplus $ \Pi_t$ not necessary equals to the backward-looking Cumulative Liabilities $\widehat\Pi_t$. 


%Let $\Xi_t $ denote the Radon-Nikodym derivative, which is a martingale under the physical measure with mean one: %	If we further introduce 
%\begin{equation}
%d\Xi_t  =  - \Xi_{t-} \sum_{k\neq s_{t-}}\eta_{t-}(k) \left(d {\mathcal J}_t^{(s_{t-}, k)} - \lambda_{s_{t-}, k} d t \right)
%\end{equation}
%We may write the debt price as: %can be written as
%\begin{equation}\label{qXi}
%q_{0, t} =  \Xi_t e^{-\int_0^{t} r_{0, u} du},
%\end{equation}
%where $r_{0, t}$ is risk-free rate over interval $[t, t+dt]$.
%Applying Ito's lemma, we obtain
%\begin{eqnarray}\label{sdfq}
%d q_{0, t} = - r_{0, t-} q_{0, t-}dt +  \sum_{k\neq s_{t-}} \eta_{t-}(k) q_{0, t-} \left(d {\mathcal J}_t^{(s_{t-}, k)} - \lambda_{s_{t-}, k} d t \right)\,,
%\end{eqnarray}
%Under the Ramsey plan, using first-order condition \eqref{foc3}, we obtain
%\begin{equation}
%\eta_{t-}(k) = \left(e^{\kappa^*(s_{t-}, k)}-1 \right)\,,
%\end{equation}
%where $\kappa^*(j, k) (k\neq j)$ is the relative jump size which satisfies 
%\[ \kappa^*(j, k)= \ln\frac{U_C^*(C(k), N(k))}{U_C^*(C(j), N(j))}. \]
%%\left(e^{\kappa(s_{t-}, s_t)}-1 \right)


%Nevertheless, in a scenario that the government chooses exposure to jump risks properly, the forward-looking cumulative surplus perfectly tracks the backward-looking Cumulative Liabilities. Specifically, when $\alpha_t=0$, we have $\Pi_t = \widehat \Pi_t$ for all $t\geq 0$.

%\paragraph{Dynamics of $\Pi_t$}

%Let $\widetilde \lambda_{i, j} = {\lambda}_{i, j}e^{\kappa^*(i, j)}$ denote the risk-adjusted jump intensity from state $i$ to state $j$, $j\neq i$.
%Using the Martingale representation theorem, we obtain
%\begin{align}
%d\Pi_t =   r_{0, t-}\Pi_{t-} dt +\left(b_{0, t-} +G_{t-}-\tau_{0, t-} N_{0, t-} \right)dt % \notag\\
%%	& \quad\quad  + 
%+ \sum_{j\neq s_{t-}} \beta(s_{t-}, j)\Pi_{t-} \left( d {\mathcal J}_t^{(s_{t-}, j)} -\widetilde{\lambda}_{s_{t-}, j} dt\right)\,.
%\end{align}
%where $\beta(i, j)=\left(e^{-\kappa^*(i, j)}-1\right)$ for $j\neq i$ is unique for the $\Pi_t$ defined in \eqref{Pi}.
%

%\newpage
%
%\section*{\WJ{Cumulative Liabilities $\Pi_t$}} 
%
%\WJ{XXX} 
%
%\paragraph{Intertemporal budget constraints.}  
%
%
%
%In a competitive equilibrium defined in Definition \ref{defn1}, %Under a  plan featuring a stochastic $\{C_{0, t} \}_{t\geq 0}$ process, the associated 
%there exists a stochastic  bond price process, $\{q_{0,t};\, t \ge 0\}$, that satisfies $q_{s,t} = e^{-\int_s^t r_{0,u}\, du}$. 
%%
%%\NW{Recall that \begin{equation}\label{eq:dPi}
%%d\Pi_t = r_{0,t}\, \Pi_t\, dt + (b_{0-,t}+G_t - \tau_{0,t}\, N_{0,t})\, dt.
%%\end{equation}
%%}
%%\eqref{eq:dPi}
%Let~$B_t$ denote the government's instantaneous debt balance at~$t$. 
%Over an infinitesimal time interval $[t, t+dt)$, %are $r_{0,t}^*\, B_t\, dt$. 
%the government  issues instantaneous debt and term debt to service  its  debts and  to  finance the continuation of its cumulative liabilities.  %$\Pi_t^*$, as follows: %
%Consequently 
%\begin{align}\label{eq:dPi1}
%%d\Pi_t^* -r_{0, t}^*\Pi_t^*dt = 
%& \underbrace{(G_t - \tau_{0,t} N_{0,t} ) d t}_{\text{primary deficits}} + \, %
%\underbrace{\biggl( b_{0-,t} + \int_{u<t}\partial_u\, b_{u,t}\, du\biggr)  \, dt}_{\text{maturing term debt repayments}}
%\notag \\
%&\quad \quad  \quad \quad =   \underbrace{dB_t -r_{0, t}B_tdt}_{\text{net instantaneous-debt issuance}}
%+ %
%\underbrace{\biggl(\mathbb{E}_t\int_{u > t} q_{t,u}\, \partial_t\, b_{t,u}\, du\biggr) dt}_{\text{gross term-debt issuance}} 
%\, . 
%%_{\text{maturing term debt repayment}} .
%\end{align}
%%Equation \eqref{eq:dPi} implies that the left side of \eqref{eq:dPi1}, which is the  net change of $\Pi_t$, 
%%%netting  the associated interest costs $r \Pi_t^* d t$.} \WJ{Equation \eqref{eq:Pi_def} implies that the left side of~\eqref{eq:dPi_star} also 
%%equals  $(b_{0-,t} + G_t - \tau_{0,t} N_{0,t} ) d t.$ Substituting this into \eqref{eq:dPi1} and re-arranging  yields: 
%%\begin{align}\label{eq:dPi2}
%%%d\Pi_t^* -r_{0, t}^*\Pi_t^*dt = 
%%& \underbrace{(G_t - \tau^\ast_{0,t} N_{0,t} ) d t}_{\text{primary deficits}} + \, %
%%\underbrace{\biggl( b_{0-,t} + \int_{s<t}\partial_s\, b_{s,t}\, ds\biggr)  \, dt}_{\text{maturing term debt repayments}}
%%\notag \\
%% &\quad \quad  \quad \quad =   \underbrace{dB_t -r_{0, t} B_tdt}_{\text{net instantaneous-debt issuance}}
%% + \quad
%% \underbrace{\biggl(\int_{s > t} q_{t,s}\, \partial_t\, b_{t,s}\, ds\biggr) dt}_{\text{gross term-debt issuance}} 
%%\, . 
%%%_{\text{maturing term debt repayment}} .
%%\end{align}
%Relation \eqref{eq:dPi1} describes the government's intertemporal budget constraint.  
%%Note that the initial term debt $(b_{0-,t} $ appears in the net change of $\Pi_t^*$. 
%Uses of government funds over $[t, t+dt)$ include  primary deficits (the first term) and  repayments for all maturing debts (the second  term), the latter of which  includes  the initial term debt $b_{0-,t}$ as well as all other term debts due at $t$.\footnote{Here  
%{$\partial_s\, b_{s,t}$} denotes the flow rate of coupon obligations issued at time~$u$ and maturing at time~$t$, so $(\int_{s<t} \mathbb{E}_t \partial_s\, b_{s,t}\, ds)\, dt$  represents  all outstanding term debts that were issued after $t=0-$ and are due at time~$t$.} 
%To finance this outflow, the government issues instantaneous debt (netting interest payments), $dB_t -r_{0, t}B_tdt$, and term debts of all possible maturities: %$\mathbb{E}_t\int_{u > t} q_{t,u}\, \partial_t\, b_{t,u}\, du$. 
%$\big(\mathbb{E}_t\int_{u > t} q_{t,u}\, \partial_t\, b_{t,u}\, du\big) dt$, where the expectation operator  sums over all  $\{s_u; u\geq t \}$ paths. %This explains t
%%on the right side of \eqref{eq:dPi1}.  %for this gross term-debt issuance term is to
%	%The second term in~\eqref{eq:dPi_star} describes net financing from newly issued term debt over the interval~$dt$, while the first term, $dB_t$, represents changes in the instantaneous debt balance that finance residual shortfalls not covered by term-debt adjustments.
%%$dB_t -r_{0, t}^*B_tdt$, and term debts of all possible maturities: $\int_{s > t}  q_{t,s}\, \partial_t\, b_{t,s}\, ds$ %\NW{Point out LS has no $B$ and the difference between DT and CT formulations?} 
%%The second term in~\eqref{eq:dPi_star} describes net financing from newly issued term debt over the interval~$dt$, while the first term, $dB_t$, represents changes in the instantaneous debt balance that finance residual shortfalls not covered by term-debt adjustments.
% %, aggregated over issuances from dates~0 through~$t$
%
%%	Uses include  primary deficits (the first term) and  repayments for all maturing debts (the second  term), the latter of which  includes  the initial term debt $b_{0-,t}$ as well as all other term debts due at $t$.\footnote{Here  
%%		{$\partial_u\, b_{u,t}$} denotes the flow rate of coupon obligations issued at time~$u$ and maturing at time~$t$, so $(\int_{u<t} \partial_u\, b_{u,t}\, du)\, dt$  represents  all outstanding term debts that were issued after $t=0-$ and are due at time~$t$.} 
%%	Sources include  (net) issuance of instantaneous debt, 
%	
%Next, we show there are infinitely many debt management policies that satisfy \eqref{eq:dPi1}. 
%%To develop a key accounting variable, we re-arrange \eqref{eq:dPi_star1}:
%\begin{lemma} \label{Lemma1}
%	For any debt management policies $\{B_t, b_{t, u}; u>t, t\geq 0\}$ that satisfies the intertemporal budget constraints \eqref{eq:dPi1}, the implementability condition \eqref{impletconstraint} always hold.
%\end{lemma}
%
%
%\paragraph{Indeterminacy of debt management policies.} 
%For a  given plan $\{C_{0,t}; t\geq 0\}$, 
%%\NW{Add a lemma and refer to appendix.}  While 
%there are infinitely many debt management policies that a government can use subject to the intertemporal budget constraints \eqref{eq:dPi1}. 
%
%$\{B_t, b_{t, u}; u>t, t\geq 0\}$ are feasible. However, as we show below 
%
%
%(see details of proof in Appendix XX), debt management policies are indeterminant and  allocations only need to satisfy the implementability condition \eqref{impletconstraint}.
%
%
%
%\paragraph{Introducing government's cumulative liabilities: $\{\Pi_t; t\geq 0\}$.} 
%We can re-write \eqref{eq:dPi1} as follows: 
%\begin{align}\label{eq:dPi_star2}
%%d\Pi_t^* -r_{0, t}^*\Pi_t^*dt = 
%(b_{0-,t}+G_t - \tau_{0,t} N_{0,t} ) d t  \, %
%=   dB_t -r_{0, t}B_tdt
%+ %
%\biggl(\mathbb{E}_t\int_{u > t} q_{t,u}\, \partial_t\, b_{t,u}\, du\biggr) dt-\biggl(   \int_{u<t}\partial_u\, b_{u,t}\, du\biggr)  \, dt
%\, . 
%%_{\text{maturing term debt repayment}} .
%\end{align} 
%While debt management policies are indeterminate, the left side of \eqref{eq:dPi_star2} is unique under a plan $\{C_{0, t}; t\geq 0\}$. 
%Next, we construct a process that tracks the cumulative effects of the outflow on  the left side of \eqref{eq:dPi1} from time $0$ to $t$. 

%
%
%Let $\Pi_t$ denote this process: 
%%
%%
%%%cumulative 
%%
%%%Now consider the debt-management policy 
%%
%%While \eqref{eq:dPi_star2} captures government's per-period liability based on initial term debt, we define
%%the government's time $t$ cumulative liabilities 
%\begin{equation}\label{eq:Pi_def}
%\Pi_t = \int_0^t \frac{q_{0,s}}{q_{0,t}}\bigl(b_{0-,s}+G_s - \tau_{0,s}\, N_{0,s}\bigr)\, ds, \qquad \Pi_0 = 0\, , 
%\end{equation} %show later $\Pi_t$ serves as a key variable to define a state variable of the problem.  
%%\NW{It is worth pointing out that leaving all existing term debts untouched is an idea that underpins our analysis, unlike in \citeauthor{LucasStokey1983}'s and \citeauthor*{debortoli2021optimal}'s  analyses.} 

%We have assumed that $\{G_t\}$ and $\{b_{0-,t}\}$ are continuous processes, so optimal allocations $\{C_{0,t}, N_{0,t};\, t \ge 0\}$ are also continuous. That implies that the zero-coupon bond price~$q_{0,t}$ is continuous in~$t$ and %. We thus can write the $\{q_{0,t}\}$ process as:
%that  $dq_{0,t} = -r_{0,t}\, q_{0,t}\, dt$, where $\{r_{0,t}\}$ is the continuously compounded risk-free interest-rate process. Differentiating~\eqref{eq:Pi_def} with respect to~$t$, we obtain the evolution of $\Pi_t$.
%\begin{equation}\label{eq:dPi}
%d\Pi_t = r_{0,t}\, \Pi_t\, dt + (b_{0-,t}+G_t - \tau_{0,t}\, N_{0,t})\, dt.
%\end{equation}

%We then show that the debt restructuring strategy $\{\partial_t b_{t, u}; u>t, t\geq 0\}$ follows Ricardian equivalence. 




%Here $(b_{0-,t}+G_t-\tau_{0,t}\, N_{0,t})\, dt$ is the (flow) increase of the government's liability based on the initial term debt over a small time interval~$dt$.
%
%\begin{equation}\label{eq:dPi}
%d\Pi_t -r_{0, t}\Pi_tdt = %\underbrace
%{dB_t -r_{0, t}B_tdt} %_{\text{net instantaneous debt issuance}}
% + %\underbrace
% {\biggl(\int_{s > t} q_{t,s}\, \partial_t\, b_{t,s}\, ds\biggr) dt}%_{\text{gross term debt issuance}} 
% -\, %\underbrace
% {\biggl( \int_{s<t}\partial_s\, b_{s,t}\, ds\biggr)  \, dt}\, . 
%%_{\text{maturing term debt repayment}} .
%\end{equation}
%One way to see \eqref{eq:dPi_star} is as follows. 
%where $b_{t-,t} dt $ is the total term debt due over the small time interval $[t, t+dt)$:
%\begin{equation}\label{eq:dPi_star} %) dt
%b_{t-,t} = \underbrace{\biggl(b_{0-,t} +  \int_0^{t-}\partial_u\, b_{u,t}\, du\biggr)}_{\text{term debt due at }t}
%\end{equation}  the instantaneous government budget constraint~

%\TS{Neng and Wei: here we need a definition of ''the contractible subspace'' in precise mathematical terms that we can also relate to Lucas and Stokey's possibly different language.}


%
%
%In this part, we study how a government manages its dynamic budget constraint for a given feasible allocation $\{C_{0, t}, N_{0, t}, \tau_{0, t}, q_{0, t} \}_{t\geq 0}$. Over instantaneous time interval $[t, t+dt)$ absent state-transitions, by equating uses and sources  of government funds, we have
%\begin{align}\label{eq:dPi_star1}
%%d\Pi_t^* -r_{0, t}^*\Pi_t^*dt = 
%& \underbrace{(G_t - \tau_{0,t} N_{0,t} ) d t}_{\text{primary deficits}} + \, %
%\underbrace{\biggl( b_{0-,t} + \int_{u<t}\partial_u\, b_{u,t}\, du\biggr)  \, dt}_{\text{maturing term debt repayments}}
%\notag \\
%&\quad \quad  \quad \quad =   \underbrace{dB_t -r_{0, t}B_tdt}_{\text{net instantaneous-debt issuance}}
%+ %
%\underbrace{\biggl(\mathbb{E}_t\int_{u > t} q_{t,u}\, \partial_t\, b_{t,u}\, du\biggr) dt}_{\text{gross term-debt issuance}} 
%\, . 
%%_{\text{maturing term debt repayment}} .
%\end{align}
%%Note that the initial term debt $(b_{0-,t} $ appears in the net change of $\Pi_t^*$. 
%


%\newpage
%
%\section{Proposal} 
%
%under Ramsey with an eye about implementation.... 
%
%\begin{itemize}
%	\item LS: once with starred Ramsey plan, then open market operations: 
%	\begin{equation}\label{gbudgetc} 
%\mathbb{E}_0 \left[ \int_{0}^{\infty} {b}_{0, t} q_{0, t}^\ast  dt \right] =  {\mathbb E}_0 \left[\int_0^{\infty}  q_{0, t}^\ast \left( \tau_{0, t}^\ast N_{0, t}^\ast -G_{t} \right) d t\right]\,,  %\leq  {\mathbb E}_0 \left[\int_0^{\infty}  q_{0, t} \left( \tau_{0, t} N_{0, t} -G_{t} \right) d t\right] \, ,
%\end{equation} 
%	then when it comes to time-$t$, LS proposed an alchemy style strategy: 
%	come up with a term debt (consol or whatever) that works from $t$ onward.... keep iterating -- a lot of wishful thinking -- the fact that it works at $t$ does not imply you can have all worked out in all future periods 
%we can also put HH's time-0 budget equation here... All goes to LS
%	\begin{itemize} 
%	\item (this point is just an observation) one strategy that works obviously but is not implementable in general is "balanced the budget at all $t$: 
%		\begin{equation}\label{bb} 
%{b}_{0, t}  + G_{t} = \tau_{0, t}^\ast N_{0, t} ^\ast\, ,
%\end{equation} 
%\end{itemize} 
%	\item another strategy 
%	\begin{eqnarray}\label{gbudgetc} 
%\mathbb{E}_0 \left[ \int_{0}^{\infty} {b}_{0, t} q_{0, t}^\ast  dt \right] &  = & \mathbb{E}_0 \left[ \int_{0}^{u} b_{0-, t} q_{0, t} ^\ast dt + \int_{u}^{\infty} b_{0-, t} q_{0, t} ^\ast dt \right] \notag \\
%&  = & {\mathbb E}_0 \left[\int_0^{\infty}  q_{0, t}^\ast \left( \tau_{0, t}^\ast N_{0, t}^\ast -G_{t} \right) d t\right]  %\leq  {\mathbb E}_0 \left[\int_0^{\infty}  q_{0, t} \left( \tau_{0, t} N_{0, t} -G_{t} \right) d t\right] \, ,
%\end{eqnarray} 
%{\bf add PV of primary surplus here.} 
%THen keep 
%\[ {b}_{0, t}  = b_{0-, t}\] for all $t$, %$t\leq u$ 
%then track cumulative liabilities under Ramsey: starred on the right side too 
%
%\begin{equation} \label{Pi}
%\Pi_t ^\ast %= \int_{0}^t\,   e^{\int_s^{t} r_ud u} \left(b_{0, s} +G_s-\tau_{0, s} N_{0, s} \right)ds 
%= \int_{0}^t\, \frac{q_{0, u}}{q_{0, t}} \left(b_{0-, u} +G_u-\tau_{0, u} N_{0, u} \right)du \, , \quad  \Pi_0 = 0\, . 
%\end{equation}   





%\end{itemize} 
	
%\newpage

\section{A Dynamic Program}\label{sec:bellmanramsey} %Via Dynamic Programming} %Optimal Commitment Policy under 
%\NW{Time inconsistency arises because of the multiple roots of Lagrangian multipliers associated with the implementability condition.}

%\subsection{Pure Short-term Debt Roll-over Policy \label{sec:irrelevance}}

%Note that the implementability condition (\ref{impletconstraint}) that time$-0$ government must comply to is independent of debt roll-over policy for the future government. 

{Via the following four steps, we  the Ramsey problem recursively: %representation of a Ramsey plan. First,
1. For a given consumption process, we define  a process, $\{\Pi_t;\, t \ge 0\}$ that tracks the government's cumulative liabilities in the absence of any debt restructuring. 2. We use the  $\{\Pi_t;\, t \ge 0\}$ process to construct a state variable process $\{X_t;\, t \ge 0\}$. 3. We frame  a dynamic program with $\{X_t;\, t \ge 0\}$ as a state variable and a  state $s_t$ dependent value function, $V(X_t,s_t)$ that  is additively separable: $V(X_t,s_t) = -\Phi_0\, X_t + H(s_t;\, \Phi_0),$ where  $\Phi_0$ is an undetermined  scalar.  4. We compute $\Phi_0$. 
} %starts from  a prescribed initial condition~$X_0$   at time $t=0$ 
%We can  use  dynamic programming to formulate and compute a Ramsey plan.
%method to solve for the optimal smooth equilibrium based on the pure short-term debt rolling strategy. 
%This section presents a recursive representation of a Ramsey plan. 
The  dynamic program  sets the stage for   a  state-contingent debt management strategy that  implements the Ramsey plan. 
 
 
 
\subsection{{Cumulative Liabilities with no Debt Restructuring}}

%A key object in our analysis is . That is, 
%Given initial debt service flow $\{b_{0-,s},\, s \ge 0\}$, in the absence of any debt restructuring the government's time $t$ liabilities are
%\begin{equation}\label{eq:Pi_def}
%\Pi_t = \int_0^t \frac{q_{0,s}}{q_{0,t}}\bigl(b_{0-,s}+G_s - \tau_{0,s}\, N_{0,s}\bigr)\, ds, \qquad \Pi_0 = 0.
%\end{equation} Here $(b_{0-,s}+G_s-\tau_{0,s}\, N_{0,s})\, ds$ is the (flow) increase of the government's liability over a small time interval~$ds$;
%% without touching the initial term debt at all; 
%the ratio $q_{0,s}/q_{0,t}$ compounds this flow's contribution to the government's time-$t$ liability stock.\footnote{When the interest rate is a constant~$r$ over $(s,t)$, $q_{0,s}/q_{0,t} = e^{r(t-s)} > 1$.} 
%\NW{It is worth pointing out that leaving all existing term debts untouched is an idea that underpins our analysis, unlike in \citeauthor{LucasStokey1983}'s and \citeauthor*{debortoli2021optimal}'s  analyses.} 
%
%We have assumed that $\{G_t\}$ and $\{b_{0-,t}\}$ are continuous processes, so optimal allocations $\{C_{0,t}, N_{0,t};\, t \ge 0\}$ are also continuous. That implies that the zero-coupon bond price~$q_{0,t}$ is continuous in~$t$ and %. We thus can write the $\{q_{0,t}\}$ process as:
%that  $dq_{0,t} = -r_{0,t}\, q_{0,t}\, dt$, where $\{r_{0,t}\}$ is the continuously compounded risk-free interest-rate process. Differentiating~\eqref{eq:Pi_def} with respect to~$t$, we obtain:
%\begin{equation}\label{eq:dPi}
%d\Pi_t = r_{0,t}\, \Pi_t\, dt + (b_{0-,t}+G_t - \tau_{0,t}\, N_{0,t})\, dt.
%\end{equation}
%
%\NW{Next, we use this $\{\Pi_t;\, t \ge 0\}$ process to construct a state variable to make the Ramsey planner's problem recursive. 
%} 

Given a  competitive equilibrium, % (see Definition \ref{defn1}),}  
%For an admissible consumption process $\{C_{0,t}\}$, %For  debt service flow  $\{b_{0, s}, s\geq 0\}$,   implied by \eqref{foch1}-\eqref{foch2}
%there exist a tax process $\{\tau_{0, t}\}$,  a state-contingent bond price process $\{q_{0, t}\}$, and a labor-supply process $\{N_{0, t}\}$. Using these processes and together with the exogenous term debt and government spending processes, $\{b_{0,t}\}$  and $\{G_t\}$,  %consider the following debt management strategy 
 we construct  % the following  $\{\Pi_t\}$ process: %a measure for the government's cumulative liability: 
%the funding gap satisfies
\begin{equation} \label{Pi}
\Pi_t %= \int_{0}^t\,   e^{\int_s^{t} r_ud u} \left(b_{0, s} +G_s-\tau_{0, s} N_{0, s} \right)ds 
= \int_{0}^t\, \frac{q_{0, u}}{q_{0, t}} \left(b_{0-, u} +G_u-\tau_{0, u} N_{0, u} \right)du \, , \quad  \Pi_0 = 0\, .
\end{equation}   
%\begin{equation} \label{Pi}
%\Pi_t %= \int_{0}^t\,   e^{\int_s^{t} r_ud u} \left(b_{0, s} +G_s-\tau_{0, s} N_{0, s} \right)ds 
%=  \left[ \int_{0}^t\, \frac{q_{0, u-}}{q_{0, t}} \left(b_{0, u-} +G_{u-}-\tau_{0, u-} N_{0, u-} \right)du \right]\, , \quad  \Pi_0 = 0\, . 
%\end{equation} 
%\NW{Do we need the $u-$?} \WJ{Edit below once we agree on notations.} 
 %  is the (flow) increase of the government's liability over a small time interval~$ds$;
% without touching the initial term debt at all; 
Evidently, the ratio $q_{0,u}/q_{0,t}$ compounds the contribution of the flow $(b_{0-,u}+G_u-\tau_{0,u}\, N_{0,u})\, du$ to the government's time-$t$ liability stock.\footnote{When the interest rate is a constant~$r$ over $(u,t)$, $q_{0,u}/q_{0,t} = e^{r(t-u)} > 1$.} %As we will Here %In (\ref{Pi}),
% $\left(b_{0, u} +G_u-\tau_{0, u} N_{0, u} \right)du $ is the   (flow) increase of the government's liability over a small time interval $du$ without touching term debt  at all and the ratio ${q_{0, u}}/{q_{0, t}} $ compounds  this flow's contribution from $t=0$ onward to the government's time-$t$ liability stock.\footnote{When the interest rate is a constant $r$ over $(s,t)$, then  $\frac{q_{0, s}}{q_{0, t}}=e^{r(t-s)}>1$.} %  $ e^{\int_s^{t} r_ud u}$  compounds interest  from $s$ to $t$.  
%Evidently, $\Pi_t$ is the time$-t$ value of the government's cumulative primary deficits. 
The $\{\Pi_t\}$ process tracks the cumulative liabilities with no debt restructuring at all from $0$ to $t$.  
%As we show below, this process plays a pivotal role in our recursive formulation of the Ramsey plan. 
%\WJ{That is, $\Pi_t$ measures time-$t$ cumulative liabilities in the absence of term-debt adjustments,    defined in equation~\eqref{eq:Pi_def}, evaluated under the Ramsey plan.}  
% the  $\Pi_t$ process as follows: 
\begin{lemma}\label{lemmaPi0}
	$\Pi_t$ defined in \eqref{Pi} evolves according to
	\begin{align}\label{dPi0}
	d\Pi_t  
	\quad =& \quad r^{}_{0, t-}\Pi_{t-}dt 
	+ \left(b_{0-, t-} + G_{t-} - \tau_{0, t-} N_{0, t-} \right)dt   \notag \\
	&\quad + \sum_{j\neq s_{t-}} \left(e^{-\eta(s_{t-}, j)} - 1 \right)\Pi_{t-} 
	\big(d {\mathcal J}_t^{(s_{t-}, j)} 
	- \lambda_{s_{t-}, j} e^{\eta(s_{t-}, j)} dt \big) \, ,
	\end{align}
	where $\eta(s_{t-}, s_t) = \ln\frac{U_{C, t}}{U_{C, t-}}$.  %\NW{Make it general: }
\end{lemma}


%We make two observations regarding 
In evolution equation  \eqref{dPi0},   $ \lambda_{s_{t-}, j} e^{\eta(s_{t-}, j)} $ in the last term is the `risk-adjusted' transition rate\footnote{Technically, $ \lambda_{s_{t-}, j} e^{\eta(s_{t-}, j)} $ is the regime-transition rate under the risk-neutral measure; see, e.g., \cite{duffie2010dynamic}, for a textbook treatment.} %which differs from the rate $ \lambda_{s_{t-}, j} $
 at which the state transitions from $s_{t-}$ to $j\neq s_{t-}$ under the physical measure. The ratio between the risk-adjusted and original transition rates, $ e^{\eta(s_{t-}, j)} $,  captures the `price' of transition risk   from state $s_{t-}$ to $j$. The representative consumer demands a positive risk premium when a  state transition increases the marginal utility of consumption:  $\eta(s_{t-}, j)>0$. %carries a positive price of risk.  
When  the state transitions from $s_{t-}$ to $j\neq s_{t-}$, i.e., when  $d {\mathcal J}_t^{(s_{t-}, j)} =1$, the cumulative liability jumps  from $\Pi_{t-}$ to $\Pi_{t-} e^{-\eta(s_{t-}, j)} =  \Pi_{t-} \frac{U_{C, t-}}{U_{C, t}}= \Pi_{t-} \frac{q_{0, t-}}{q_{0, t}} = \Pi_t .$ Consequently,  while $  \Pi_t$ and $q_{0, t}$ can both jump, their product $q_{0, t}  \Pi_t$ is  continuous at all $t$. %
%This part corresponds to the `quantity' of this regime-transition risk. 
%
%In summary, both the `price' and the `quantity' of risk arise from the household's demand for bearing systematic regime-transition risk in equilibrium.

%\newpage 

%Next, we use the $\{\Pi_t\}$ process to introduce a state variable $\{X_t; t\geq 0\}$.

% in subsection \ref{sec:X}, then formulate a dynamic programming problem  for a given $X_0$ in subsection \ref{sec:DPdefine}, and finally% ({\bf  DP Problem}). 
% characterize the Ramsey plan in subsection \ref{sec:DPRamsey}. 


% using the solution of this {\bf  DP Problem}. %we fully characterize the Ramsey plan.} %er's time-0 decision involving $X_0$.} %solve for $X_0$ and $C_0$ 


%Bellmanize  and  analytically characterize the  Ramsey plan}  given , we solve for a dynamic programming problem with $X_t$ and time $t$ being the state variables. Then, 

%Anticipating that it is sometimes more convenient and intuitive to work with the risk-free rate 

%\paragraph{Cumulative primary deficits.} 

%\NW{Discuss state contingent debt and hedging equivalents.... link to the literature.} 

%Let $r_t dt$ denote a % instantaneous
%interest paid  over a small time interval $dt$. % where $r_t$ will be endogenously determined. 
%The  time-$0$  price of a  zero-coupon bond with maturity $t$ is %$q_{0,t}$ is given by %$q_{0, t}$ implies that the $r_t$ sa
%\begin{equation} \label{qprice}
%	q_{0,t} = e^{-\int_0^{t} r_s d s}\, .
%\end{equation}
\subsection{A State Variable $X_t$} \label{sec:X}
%Ramsey plan 
%We have shown that Ramsey allocations are Markovian in $s_t$ under Assumption \ref{assump1}. Thus, when constructing the state variable  $X_t$ in this section, we focus on the set of admissible Markovian consumption process 
%%Recall that under Assumption \ref{assump1}, there exists a Markov process for consumption $\{C_{0,t}=C(s_t)\}$,  For a given consumption process 
%$\{C_{0,t}\}$.  %For  debt service flow  $\{b_{0, s}, s\geq 0\}$,   implied by \eqref{foch1}-\eqref{foch2}
%

The continuity of the $\{q_{0, t}  \Pi_t\}$ process induces us to define %Next, we introduce $X_t$  which is $\Pi_t$ multiplied by the household's marginal utility of consumption $U_{C, t}$: %multiplied by the cumulative primary deficits 
\begin{align}\label{X}
X_t = \Pi_t U_{C, t}\, .
\end{align}
%In effect, 
%We later show that $X_t$ is the 
Evidently,  $X_t$ incorporates both information about  $\Pi_t$ and the marginal utility of consumption  $U_{C, t}$. 
%The reason that $X_t$ allows us to recursify our problem is can be viewed as the new worth of cumulative primary deficit measured by the marginal utility up to time $t$.  
%To highlight the mechanism, 
Combining \eqref{X}, \eqref{Pi} and \eqref{foch2} yields
\begin{equation}
X_t =\Pi_t e^{\rho t}q_{0, t}U_{C, 0} = e^{\rho t} U_{C, 0} \int_{0}^t\, {q_{0, u}} \left(b_{0-, u} +G_u-\tau_{0, u} N_{0, u} \right)du \,.
\end{equation}
Differentiating both sides of the above equation  and using $C_t = N_t - G_t$, $e^{\rho t} U_{C, 0} q_{0, t}=U_{C, t}$, and $(1-\tau_{0, t}) U_{C, t}  = -  U_{N, t}$,  implied by \eqref{resourceconstraint} and \eqref{foch1}-\eqref{foch2}, we have
\begin{eqnarray}\label{dX} 
dX_t  &= & \rho  X_t dt +  e^{\rho t} U_{C, 0} q_{0, t} \left(b_{0-, t} +G_t-\tau_{0, t} N_{0, t} \right)dt   \notag\\
&= & \rho  X_t dt +  U_{C, t} \left(b_{0-, t} +G_t - N_{0,t} + (1-\tau_{0, t}) N_{0, t} \right)dt   \notag\\& = & \left(\rho X_{t}+b_{0-, t}U_{C, t}-C_{0, t}U_{C, t} -N_{0, t}U_{N, t}  \right)dt\,.
\end{eqnarray}


%Because of jump shocks, $\{C_{0, t}; t\geq 0\}$ are discontinuous. 
%Let   $\mathcal T = \{0<t_1<t_2<\cdots \}$ denote a set of countable (possibly infinite)  points where the $\{C_{0, t}; t\geq 0\}$ process jumps.
%
%\paragraph{What happens absent jumps.} 
%For $t\notin \mathcal T$ where $\{\Pi_t\}$ does not jump, $q_{0, t}$ and $\Pi_t$ evolve continuously: 
%% is  continuous so that 
% $d q_{0, t}  = -r_{0,t} q_{0, t} dt $ and % as follows: 
%\begin{equation} \label{eqPi}
%%\left\{
%%\begin{aligned}
%d\Pi_t  %=  -\Pi_t\frac{dq_{0, t}}{q_{0, t}} + \left(b_{0, t} +G_t-\tau_{0, t} N_{0, t}\right)dt 
%=  {r_{0,t}}\Pi_t d t  + \left(b_{0-, t} +G_t-\tau_{0, t} N_{0, t}\right)dt\, .  % \\
%%\end{aligned}
%%\right.	
%\end{equation}
%
%
%For $t\notin \mathcal T$ where $\{\Pi_t\}$ does not jump, $X_t$ is continuous because $U_{C,t}$ is also continuous. Differentiating $X_t$ both sides of  (\ref{X}) with respect to  $t$ and using (\ref{eqPi}) together with (\ref{foch1}) and (\ref{foch2}), we obtain: 
%\begin{eqnarray}\label{dX}  
%		d X_t  &=& U_{C, t}d\Pi_t + \Pi_t d\left(e^{\rho t}q_{0, t}U_{C, 0} \right)\notag\\
%		& = & U_{C, t} \left[ -\Pi_t\frac{dq_{0, t}}{q_{0, t}} + \left(b_{0-, t} +G_t-\tau_{0, t} N_{0, t}\right)dt\right]  + \Pi_t \left( q_{0, t}U_{C, 0}\rho e^{\rho t}dt +e^{\rho t} U_{C, 0}q_{0, t}\frac{dq_{0, t}}{q_{0, t}} \right) \notag\\
%		& = & -X_t \frac{dq_{0, t}}{q_{0, t}} + \left( b_{0-, t}U_{C, t}-C_{0, t}U_{C, t} -N_{0, t}U_{N, t} \right) dt + \rho X_t dt +X_t \frac{dq_{0, t}}{q_{0, t}} \notag\\ 
%%		& = &   \left[\rho X_{t}-C_{0, t}U_{C, t} -N_{0, t}U_{N, t}  +b_{0, t}U_{C, t}\right]dt \, . % \\
%%		%&  & +\, .
%%		\end{eqnarray*} 
%%\begin{eqnarray}
%%d X_t 
%& = &   \left(\rho X_{t}-C_{0, t}U_{C, t} -N_{0, t}U_{N, t}  +b_{0, t}U_{C, t}\right)dt \, .    % \\
%%&  & +\, .
%\end{eqnarray}
%
%\paragraph{What happens when a jump occurs.} 
%
%%the Absent jumps, i.e., for
%Where  $\{\Pi_t\}$ jumps at  $t\in \mathcal T$, we have 
%\begin{equation} \label{eqPi2}
%\Pi_{t}  =  \Pi_{t-} \frac{q_{0, t-}}{q_{0, t}}\, , \quad t \in \mathcal T \,.
%\end{equation}
%Thus, the time-0 value of the government's time $t$ cumulative liability $ q_{0, t} \Pi_{t} $ is the same before and after a jump. 
%% in that  
%
%
%%\NW{Since we assume that  $\{G_t\}$ and $\{b_{0,t}\}$ are {adapted processes},  optimal allocations  $\{C_{0, t}, N_{0, t}; t\geq 0\}$ may not be continuous. } 
%%Starting from  the    evolves  as
%
%
%%\begin{equation} \label{eqPi}
%%\left\{
%%\begin{aligned}
%%d\Pi_t  &=  -\Pi_t\frac{dq_{0, t}}{q_{0, t}} + \left(b_{0, t} +G_t-\tau_{0, t} N_{0, t}\right)dt\,, && t\in(t_i, t_{i+1}) \\
%% \Pi_{t_i} & =  \Pi_{t_i-} \frac{q_{0, t_i-}}{q_{0, t_i}}\,, && t=t_i \,.
%%\end{aligned}
%%\right.	
%%\end{equation}
%
%
%
%
%What happens where  $\{\Pi_t\}$ jumps at  $t\in \mathcal T$? Using (\ref{eqPi2}) and  first-order necessary condition (\ref{foch2}), we obtain $X_{t-} = \Pi_{t-}U_{C, t-} = \Pi_{t-}q_{0, t-}e^{\rho t-} U_{C, 0} =\Pi_{t}q_{0, t}e^{\rho t} U_{C, 0} = \Pi_{t}U_{C, t}=X_t$.  

We have established: %proven the next lemma. %we show that $X_t$ is still continuous in $t$. 
	\begin{proposition} 
	%Regardless whether a jump arrives or not at $t$, 
		{For a given competitive equilibrium, $X_t$ defined in (\ref{X}) is continuous in time and for all $t \geq 0$ evolves according to  (\ref{dX}).} 
	\end{proposition}
	%Because  welfare maximization calls for  continuity of the household's net worth $\Pi_t$ multiplied by the marginal utility $U_{C, t}$, 
%	\TS{Neng and Wei: the following sentence needs editing. To be valid, we'd have to point out that the evolution equation for $X_t$ has built in already that the Ramsey plan has already been computed -- say
%	by the Lagrangian method. This is probably built in with some of the equations pointed to. I'll have to figure this out. Let's talk.} \NW{Tom: this continuity of $X$ actually holds in any competitive equilibrium. We suggest deleting the following paragraph in red.} % below.} 
%
%%A key takeaway is that i
%	\WJ{%In a  competitive equilibrium, when state transitions from $s_{t-}$ to $s_{t}\neq s_{t-}$, the cumulative liability  jumps from $\Pi_{t-}$ to  $\Pi_{t}\neq \Pi_{t-}$. Importantly, the jump of $\Pi_{t}$ is exactly offset by the jump of $U_{C,t}$. 
%	That is, the Ramsey planner wants to  make $X_t$ continuous,  even where  the initial debt structure $\{b_{0-,t}\}$ or  the government spending process $\{G_t\}$ jumps. 
%	}
	
%	That is, the Ramsey planner wants to  make $X_t$ continuous,  even where  the initial debt structure $\{b_{0-,t}\}$ or  the government spending process $\{G_t\}$ jumps. 

%\WJ{	
%\begin{lemma}\label{lemmaPi}
%$\Pi_t$ defined in \eqref{Pi} evolves as follows:
%\begin{align}\label{dPi}
%d\Pi_t  
% =& r^{}_{0, t-}\Pi_{t-}dt 
% + \left(b_{0-, t-} + G_{t-} - \tau_{0, t-} N_{0, t-} \right)dt   \notag \\
% &\quad + \sum_{j\neq s_{t-}} \left(e^{-\eta(s_{t-}, j)} - 1 \right)\Pi_{t-} 
% \big(d {\mathcal J}_t^{(s_{t-}, j)} 
% - \lambda_{s_{t-}, j} e^{\eta(s_{t-}, j)} dt \big) \, ,
%\end{align}
%where $\eta(s_{t-}, s_t) = \ln\frac{U_{C, t}}{U_{C, t-}}$.  %\NW{Make it general: }
%\end{lemma}
%}
%


%
%\begin{lemma}\label{lemmaPi}
%$\Pi_t$ defined in \eqref{Pi} evolves as: %follows the dynamics
%\begin{align}\label{dPi}
%d\Pi_t  
% =&r^{}_{0, t-}\Pi_{t-}dt +\left(b_{0-, t-} +G_{t-}-\tau_{0, t-} N_{0, t-} \right)dt   \notag \\
% &\quad   +    \sum_{j\neq s_{t-}} \left(e^{-\eta(s_{t-}, j)}-1 \right)\Pi_{t-} \big(d {\mathcal J}_t^{(s_{t-}, j)} - \lambda_{s_{t-}, j}e^{\eta(s_{t-}, j)} d t \big) \, , 
%\end{align}
%where \NW{Make it general: } $\eta(i, j) = \ln\frac{U_C(C(j), N(j))}{U_C(C(i), N(i))} $.  
%\end{lemma}
%We make two observations for \eqref{}. First, note that $ \lambda_{s_{t-}, j}e^{\eta(s_{t-}, j)} $ in the last term is the `risk-adjusted' transition rate,\footnote{Technically speaking, $ \lambda_{s_{t-}, j}e^{\eta(s_{t-}, j)} $ is the regime-transition rate under the risk-neutral measure, see, e.g., \cite{}, for a textbook treatment.} which is different from the  rate $ \lambda_{s_{t-}, j}$ at which the state transitions from $s_{t-}$ to $j\neq s_{t-}$ under the physical measure. This part captures the `price' of  risk associated with the regime transition from $s_{t-}$ to $j$. Second, $\Pi_t$ jumps discretely from  $\Pi_{t-}$ to $\Pi_{t-} e^{-\eta(s_{t-}, j)}$ when the state transitions from $s_{t-}$ to $j\neq s_{t-}$. This 
%part corresponds to the `quantity' of  this regime-transition risk. 
%In summary, both `price' and `quantity' of risk arise from the household's demand  for bearing systematic regime-transition risk in equilibrium. 

%	\begin{proof}
%				First, for $t \in \mathcal T$, $X_{t_-}=X_{t}$. 
%Second, for $t\notin \mathcal T$, 
%	\end{proof}
%

%As the government's spending  processes can arbitrarily depend on time,   

%although  hough $b_{0,t}$ and   $G_t$ may be  discontinuous in $t$, the Ramsey planner chooses to make $X_t$ continuous in $t$. 

 
%Since   the initial debt structure $b_{0-,t}$ and  government spending $G_t$ are generally time dependent, the Ramsey problem is time inhomogeneous, so $t$ will appear in  Ramsey planner's value function. %Below we show  that  together allow us to Bellmanize the Ramsey problem. %optimally makes nsures the net worth of cumulative primary deficit is continuous in time such that it is not distorted. 



%In addition to calendar time $t$, the product  $X_t$ can serve as a component of 


%Nonetheless, we can  recursively formulate the  Ramsey plan using  $X_t$  as a state variable together with calendar time $t$. 
%are not homogeneous.  %not Markovian. %Let $X_t$ denote $B_t$  multiplied by the marginal utility $U_{C, t}$ of consumption: 


%\footnote{It can be verified that 
%\begin{align}\label{X}
%\int_0^t e^{-\rho s}  \left[ \big[ C_{0, s} - b_{0, s} \big] U_{C, s} + N_{0, s}U_{N, s} \right]  d s + e^{-\rho t}  X_t %= e^{-\rho t} B_t U_{C, t}\,.
%\end{align}
%is a constant. Its counterpart is a martingale in a stochastic version of our model to be studied in a sequel to this paper.}

%
%
%Definition (\ref{qprice}) of the time $0$ price  of a time $t$ consumption flow   and first-order necessary condition  (\ref{foch2}) imply that%The dynamics of $X_t$ satisfies

%
%
%	\begin{align}\label{dB}
% &d X_t  \nonumber  \\   = & \left[\rho (B_t U_C(C_{0, t}, N_{0, t}))-C_{0, t}U_C(C_{0, t}, N_{0, t})-N_{0, t}U_{N}(C_{0, t}, N_{0, t}) +b_{0, t}U_C(C_{0, t}, N_{0, t})\right]dt. 
%\end{align}
%%where $M_{0,t}= U_C(C_{0, t}, N_{0, t})$ is the marginal utility of consumption determined at time $0$: 
%
%\newpage 


%\NW{In the spirit of Ljungqvist and Sargent (2018, ch.XX), we formulate a Ramsey plan recursively by first taking $(X_t, t)$ as a state vector	that confronts a time $t$ continuation Ramsey planner.  Then we study  a time $0$ Ramsey planner who chooses $X_0$.  } 

%Next, we provide a procedure to obtain the time-$0$ optimal smooth equilibrium using dynamic programming method. According to the two steps characterization in section 2, we will first derive HJB equation for the optimization problem, taking $C_{0, 0}$ as given, and then choose $C_{0, 0}$ such that the smooth equilibrium condition holds and (\ref{obj}) is solved. With the newly introduced state variable $X_t$, we characterize the optimal smooth equilibrium by the following two steps:
%
%
%\begin{itemize}
%	\item The continuation Ramsey problem: Taking $C_{0, 0}$ as given, we have a fixed value of $X_0$. Then (\ref{dy}) fully describes the evolution of $X$ given policies. the planner chooses $\{C_{0, t}; t> 0\}$ for all time$-t$ ($t> 0$) onward government. Note that $X_t$ is the state variable. 
%	\item The time$-0$ Ramsey problem: The time$-0$ government chooses time$-0$ consumption $C_{0, 0}$ and the marginal value of $X_0$, subject to the smooth equilibrium condition, to optimize the time$-0$ objective.
%\end{itemize}

\subsection{Ramsey Problem as a Dynamic Program} \label{sec:DPdefine}



We  reformulate  the Ramsey problem recursively in   four steps. First, for a given $X_0$, we define and solve  a dynamic programming (DP) problem with $X_t$ and state $s_t$ being the state variables in subsection \ref{sec:DPdefine}. Then, in subsection \ref{sec:DPRamsey} we use the DP problem  to  characterize the Ramsey plan. %solve for $X_0$ and $C_0$ 



\begin{definition} [{\bf  DP Problem}] \label{definDP}
	Confronting  $X_0$ at $t=0$, 
	%\paragraph{Dynamic programming for the continuation planner}  at time $0$ confronts $X_0$ as a state variable and
	the Ramsey planner  solves 
	%. Then the continuation Ramsey planer chooses $\{C_{0, t}; t\geq \varepsilon\}$ to optimize
	\begin{equation}\label{continuation_obj}
		\max_{\{C_{0, t}; t\geq 0 \}}\; 
		\mathbb{E}_0 \,\bigg[\int_{0}^{\infty} e^{-\rho t} U({C}_{0, t}, {C}_{0, t}+G_t) d t  \bigg],
	\end{equation}
	where maximization is subject to the evolution equation  (\ref{dX}). % of   state variable $X_t$. 
\end{definition} 
%
%As we have effectively coded the history-relevant information via $X_t$ given in \eqref{dX}, we next use dynamic programming to formulate and solve this problem. %(see e.g., Fleming and Soner, 2006).
% casts in terms of  a family of optimization problems. 
%Starting from   states vector $(X_t, t)$ at $t\in[0, \infty)$, 
Let $V(X_0, s_0)$ denote the (optimal) value function for this problem. %To solve $V(X_0, s_0)$, we first 
%Construct the following time$-t$ optimization problem: %value function 
%\NW{\begin{equation}
%	e^{-\rho t}  {V}(X_t, t) = \max_{\{C_{0, s}; s\geq t \}}\; 
%	\int_t^{\infty} e^{-\rho s} U({C}_{0, s}, {C}_{0, s}+G_s) d s .
%\end{equation}}
Let
\begin{equation}\label{Vcontinue}
	%{V}(X_t, t) = 
V(X_t, s_t) =	\max_{\{C_{0, u}; u\geq t \}}\; 
	\mathbb{E}_t	\, \bigg[\int_t^{\infty} e^{-\rho(u-t)} U({C}_{0, u}, {C}_{0, u}+G_u) d u \bigg] \,  
\end{equation}
where maximization is subject to (\ref{dX}).
%Let $V(X_t, s_t)$ denote the value function where $(X_t, s_t)$ are the state variables. 
%
%Next, we construct a martingale that will help us to derive the process defined by
%\begin{equation*}
%	\int_{0}^t e^{-\rho t} U(C_{0, t}, C_{0, t}+G_t)dt + e^{\rho t}V(X_t, s_t) 
%\end{equation*}
%is a martingale under physical measure so its drift is zero
%\begin{equation} \label{martingaleV}
%	\mathbb{E}_t \left[d \left(e^{\rho t}V(X_t, s_t)  \right) \right] + e^{-\rho t} U(C_{0, t}, C_{0, t}+G_t)dt =0
%\end{equation}
%Simplifying \eqref{martingaleV} for all $t\geq 0$, 
{Appendix \ref{appendix:A}} verifies that $V(X_t, s_t)$ satisfies the following HJB equation: 
%\begin{eqnarray}\label{HJBV} \nonumber
%\rho V (X, s) & = &  \max_{N\in[G, 1], \beta} \quad U(N-G, 1-N) \\
%& & \quad  + \frac{\partial V}{\partial X} \left(\rho X+N U_{L}(N-G, 1-N)-C U_{C}(N-G, 1-N) \WJ{- \sum_{k\neq s}\lambda_{s k}\beta(s, k)}\right)  \nonumber \\
%&& \quad +\sum_{k\neq s}\lambda_{sk}\left[ \WJ{V(X+\beta(s, k), k)}-V(X, s) \right],
%\end{eqnarray}
%\begin{eqnarray}\label{HJBV} \nonumber
%\rho V (X, s) & = &  \max_{N\in[G, 1], \beta} \quad U(N-G, 1-N) \\
%& & \quad\quad \quad \quad+ \frac{\partial V}{\partial X} \left(\rho X+N U_{L}(N-G, 1-N)-C U_{C}(N-G, 1-N) \right)  \nonumber \\
%&& \quad \quad \WJ{- \frac{\partial V}{\partial X} \left(\sum_{k\neq s}\lambda_{s k}\beta(s, k) \right) }+\sum_{k\neq s}\lambda_{sk}\left[ \WJ{V(X+\beta(s, k), k)}-V(X, s) \right],
%\end{eqnarray}
\begin{align}\label{HJBV0} 
	\rho V = &   \max_{C_{0, t}} \;\; U(C_{0, t}, C_{0, t}+G_t) 
	+ \frac{\partial V }{\partial X_t} \left[\rho X_t- C_{0, t} U_{C, t}-(C_{0, t}+G_t) U_{N, t}+b_{0-, t} U_{C, t} \right]   \notag \\
	& \quad\quad \,\, +\sum_{k\neq s_t} \lambda_{s_t, k}\left[V(X, k)-V(X, s_t) \right]	 \, . 
	%	&& \quad   \;\;\; %\WJ{- \frac{\partial V}{\partial X} \left(\sum_{k\neq s}\lambda_{s k}\beta(s, k) \right) }
	%& \quad\quad \,\,	+\WJ{\sum_{k\neq s}\lambda_{s_tk}\left[ V(X+\beta{(s_t, k)}X, k)-V(X, s_t) - \frac{\partial V(X, s_t)}{\partial X} \beta{(s_t, k)}X   \right] }.
\end{align} 
%
%
Now guess and  verify that $V(X_t, s_t)$ takes the following  form:
\begin{eqnarray}\label{Vadditive0} 
	V(X_t, s_t)=-\Phi_0 X_t+H(s_t; \Phi_0)\, ,
\end{eqnarray}
where the scalar  $\Phi_0$ is to be chosen by the Ramsey planner. 
%The intuition for this result is as follows. The government optimally chooses the allocation so that the 
Substituting (\ref{Vadditive0}) into  (\ref{HJBV0}), we obtain %that $H(s)$, for $G=g(s) \left(s=1, \cdots, n\right)$, satisfies 
the following interconnected system of equations for $H(s_t; \Phi_0)$ for each $s_t=1, \cdots, n$: 
\begin{align}\label{HJBH00} %\nonumber
	\rho H(s_t; \Phi_0)  = &\;
	\max_{C_{0, t}} \, U(C_{0, t},C_{0, t}+G_t) 
	+ \Phi_0 \big[C_{0, t} U_{C, t}+(C_{0, t}+G_t) U_{N, t}-b_{0-, t} U_{C, t} \big]  \notag\\
	& \quad\,	+\sum_{k\neq s_t}\lambda_{s_t , \, k}\left[ {H(k; \Phi_0)}-H(s_t; \Phi_0) \right].
\end{align}

%If the government with full commitment chooses $\beta=0$, then the value function $V(X_t, s_t)$ satisfies
%\begin{align}\label{HJBV01} 
%	\rho V = &   \max_{C_{0, t}} \;\; U(C_{0, t}, C_{0, t}+G_t) 
%	+ \frac{\partial V }{\partial X_t} \left(\rho X_t- C_{0, t} U_{C, t}-(C_{0, t}+G_t) U_{N, t}+b_{0, t} U_{C, t} \right)   \notag \\
%	% & \quad\quad \,\, +\sum_{k\neq s} \lambda_{ks}\left[V(X, k)-V(X, s) \right]	 
%	%	&& \quad   \;\;\; %\WJ{- \frac{\partial V}{\partial X} \left(\sum_{k\neq s}\lambda_{s k}\beta(s, k) \right) }
%	& \quad\quad \,\,	+\WJ{\sum_{k\neq s}\lambda_{s_tk}\left[ V(X, k)-V(X, s_t) - \frac{\partial V(X, s_t)}{\partial X} \beta{(s_t, k)}X   \right] }.
%\end{align} 
%and the associated $H(s_t; \Phi_0)$ also satisfies \eqref{HJBH00}.


%Let $C(s_t; \Phi_0)$  denote the function for $C_{0, t}$ that maximizes \eqref{HJBH00}:
Define
\begin{equation}\label{optC00}
	C(s_t; \Phi_0) = \arg\max_{C_{0, t}} \, U(C_{0, t}, C_{0, t}+G_t) 
	+ \Phi_0 \left[ C_{0, t} U_{C, t}+(C_{0, t}+G_t) U_{N, t}-b_{0-, t} U_{C, t} \right]\, 
\end{equation}
and let %Let $N(s_t; \Phi_0)$ denote the corresponding optimal labor supply: 
\begin{eqnarray}  \label{optN00} % \nonumber
	%N^*(s) & = & \arg\max_{N\in[G(s), 1]} \left\{ U(N-G(s), 1-N) \right. \\
	%& &  \left. + \Phi_1 \left(\rho X+N U_{L}(N-G(s), 1-N)-C U_{C}(N-G(s), 1-N) \right)\right\}. \\
	N(s; \Phi_0) & = & C(s; \Phi_0)+G(s).  \label{optC0}
\end{eqnarray}
According to  (\ref{optC00})-\ref{optN00},  consumption $C_{0, t}$ and labor $N_{0, t}$  depend only on the contemporaneous exogenous state $s_t$ and  not  on $X_t$.

%The hedging policy $\beta_t$ is indeterminant. The most efficient hedging policy is the one that ensures the equilibrium $X_t$ to be continuous in time, $\beta_t=0$.
%Moreover, we will later show the optimal hedging ensures the equilibrium $X_t$ to be time invariant as long as the exogenous state $s_t$ does not switch. 



%We summarize the  planner's value function and the optimal policies for $t>0$ in the following proposition. For convenience, we introduce the following  vectors:
\begin{eqnarray*}
	\mathbb{V} &= & \left[V(X, 1; \Phi_0), \cdots, V(X, n; \Phi_0)  \right]^{\intercal} \\
	\mathbb{H}  & = &  \left[{H}(1; \Phi_0), \cdots, {H}(n; \Phi_0)\right]^{\intercal}  \,.
	%	\mathbb{A} & = &  ({A}(1; \Phi_0), \cdots, {A}(n; \Phi_0))^{\intercal} \\
	%	\mathbb{\widehat{A}} & = &   (\widehat{A}(1; \Phi_0), \cdots, \widehat{A}(n; \Phi_0))^{\intercal} \\
	%	\mathbb{\widetilde{A}} & = &   (\widetilde{A}(1; \Phi_0), \cdots, \widetilde{A}(n; \Phi_0))^{\intercal},
\end{eqnarray*}
%where $A(s; \Phi_0), \widehat{A}(s; \Phi_0), $ and $\widetilde{A}(s; \Phi_0) $ are given by %follows
%\begin{eqnarray*}
%	{A}(s; \Phi_0) & = & U(C(s; \Phi_0), N(s; \Phi_0)); \\
%	\widehat{A}(s; \Phi_0) & = & -N(s; \Phi_0) U_{N}-( C(s; \Phi_0) -b(s) ) U_{C}; \\
%	\widetilde{A}(s; \Phi_0) & = &U(C(s; \Phi_0), N(s; \Phi_0))+\Phi_0 \left[ -N(s; \Phi_0) U_{N}-( C(s; \Phi_0)-b(s) ) U_{C} \right] .
%\end{eqnarray*} 


\begin{proposition} \label{proH}
For  $s\in\{1, 2, \cdots, n\}$, let
	\begin{align*}
		{A}(s; \Phi_0)  =   U(C(s; \Phi_0),C(s; \Phi_0)+G(s)) 
	+ \Phi_0 \big[C(s; \Phi_0) U_{C}+(C(s; \Phi_0)+G(s)) U_{N}-b(s) U_{C} \big] 
	. \end{align*}  
Let \begin{eqnarray*}
	\mathbb{V} &= & \left[V(X, 1; \Phi_0), \cdots, V(X, n; \Phi_0)  \right]^{\intercal} \\
	\mathbb{H}  & = &  \left[{H}(1; \Phi_0), \cdots, {H}(n; \Phi_0)\right]^{\intercal}  \\
	\mathbb{A}  & =   & \left[{A}(1; \Phi_0), \cdots, {A}(n; \Phi_0)\right]^{\intercal}  \,.
	%	\mathbb{A} & = &  ({A}(1; \Phi_0), \cdots, {A}(n; \Phi_0))^{\intercal} \\
	%	\mathbb{\widehat{A}} & = &   (\widehat{A}(1; \Phi_0), \cdots, \widehat{A}(n; \Phi_0))^{\intercal} \\
	%	\mathbb{\widetilde{A}} & = &   (\widetilde{A}(1; \Phi_0), \cdots, \widetilde{A}(n; \Phi_0))^{\intercal},
\end{eqnarray*}
	Choices of $\{C_{0,t}\}$  for the  {\bf  DP Problem} indexed by  $\Phi_0\geq 0$ satisfy  (\ref{optC00})  and attain  value function \eqref{Vadditive0},
%	\NW{the following eqn repeated?} 
%	\begin{equation}
%	\NW{	V(X_t, s_t; \Phi_0) = -\Phi_0X_t+H(s_t; \Phi_0)\,.} 
%	\end{equation}
	where $H(s_t; \Phi_0)$ is given by
	\begin{equation}\label{HH}
		\mathbb{H} = \left(\rho \mathbb{I}-{\Lambda}  \right)^{-1}\mathbb{A} \, , 
	\end{equation}
	 $
	%	\mathbb{V} &= & \left(V(X, 1; \Phi_0), \cdots, V(X, n; \Phi_0)  \right)^{\intercal} \\
	%	\mathbb{H}  & = &  ({H}(1; \Phi_0), \cdots, {H}(n; \Phi_0))^{\intercal}  \\
%	\mathbb{A}  =   \left[{A}(1; \Phi_0), \cdots, {A}(n; \Phi_0)\right]^{\intercal} \, , 
	%	\mathbb{\widehat{A}} & = &   (\widehat{A}(1; \Phi_0), \cdots, \widehat{A}(n; \Phi_0))^{\intercal} \\
	%	\mathbb{\widetilde{A}} & = &   (\widetilde{A}(1; \Phi_0), \cdots, \widetilde{A}(n; \Phi_0))^{\intercal},
	$
%	and   for $s\in\{1, 2, \cdots, n\}$, 
%	\begin{align}
%		{A}(s; \Phi_0)  =   U(C(s; \Phi_0),C(s; \Phi_0)+G(s)) 
%	+ \Phi_0 \big[C(s; \Phi_0) U_{C}+(C(s; \Phi_0)+G(s)) U_{N}-b(s) U_{C} \big]
%	. \end{align}  
	
\end{proposition}





\subsection{Characterizing the Ramsey Plan} \label{sec:DPRamsey}

%The consumption stream  $\{C_{0,t}\}$ is a {adapted}  process: % so that,  %meaning that, at $t=0$,  %implies that %The smooth 
%$C_{0, 0} = \lim_{t\downarrow 0} C_{0, t} = C(s_0; \Phi_0)\, , $ which  implies that 
%\end{equation}
%$C_{0, 0}$ is a function of $\Phi_{0}$. %, which we write as $C_{0,0}(\Phi_0)$. 
%Subject to the implementability constraint (\ref{impletconstraint}), in addition to choosing $\{C_{0, t}, t >0\}$ by using \eqref{optC00},
Setting $X_0=0$ and  $\Pi_0=0$ and keeping   the initial term-debt  structure $\{b_{0-, t}=b(s_t); t\geq 0\}$ fixed,
the  expected utility attained by the representative agent under the Ramsey plan at state $s_0$ is  %chooses $\Phi_{0}$ by solving % to maximize the household's welfare by solving%, or equivalently consumption $C_{0, 0}$, by solving:
\begin{eqnarray}\label{W}
	W(s_0) %(\{b_{0, t}; t\geq 0\})
	\; = \; \max_{\Phi_{0}} \; V(X_{0}, s_0; \Phi_0)  \, ,
\end{eqnarray}
%where   $  V(X_{0}, s_0; \Phi_0) $ is given in (\ref{Vadditive0}). %the optimization is 
where the maximization on the right side is subject to implementability constraint (\ref{impletconstraint}). 
%\TS{Tom and Neng and Wei -- let's talk about how to improve the  preceding sentence. } \WJ{Tom: please see our edits in blue.} %and solves 
Define%  the following function $\Psi_{}(\Phi)$: Where 
\begin{eqnarray}\label{eq:Psi0}
	I_0(s_0; \Phi_{0}) = \mathbb{E}_0 \, \bigg[ \int_{0}^{\infty} e^{-\rho t} \left[ C(s_t; \Phi_{0})U_{C, t}+(C(s_t; \Phi_{0})+G(s_t))U_{N, t} -b(s_t) U_{C, t}  \right]  dt  \bigg]  \, .%\nonumber \\
	%&& - \int_{0}^{\infty} e^{-\rho t}b_{0, t} U_C dt%+B_{0}U_{C}( C(\Phi_{}, b_{0, 0}, G_0) ,C(\Phi_{}, b_{0, 0}, G_0)+G_0) 
\end{eqnarray}
implementability constraint \eqref{impletconstraint} requires that $\Phi_0$ satisfies $I_0(s_0; \Phi_{0}) =0$.
For a given initial state $s_0$, let $\mathcal A=\{\Phi_0\geq 0: I_0(s_0; \Phi_0)=0\}$ denote an admissible set for $\Phi_0$.
%\NW{Introduce the admissible set for $\Phi_0$?} 
Using (\ref{Vadditive0}) and $X_0=0$, the representative agent attains utility criterion %equivalent to
\begin{eqnarray}\label{WWW} 
	W(s_0)  & = & \max_{\Phi_0\in \mathcal A}\, V(0, s_0) \;\;
	= \;\; \max_{\Phi_0\in \mathcal A} \, H(s_0; \Phi_0)  \,.
\end{eqnarray}
Note that  $\Phi_0$ does not appear explicitly in \eqref{WWW}, since $X_0 =0$ because $\Pi_0 =0$; nevertheless  the first term in \eqref{Vadditive0}, % as $X_0=0$, 
$\Phi_0$  appears implicitly in \eqref{WWW} via the  $H(s_0; \Phi_0)$ function shown in \eqref{HH}.
Evidently,  $\Phi_0$ depends on $s_0$. %That is, which state the economy starts from matters for the Ramsey plan not
%Next, we summarize the recursive solution for the Ramsey plan. 
We are now equipped to state:



%Then we can show that an initial debt is feasible depends on whether there exists a parameter $\Phi_0$ for the time$-0$ government.  
%\begin{proposition}
%	A debt and spending scheme $\{b_{0, t}, B_{0}, G_t, t\geq 0\}$ is feasible if and only if there exists a $\Phi_{}$ that solves
%	\begin{equation}
%	\Psi_0(\Phi_{})=0.
%	\end{equation}
%\end{proposition}

%The optimal smooth equilibrium can be summarized in the following propositions.

%Next, we summarize the Ramsey planner's solution. 
\begin{proposition} \label{propbellman} 
	When it exists, a Ramsey plan satisfies
	\begin{eqnarray*}
		& C_{0, t}^{*}  =  C(s_t; \Phi_{0}^{*}); \quad 
		N_{0, t}^{*}  =  C_{0, t}^{*}+G(s_t) ;  \\
		& 	\tau_{0, t}^{*}   =   1+\frac{U_{N, t}(C_{0, t}^{*}, N_{0, t}^{*})}{U_{C, t}(C_{0, t}^{*}, N_{0, t}^{*})}; \quad
		q^*_{0,t}  = 
		e^{-\rho t}\frac{U_{C, t}(C_{0, t}^{*}, N_{0, t}^{*})}{U_{C, 0}(C_{0, 0}^{*}, N_{0, 0}^{*})}
	\end{eqnarray*}
	where $\Phi_{0}^*= \arg\max_{\Phi_0\in \mathcal A} \; V(0, s_0).$ Under a Ramsey plan, % $X_t = X^*_t$ where %and  satisfies
	%	\begin{equation}
	%		X_t^* = \mathbb{E}_t \int_{t}^{\infty}  e^{-\rho (u-t)}\left[C_{0, u}^*U_{C, u}(C_{0, u}^*, N_{0, u}^*)+N_{0, u}^*U_{N, u}(C_{0, u}^*, N_{0, u}^*) -b(s_u)U_{C, u}(C_{0, u}^*, N_{0, u}^*)\right]du	
	%	\end{equation}
	%	Under the Ramsey plan, 
	the representative agent's value  is 
	\begin{equation}
		W(s_0) = H(s_0; \Phi_0^*).
	\end{equation}
\end{proposition}
%Next, we show that  a %The following proposition states that the
% Ramsey plan  obtained via  the section \ref{sec:environ} Lagrangian method equals the plan described in Proposition \ref{propbellman}. %using the Bellmanization method in this section. 
%s generate the same Ramsey plan and associated allocation, tax rate process, and price system.
% and that the optimal policy obtained from the dynamic programming method that based on a pure short-term debt rolling strategy is equivalent to the optimal smooth equilibrium obtained from the Lagrangian method. 
%\begin{proposition}
%	A %The following proposition states that the
%	Ramsey plan computed with  the Proposition   \ref{sec:environ} Lagrangian method    equals the Ramsey plan  computed with  the Proposition \ref{propbellman} Bellman equation. % including $ \Phi_{0}^* = \Lambda_0^{*}$. } 
%%	The solution summarized in Proposition \ref{propbellman} is the same as the solution summarized in Proposition \ref{proplag} (via the Lagrangian method).
%%	\begin{eqnarray*}
%	%		& W(\{b_{0, t}; t\geq 0\}) = \mathcal{L}_0^*; \,  \Phi_{0}^* = \Lambda_0^{*}\,, \\
%	%		& C^*_{0, t} = \widetilde{C}_{0, t}; N^*_{0, t} = \widetilde{N}_{0, t}; \tau^*_{0, t} = \widetilde{\tau}_{0, t}; q^*_{0, t} = \widetilde{q}_{0, t}\,.
%	%	\end{eqnarray*}	
%	\end{proposition}


%\TS{The next section studies Ramsey plans for settings in which  $ \{b_{0,t}, G_t;t\geq 0\}$ are continuous in $t$.} %Second, it prepares us for } 


%\WJ{Then with continuous settings, $B$ is automatic continuous under Ramsey plan but then with minimal commitment in that $C$ is continuous, it is automatic that Ramsey plan is confirmed by future governments. That is, future government simply cannot immediately renege on its debt obligation in that it has to honor the purchasing power of $B$.} 





\section{Financing Deficits}\label{sec:financing}

%In this section, we discuss how  the government finances its cumulative liabilities, $\Pi^\ast$ given in \eqref{Pi}, under the Ramsey plan characterized in the preceding section. 
We alter \citeauthor{LucasStokey1983}'s formulation by 
%expanding the class of securities that the government can trade.  Technically,  we expand  \citeauthor{LucasStokey1983}'s  ''contractible subspace'' by 
allowing the government
to trade  ''instantaneous debt'', a security typically  included  in complete markets models in finance,  e.g., \cite{BlackScholes1973}, \cite{merton1973theory}, and \cite{HarrisonandKreps1979}. In this way,
we extend  \citeauthor{LucasStokey1983}'s assumption of complete markets to our continuous time setting.
%
%In addition to the term debts in the original \citeauthor{LucasStokey1983}'s formulation, we introduce instantaneous debt to the {contractible subspace}. 
%This debt plays a crucial role in our financing arrangement and is the same type of short-term debt widely used in asset pricing, e.g., \cite{BlackScholes1973}, \cite{merton1973theory}, and \cite{HarrisonandKreps1979}. 


\subsection{Liability Management}\label{subsec:financing1}
%Let the current state be $s_t$. 
We seek a  dynamic portfolio management strategy that finances the government's cumulative liabilities, $\Pi_t^\ast$, under the Ramsey plan. 
Evaluating	 $\Pi_t$ given in \eqref{dPi0} under the Ramsey plan yields: 
	\begin{align} \label{dPiast}
	d\Pi_t  ^\ast - 
	 r^{\ast}_{0, t-}\Pi_{t-}^\ast  dt 
&= \left(b_{0-, t-} + G_{t-} - \tau_{0, t-} N_{0, t-} ^\ast  \right)dt   \notag \\
	&\quad + \sum_{j\neq s_{t-}} \left(e^{-\eta^\ast (s_{t-}, j)} - 1 \right)\Pi_{t-} ^\ast 
	\big(d {\mathcal J}_t^{(s_{t-}, j)} 
	- \lambda_{s_{t-}, j} e^{\eta^\ast (s_{t-}, j)} dt \big) \, . 
	\end{align}
We require 
 a set of financial claims that allow us to finance $d\Pi_t  ^\ast - 
	 r^{\ast}_{0, t-}\Pi_{t-}^\ast  dt $, the net change of $\Pi_t  ^\ast$ over a small $dt$ interval. The government can issue  instantaneous debt that is locally risk free in the sense that   over an infinitesimal time interval $dt$ it returns %$(t, t+dt)$ are  
$r^\ast_{0, t}dt$.  %which is known at the beginning of 
%That is, whether the state transitions or not over this $dt$ interval, the interest rate for this instantaneous debt  is the same. 
Such  instantaneous debt is  not state-contingent.  
The government also has access to competitively priced hedging contracts that allow it to manage  state-transition risk exposure. 
Because $s_t$ resides in one of  $n$ states, starting from $s_{t-}$, we have $n-1$ distinct state transitions. For each possible transition from $s_{t-}$ to $j\neq s_{t-}$,  there is a hedging contract. Over a small time interval  $dt$, the holder of a unit of such a  hedging contract receives a  unit  payoff at $t$ if and only if the state transitions from $s_{t-}$ to another state $j\neq s_{t-}$ at time $t$, i.e., $d {\mathcal J}_t^{(s_{t-}, j)}=1$. The holder pays an actuarially fair hedging premium, $\lambda_{s_{t-}, j}e^{\eta^*(s_{t-}, j)}$,  per unit of time before a transition occurs. The premium equals the transition rate $\lambda_{s_{t-}, j}$ multiplied by a factor $e^{\eta^*(s_{t-}, j)}$ that compensates for  state-transition risk.  If the transition to state $j$ increases the marginal utility $U_C$, $\eta^*(s_{t-}, j)>0$, then investors receive  a positive hedging premium: note that $\lambda_{s_{t-}, j}e^{\eta^*(s_{t-}, j)}$ exceeds  state transition probability $\lambda_{s_{t-}, j}$. 

%a hedging contract that pays off one dollar if and only if $d {\mathcal J}^{}=1$. The hedging premium for this contract is 

Together with the state non-contingent instantaneous  debt, this  collection of  $(n-1)$ distinct hedging contracts  spans all possible contingencies over a small time interval $dt$, so  markets  are  dynamically complete %\footnote{\WJ{Finally, the government can also issue state-contingent term debt as markets are complete. TO PLACE LATER.}}
 in the tradition of  \cite{BlackScholes1973}, \cite{merton1973theory}, and \cite{HarrisonandKreps1979}.  
%%as markets are complete and competitive,
%
%%a stochastic transition from the current
%
%%consider a liability and risk management strategy where the 
% %Interest payments  %where $r^\ast_{0, t}$ is the instantaneous interest rate given in (\ref{interestrate}).  %we  obtain the   issuance of instantaneous debt as follows: With no new issuance and no  principal repayment, the change of the instantaneous debt balance is $dB_t = 
%%  The first term of the right hand side, 
%%Below we describe how a government finances its cumulative deficits $\Pi_t^\ast$. %debt dynamics %using instantaneous debt.
%%\paragraph{Financing cumulative deficits $\Pi_t$.} 
%%When $t\notin \mathcal T$, 
%%Under the Ramsey plan, 
%
%%\[ \Delta S  + B \] 
%
Let $B^{(0)}_t$  denote the  balance of the state non-contingent   instantaneous debt  at $t$. Let $\Delta_{t-}^ {(s_{t-}, j)}$ denote the hedging demand for state transition from $s_{t-}$ to $j\neq s_{t-}$. Let 
$\partial_t b_{t-, u}$  denote  the (incremental) term-debt coupon   issued at time $t-$ and due  at $u$. %The third term {$\left(\int_{0}^{t-} \partial_u b_{u, t-} du\right)dt $} is the  (incremental) term debt due at $t-$,     issued cumulatively from 0 to $t$. 
%The government's debt portfolio at $t$ includes state-contingent term debts and this newly introduced instantaneous debt. %and a set of hedging contracts
%
%
%
%
%\begin{eqnarray} \label{finB0}
%d B_t  %=  -\Pi_t\frac{dq_{0, t}}{q_{0, t}} + \left(b_{0, t} +G_t-\tau_{0, t} N_{0, t}\right)dt 
%& = &   {r_{0,t-}^*} B_{t-} d t  + \big(b_{0, t-} +G_{t-}-\tau_{0, t-} N_{0, t-}\big)dt \notag \\
%&  & \quad\quad + \sum_{j\neq s_{t-}} B_{t-}\big(e^{-\eta^*(s_{t-}, j)}-1\big) \left(d {\mathcal J}_t^{(s_{t-}, j)} -  \lambda_{s_{t-}, j}e^{\eta^*(s_{t-}, j)}d t \right)    \, . 
%\end{eqnarray}
%
%
%
%\newpage 
%

%Next we equate the government's uses and sources of funds by forming a portfolio. Specifically, the following dynamic law of motion holds: 
The evolution of sources and uses of government funds satisfies % the following evolution equation:
%\footnote{%\WJ{Interpretation for the jumps.} 
%		When $t \in \mathcal T$, all instantaneous debt and term debt must be re-evaluated under the Ramsey plan from time $t-$ to time $t$ in order to preserve the purchasing power of  $\Pi_t^\ast$ under the Ramsey plan. That is, the instantaneous debt balance shall change to $B_{t}=B_{t-}{q^\ast_{0, t-}}/{q^\ast_{0, t}}$ from $B_{t-}$ and the term debt  service flows change to $b_{t, s} = b_{t-, s}{q^\ast_{0, t-}}/{q^\ast_{0, t}}$ from $b_{t-, s}$ for all $s\geq t$. %Doing so ensures that $\Pi^\ast_{t}= \Pi^\ast_{t-}\frac{q^\ast_{0, t-}}{q^\ast_{0, t}}$,  
\begin{align} \label{dynamicbudgetconstraint5}
d\Pi_t ^\ast -r_{0, t-}^\ast \Pi_{t-}^\ast dt  = &  {\underbrace{dB^{(0)}_t  -r_{0, t-}^\ast B^{(0)}_{t-} dt}_{\text{net non-contingent instantaneous debt issuance}}}   +  {\underbrace{ \left( \mathbb{E}_{t-}\int_{u> t} q^\ast_{t-, u}\partial_tb_{t-, u} du \right) dt}_{\text{gross state-contingent term-debt issuance}}}  \notag \\
 & \quad -{\underbrace{ \left(\int_{u<t} \partial_u b_{u, t-} du\right)dt}_{\text{incremental term debt due}}}  
\; - {\underbrace{\sum_{j\neq s_{t-}}\Delta_{t-}^{(s_{t-}, j)} \left(d {\mathcal J}_t^{(s_{t-}, j)}-\lambda_{s_{t-}, j}e^{\eta^*(s_{t-}, j)}d t \right)}_{\text{stochastic hedging payoff netting premium payments}} } \,, %\notag \\
%  & \quad + {\underbrace{\sum_{j\neq s_{t-}} \left(e^{-\eta^*(s_{t-}, j)}-1 \right) \Delta_{}^{(s_{t-}, j)} d {\mathcal J}_t^{(s_{t-}, j)}}_{\text{payoff from hedging contracts at $t$}} } \,, %\notag \\
 %- {\underbrace{\sum_{j\neq s_{t-}} \left(e^{-\eta^*(s_{t-}, j)}-1 \right)\Pi_{t-}^* \left(\lambda_{s_{t-}, j}e^{\eta^*(s_{t-}, j)}d t \right)}_{\text{hedging premium payment rate}} } \, , 
\end{align}
where $\Pi_t^*$ defined in (\ref{Pi}) is the cumulative liability  under the Ramsey plan absent  term debt adjustments. 
The first term on the right side of \eqref{dynamicbudgetconstraint5} describes the net instantaneous non-contingent debt issuances. 
The second term is  gross state-contingent term-debt issuance over  $dt$. %where $\partial_t b_{t-, u}$  denotes a rate of increase the (incremental) term-debt coupon   issued at time $t-$ and due  at $u$. 
The third term is the sum of maturing term debts that had been    issued cumulatively from 0 to $t$.  %netting repayment to the due debt that issued in the past . %so that the total term debt due at $t$ is $b_{t,t} dt = b_{0,t} dt +\left(\int_{0}^t \partial_u b_{u, t} du\right) dt$.} 
%The left  side is the change of $\Pi_t$ over this period.  d\Pi_t ^\ast
The last  term equals the sum of contingent payoffs from all its $(n-1)$ hedging contracts, $\sum_{j\neq s_{t-}} \Delta_{t-}^{(s_{t-}, j)} d {\mathcal J}_t^{(s_{t-}, j)}$, minus the sum of hedging premium payments for all these contracts:  $\sum_{j\neq s_{t-}} \Delta_{t-}^{(s_{t-}, j)} \lambda_{s_{t-}, j}e^{\eta^*(s_{t-}, j)}d t . $ As the government's cumulative liability $\Pi_t^\ast$ increases with the hedging premium payments and decreases by $\Delta_{t-}^{(s_{t-}, j)}$ when a state transition from $s_{t-}$ to $j\neq s_{t-}$ occurs, i.e., when $d {\mathcal J}_t^{(s_{t-}, j)}=1$, there is a minus sign in front of this underbraced term.  
Note that over a $dt$ interval, either zero  or one state transition occurs. 

%hedging contract netting  hedging premium payments for 

%state non-contingent instantaneous and state-contingent term debts together with $(n-1)$ distinct hedging contracts to finance its primary deficits, so  
Equating the right sides of \eqref{dPiast} and \eqref{dynamicbudgetconstraint5} yields: 
\begin{align}\label{eq:dPi_star3}
%d\Pi_t^* -r_{0, t}^*\Pi_t^*dt = 
& \left(b_{0-, t-} + G_{t-} - \tau_{0, t-} N_{0, t-} ^\ast  \right)dt   
 + \sum_{j\neq s_{t-}} \left(e^{-\eta^\ast (s_{t-}, j)} - 1 \right)\Pi_{t-} ^\ast 
	\big(d {\mathcal J}_t^{(s_{t-}, j)} 
	- \lambda_{s_{t-}, j} e^{\eta^\ast (s_{t-}, j)} dt \big)  \notag \\
&	= \quad 
d B^{(0)}_t  -r_{0, t-}^\ast B^{(0)}_{t-} dt   +   \left( \mathbb{E}_{t-}\int_{u> t} q^\ast_{t-, u}\partial_tb_{t-, u} du \right) dt - \left(\int_{u<t} \partial_u b_{u, t-} du\right)dt  \notag \\
&  \quad\quad  
- \sum_{j\neq s_{t-}}\Delta_{t-}^{(s_{t-}, j)} \left(d {\mathcal J}_t^{(s_{t-}, j)}-\lambda_{s_{t-}, j}e^{\eta^*(s_{t-}, j)}d t \right)
	\end{align}
Since \eqref{eq:dPi_star3}  holds at all $t$, matching the term for the   transition from the current state $s_{t-}$ to $j\neq s_{t-}$ yields the following hedging demand $\Delta_{t-}^{(s_{t-}, j)}$: 
\begin{equation} \label{Delta} 
	\Delta_{t-}^{(s_{t-}, j)} =   \left( 1 - e^{-\eta^*(s_{t-}, j)} \right) \Pi_{t-}^\ast\, . 
\end{equation} 
Evidently, the government uses the   Delta hedging presented by \cite{BlackScholes1973} and \cite{merton1973theory}, appropriately adapted to our state-transition setting. 
Note that the  hedging demand associated with a  transition from $s_{t-}$ to $j\neq s_{t-}$ is unique. %has to be unique and is given by 
%For each state transition from $s_{t-}$ to $j\neq s_{t-}$, to finance the continuation of $\Pi_{t}^\ast$, we must have 
%For \eqref{eq:dPi_star3} to hold at all $t$, the hedging demand $\Delta_{}^{(s_{t-}, j)}$ has to be unique and is given by 
%Relation \eqref{eq:dPi_star2} equates uses and sources  of government funds. 
Cancelling the last terms on both sides of \eqref{eq:dPi_star3}, we obtain the following accounting identity:  
%hedging demand  given in \eqref{Delta}, the following has to hold: 
\begin{align}\label{eq:dPi_star22}
%d\Pi_t^* -r_{0, t}^*\Pi_t^*dt = 
& \underbrace{(G_{t-} - \tau^\ast_{0,t-} N^\ast_{0,t-} ) d t}_{\text{primary deficits}} + \, %
\underbrace{\biggl( b_{0-,t} + \int_{u<t}\partial_u\, b_{u,t-}\, du\biggr)  \, dt}_{\text{all maturing term debt repayments}}
\notag \\
 &\quad \quad  \quad \quad =   \underbrace{dB^{(0)}_t  -r_{0, t-}^* B^{(0)}_{t-}dt}_{\text{net instantaneous-debt issuance}}
 + \quad
 \underbrace{\biggl(\mathbb E_{t-}\int_{u > t} q_{t-,u}^*\, \partial_t\, b_{t-,u}\, du\biggr) dt}_{\text{gross term-debt issuance}} 
\, . 
%_{\text{maturing term debt repayment}} .
\end{align}
%netting  the associated interest costs $r \Pi_t^* d t$.} \WJ{Equation \eqref{eq:Pi_def} implies that the left side of~\eqref{eq:dPi_star} also 
%Eequals  $(b_{0-,t} + G_t - \tau^\ast_{0,t} N^\ast_{0,t} ) d t.$ Substituting this into \eqref{eq:dPi_star} and re-arranging  yields..... 


%On the right side of \eqref{dynamicbudgetconstraint5}, the first term describes the net instantaneous non-contingent debt issuances. 
%The second term is the net financing from new  debt issuance over  $dt$, where $\partial_t b_{t-, u}$  denotes a rate of increase the (incremental) term-debt coupon   issued at time $t-$ and due  at $u$. The third term {$\left(\int_{0}^{t-} \partial_u b_{u, t-} du\right)dt $} is the  (incremental) term debt due at $t-$,     issued cumulatively from 0 to $t$.  %netting repayment to the due debt that issued in the past . %so that the total term debt due at $t$ is $b_{t,t} dt = b_{0,t} dt +\left(\int_{0}^t \partial_u b_{u, t} du\right) dt$.} 
%%The left  side is the change of $\Pi_t$ over this period.  d\Pi_t ^\ast
%The fourth term is the sum of  premium payments for all $(n-1)$ hedging contracts. 
%The fifth term captures the effect of hedging contract payments on the change of $\Pi_t^\ast$ when a state transition occurs. 
%\NW{Need editing.} 
%
%Using the dynamics of $\Pi_t^\ast$ given in \eqref{dPi}, we obtain the state non-contingent instantaneous debt follows:
%\begin{equation}
%dB_t =r_{0, t-}^\ast B_{t-} dt +()
%\end{equation}
%
%
%We can also equivalently use state-contingent debts to finance the government's cumulative liabilities. 
%
%\NW{Use $N$ state contingent debt .....} 
%
%
%\begin{eqnarray} \label{dynamicbudgetconstraint5}
%d\Pi_t ^\ast -r_{0, t-}^\ast \Pi_{t-}^\ast dt  &= &  {\underbrace{dB_t -r_{0, t-}^\ast B_{t-} dt}_{\text{instantaneous non-contingent debt issuances}}}   +  {\underbrace{ \left( \int_{u> t} q^\ast_{t-, u}\partial_tb_{t-, u} du \right) dt}_{\text{term debt issued at $t-$}}}  -{\underbrace{ \left(\int_{u<t} \partial_u b_{u, t-} du\right)dt}_{\text{term debt due at $t$}}} \notag \\
%& & \quad\quad  + {\underbrace{\sum_{j\neq s_{t-}} \left(e^{-\eta^*(s_{t-}, j)}-1 \right)\Pi_{t-}^* d {\mathcal J}_t^{(s_{t-}, j)}}_{\text{payoff from hedging contracts at $t$}} }  \notag \\
%& & \quad\quad  - {\underbrace{\sum_{j\neq s_{t-}} \left(e^{-\eta^*(s_{t-}, j)}-1 \right)\Pi_{t-}^* \left(\lambda_{s_{t-}, j}e^{\eta^*(s_{t-}, j)}d t \right)}_{\text{hedging premium payment rate}} } \, , 
%\end{eqnarray}
%
%Next, we propose a {\it self-financing} dynamic portfolio strategy that  actively adjusts the instantaneous state non-contingent debt balance $B^{(0)}_t $ and  the $(n-1)$ state-contingent hedging demand functions, $	\Delta_{t-}^{(s_{t-}, j)}$ for $j\neq s_{t-}$,  to finance  the government's cumulative liability $\{\Pi_t^*\}$ process. 

%\subsection{Managing $\Pi_t^*$ without adjusting term debts}
\subsection{No rescheduling of initial term debts}\label{subsec:financing2}


%We propose the use of instantaneous debt. At time $t$, the cumulative liabilities under the Ramsey plan is $\Pi_t^*$.  
We describe a  dynamic portfolio strategy that finances  the government's cumulative liability $\{\Pi_t^*\}$ process by automatically adjusting the instantaneous state non-contingent debt balance $B^{(0)}_t $ and  the $(n-1)$ state-contingent hedging demand functions $	\Delta_{t-}^{(s_{t-}, j)}$ for $j\neq s_{t-}$ while leaving all exiting term debts untouched and issuing no new term debts.
Under this strategy,  $\partial_t\, b_{t, u}=0$ for all $(t\geq 0, u> t)$ and $t\geq 0$.  Simplifying \eqref{eq:dPi_star22}, we obtain the following process for the balance of the instantaneous state non-contingent debt $B_t^{(0)}$: %(excluding the hedging premium payments): 
%The government finances $\Pi_t^*$ using by setting $B_t= \Pi_t^*$ and \NW{XXX} so that 
\begin{eqnarray} \label{finB}
d B^{(0)}_t  %=  -\Pi_t\frac{dq_{0, t}}{q_{0, t}} + \left(b_{0, t} +G_t-\tau_{0, t} N_{0, t}\right)dt 
& = &   {r_{0,t-}^*} B_{t-}^{(0)} d t  + \big(b_{0, t-} +G_{t-}-\tau_{0, t-} ^\ast N_{0, t-}^\ast\big)dt   \, , \;\; \text{with}\; \; B_{0}=\Pi_0 = 0\, .  %\notag \\
%&  & \quad\quad + \sum_{j\neq s_{t-}} B_{t-}\big(e^{-\eta^*(s_{t-}, j)}-1\big) \left(d {\mathcal J}_t^{(s_{t-}, j)} -  \lambda_{s_{t-}, j}e^{\eta^*(s_{t-}, j)}d t \right)    \, . 
\end{eqnarray} % with zero balance, which increases when the government incurs a hedging loss.
A  hedging account tracks the cumulative losses due to hedging.  Let $B^{(1)}_t $ denote the balance of this account starting with $B^{(1)}_{0}=0$. If $B^{(1)}_t>0$, the government has a cumulative loss at $t$. If $B^{(1)}_t<0$, the government has a cumulative hedging gain at $t$. 
As  this account tracks the cumulative hedging  losses under the optimal hedging policy,  the government  can roll over the entire (loss) balance of this hedging account at the locally risk-free rate $r_{0,t-}^\ast$. This policy implies the following account balance dynamics:  
\begin{eqnarray} \label{finH}
d B^{(1)}_t  %=  -\Pi_t\frac{dq_{0, t}}{q_{0, t}} + \left(b_{0, t} +G_t-\tau_{0, t} N_{0, t}\right)dt 
& = &   {r_{0,t-}^*} B^{(1)}_{t-}d t \,  -  {\sum_{j\neq s_{t-}}
 %\left(e^{-\eta^*(s_{t-}, j)}-1 \right)\Pi_{t-}^*
 	\Delta_{t-}^{(s_{t-}, j)}
	 \left(d {\mathcal J}_t^{(s_{t-}, j)} - \lambda_{s_{t-}, j}e^{\eta^*(s_{t-}, j)}d t \right)}  \, . %  \text{with} \; H_{0-}=0\, .  %\notag \\
%&  & \quad\quad + \sum_{j\neq s_{t-}} B_{t-}\big(e^{-\eta^*(s_{t-}, j)}-1\big) \left(d {\mathcal J}_t^{(s_{t-}, j)} -  \lambda_{s_{t-}, j}e^{\eta^*(s_{t-}, j)}d t \right)    \, . 
\end{eqnarray}
where $\Delta_{t-}^{(s_{t-}, j)}$ is given in  \eqref{Delta}. 

Let $B_t = B^{(0)}_t + B^{(1)}_t$ denote the  combined balances in  the two accounts. Forward induction from $t=0$ path by path implies $B_t= \Pi_t^\ast$ for all $t$. Thus, combining \eqref{finB} with \eqref{finH} and using  \eqref{Delta}, we obtain 
 \begin{eqnarray} \label{finH}
d B_t  %=  -\Pi_t\frac{dq_{0, t}}{q_{0, t}} + \left(b_{0, t} +G_t-\tau_{0, t} N_{0, t}\right)dt 
& = &   {r_{0,t-}^*} B_{t-} d t  + \big(b_{0, t-} +G_{t-}-\tau_{0, t-} ^\ast N_{0, t-}^\ast\big)dt  \notag \\
&&  + {\sum_{j\neq s_{t-}} \big(e^{-\eta^*(s_{t-}, j)}-1 \big)B_{t-}^{} \left(d {\mathcal J}_t^{(s_{t-}, j)} - \lambda_{s_{t-}, j}e^{\eta^*(s_{t-}, j)}d t \right)}  \, . %  \text{with} \; H_{0-}=0\, .  %\notag \\
%&  & \quad\quad + \sum_{j\neq s_{t-}} B_{t-}\big(e^{-\eta^*(s_{t-}, j)}-1\big) \left(d {\mathcal J}_t^{(s_{t-}, j)} -  \lambda_{s_{t-}, j}e^{\eta^*(s_{t-}, j)}d t \right)    \, . 
\end{eqnarray}
Note that $B_t $ incorporates the cumulative effects of the government's primary deficits and   hedging losses (gains). Consequently,  $B_t$ perfectly tracks $\Pi_t^\ast$. 

Thus, our  dynamic portfolio strategy for financing $\Pi_t^\ast$ over time is composed of the instantaneous state non-contingent debt $B^{(0)}_t $ given in \eqref{finB} and  the $(n-1)$ state-contingent hedging demand functions, $	\Delta_{}^{(s_{t-}, j)}$ for $j\neq s_{t-}$ given in \eqref{Delta}. 
%This financing strategy ensures that the cumulative liabilities will be fully financed with instantaneous debt and purchasing hedging contracts.
%
%\WJ{
%%We make two observations regarding \eqref{dPi0}. 
%%First,  $ \lambda_{s_{t-}, j} e^{\eta(s_{t-}, j)} $ in the last term is the `risk-adjusted' transition rate,\footnote{Technically, $ \lambda_{s_{t-}, j} e^{\eta(s_{t-}, j)} $ is the regime-transition rate under the risk-neutral measure; see, e.g., \cite{}, for a textbook treatment.} which differs from the rate $ \lambda_{s_{t-}, j} $ at which the state transitions from $s_{t-}$ to $j\neq s_{t-}$ under the physical measure. Th ratio between the two transition rates, $ e^{\eta(s_{t-}, j)} $,  captures the `price' of the transition risk   from state $s_{t-}$ to $j$. The households demands a positive risk premium if the state transition increases the marginal utility:  $\eta(s_{t-}, j)>0$. %carries a positive price of risk.  
%%Second, as the state transitions from $s_{t-}$ to $j\neq s_{t-}$, i.e., when  $d {\mathcal J}_t^{(s_{t-}, j)} =1$, the cumulative liability jumps  from 
%$\Pi_{t-}$ to $\Pi_{t-} e^{-\eta(s_{t-}, j)} =  \Pi_{t-} \frac{U_{C, t-}}{U_{C, t}}= \Pi_{t-} \frac{q_{0, t-}}{q_{0, t}} = \Pi_t .$ That is, while $  \Pi_t$ and $q_{0, t}$ can both jump, the product, $q_{0, t}  \Pi_t$, is  continuous at all $t$. %
%%This part corresponds to the `quantity' of this regime-transition risk. 
%%
%%In summary, both the `price' and the `quantity' of risk arise from the household's demand for bearing systematic regime-transition risk in equilibrium.
%}
%
%\NW{
%%\[ XX B_t^L =  \Pi_t  - \Pi_{t-}   \] 
%\[ B_t^L =  \frac{q_{0, t-}}{q_{0, t}}  \big( \Pi_t  - \Pi_{t-} \big) 
%\] 
%} 
%\WJ{
%\[ B_t^L  \big( d J_t  - \big) 
%\]
%} 
%
%\NW{Imagine two states 
%\begin{eqnarray} \label{finB}
%d B_t^H %=  -\Pi_t\frac{dq_{0, t}}{q_{0, t}} + \left(b_{0, t} +G_t-\tau_{0, t} N_{0, t}\right)dt 
%& = &   {r_{0,t-}^*} B_{t-}^H d t   \, . 
%\end{eqnarray}
%\begin{eqnarray} \label{finB}
%d B_t^L %=  -\Pi_t\frac{dq_{0, t}}{q_{0, t}} + \left(b_{0, t} +G_t-\tau_{0, t} N_{0, t}\right)dt 
%& = &B_t^L  \left( d J_t  + {r_{0,t-}^*}  d t  - \lambda e^{\eta } d t   \right)\, . 
%\end{eqnarray}
%\[ B_t^L = \Pi_{LH} \]  
%}

By contrast, the debt-management strategy of \cite{LucasStokey1983} uses  active management of state-contingent term debts over time. 
\begin{remark}
	The debt rescheduling policies of  \cite{LucasStokey1983} use only term debt. A coupon flow management policy of $\{\partial_t b_{t-, u}; u\geq t,  t\geq 0 \}$ and $\{B_t^{(0)}=0; t\geq 0\}$  is used  to satisfy the dynamic budget constraint: 
	\begin{eqnarray} \label{dynamicbudgetconstraint-t}
	d\Pi ^\ast_t -r^*_{0, t-}\Pi_{t-}^\ast dt& = &  \left( \mathbb E_{t-} \int_{u> t} q^\ast_{t-, u}\partial_tb_{t-, u} du \right) dt  -\left(\int_{u<t} \partial_u b_{u, t-} du\right)dt  \notag \\
%	\left(\int_{u\geq t} q^\ast_{t, u}\partial_tb_{t, u} du \right) dt %}_{\text{term debt issued at $t$}} }
	%-%\NW{ \underbrace{ 
%	\left(\int_{0}^t \partial_v b_{v, t} dv\right)dt  \notag \\ %} %_{\text{term debt due at $t$}} %} 
	& & \quad + {\sum_{j\neq s_{t-}} \left(e^{-\eta^*(s_{t-}, j)}-1 \right)\Pi_{t-}^* \left(d {\mathcal J}_t^{(s_{t-}, j)} - \lambda_{s_{t-}, j}e^{\eta^*(s_{t-}, j)}d t \right)} \,.
	%r_{0, t}\Pi_t dt + \left( \left(\int_{s\geq t} q_{t, s}\partial_tb_{t, s} ds \right) dt - \left(\int_{0}^t \partial_u b_{u, t} du\right)dt \right) \,.
	%		&	\tau_tN_{0, t} dt +   \left(\int_{s\geq t} q_{t, s}\partial_tb_{t, s} ds \right) dt  = G_t dt +  b_{t, t}dt \, . 
	\end{eqnarray} 
	This policy requires  continuous adjustments to the term debt structure. That makes it challenging to formulate the Ramsey problem recursively. %to persistently adjust its term debt for all maturities. 
	%\NW{For cadlag processes, fiscal deficits always equal zero at all $t$. Incorrect} 
\end{remark}
%his comparison between our and \cite{LucasStokey1983}'s portfolio management strategies are fundamentally different.

Relative to   \cite{LucasStokey1983}, our dynamic portfolio strategy  has two key features: (1) distinguishing sharply  between instantaneous debt and term debts, and (2) a menu of state-contingent hedging contracts involving minimal number of securities in the spirit of one-period Arrow securities \citep{arrow1964role}.
%
%
%\TS{Tom edit and move next sentence.}
%As we show in the next section, these two features of our state-contingent dynamic portfolio strategy motivate us to introduce a minimal degree of commitment into our setting, which allows us to implement the Ramsey plan, provided that it exists.  




\section{Implementing a Ramsey Plan \label{implementation}}
We now  answer  a question that concerned \cite{LucasStokey1983} and \cite*{debortoli2021optimal}: how to implement the Ramsey plan sequentially under a timing protocol
that attributes ``partial commitment'' to a sequence of government decision makers.\footnote{The concluding section of \cite{lucas1982optimal}, the original working paper version of \cite{LucasStokey1983}, contains a fascinating discussion
of ``partial commitment'' and the way it shaped monetary-fiscal policy decision making arrangements in the U.S.~during the first Washington administration.}
%\footnote{Our work is  related  to   \citet{clymo2020fiscal}  who study how to implement a Ramsey plan  under ``Limited-Time Commitment" (LTC). They show that  $N$ periods of commitment    are required to   implement a Ramsey plan in a 
%setting with  $N$-period debt. Since we let   the  initial debt structure $\{ b_{-1,i} \}_{i=0}^\infty$ be an arbitrary sequence as in  \cite{LucasStokey1983},  \citeauthor{clymo2020fiscal}'s  LTC  implementation requires the number of periods  of commitment equal  
%the longest debt maturity. 
%By way of contrast,  our  implementation of the Ramsey plan using a short-term debt management strategy  works  for any initial term debt  structure.} 
%\TS{Tom: bring in quote from GGHH LS's unpublished concluding section.}
We adopt a timing protocol that    adds a speck of  partial commitment to a continuation Ramsey planner's choice of a tax policy.
%\footnote{This is a continuous time counterpart of a timing protocol used in the discrete-time model of \cite*{JiangSargentWang2025}.} 
This lets us implement  Ramsey plans
that Lucas and Stokey can  not.% 
%\TS{GGHH {cite the JPE submission paper}}.}  
%Next, we turn to implementations of Ramsey plans and show that using instantaneous debt only implements a Ramsey plan, provided that it exists.


%
%\begin{itemize}
%	\item We provide an implementation strategy that uses only instantaneous debt to roll over and an associate unique hedging policy to implement the Ramsey plan
%	\item We provide alternative implementation strategy that uses both instantaneous debt and term debt to roll over.
%	\item We argue the most efficient implementation strategy requires that $X_t$ is smooth in time
%\end{itemize}


\subsection{Partial Commitment  and Implementation}
In order to  induce  an immediate successor government to  confirm a current government's planned  tax, 
% an instance ahead
%We now introduce an amendment to Lucas and Stokey's  timing protocol 
we infinitesimally constrain  Lucas and Stokey's (1983)   governments in the following way. %\footnote{\cite*{jiang2025implementing} analyze a  discrete-time counterpart of this timing protocol.} %Below we introduce a new timing protocol that is necessary for our new implementation.
Let  $\tau_{t-,t}$ denote a time-$t$ tax rate applied over the interval $(t, t+dt)$ that is preferred by the government in office for the interval $(t-dt, t)$, which we index as the time $(t-)$ government. 
%
%We now introduce an amendment to Lucas and Stokey's  timing protocol by strengthening their government's ability to commit infinitesimally in the following way. %Below we introduce a new timing protocol that is necessary for our new implementation.
%Let  $\tau_{t-,t}$ be a time-$t$ tax rate that is preferred by a time-$t-$ government.
\begin{assumption}[Instantaneous partial commitment]\label{ass:timing}
	{A  time-$t-$} government sets a time-$t$ government's current period tax rate~$\tau_{t,t}$: % is dictated by its predecessor
	\begin{equation}\label{eq:timing}
	\tau_{t,t} = \tau_{t-,t}.
	\end{equation}
	
\end{assumption}

%\newpage 
%
%\begin{eqnarray}
%	V_t &=& E_t \int_t^{\infty} e^{-\rho (s-t)} U(C_s) d s  \\
%\int_0^t U(C_s) d s + e^{-rt}	V_t &=& E_0 \int_t^{\infty} e^{-\rho s} U(C_s) d s  \\ 
%d V_t &=& d t + d Z_t  \\  
%\Pi_t &=&\int_{0}^t\, \frac{q_{0, u}}{q_{0, t}} \left(b_{0-, u} +G_u-\tau_{0, u} N_{0, u} \right)du \, , \\ 
%	d\Pi_t  
%	&=& \quad r^{}_{0, t-}\Pi_{t-}dt 
%	+ \left(b_{0-, t-} + G_{t-} - \tau_{0, t-} N_{0, t-} \right)dt   \notag \\
%	&+&  \sum_{j\neq s_{t-}} \left(e^{-\eta(s_{t-}, j)} - 1 \right)\Pi_{t-} 
%	\big(d {\mathcal J}_t^{(s_{t-}, j)} 
%	- \lambda_{s_{t-}, j} e^{\eta(s_{t-}, j)} dt \big) \, , \\ 
%	\Psi_t &=&E_t \int_t^\infty\, \frac{q_{0, u}}{q_{0, t}} \left(b_{0-, u} +G_u-\tau_{0, u} N_{0, u} \right)du \, , \\ 
%	d\Psi_t  
%	&=& \quad r^{}_{0, t-}\Psi_{t-}dt 
%	+ \left(b_{0-, t-} + G_{t-} - \tau_{0, t-} N_{0, t-} \right)dt   \notag \\
%	&+&  \sum_{j\neq s_{t-}} \beta(s_{t-}, j)% \left(e^{-\eta(s_{t-}, j)} - 1 \right)\Pi_{t-} 
%	\big(d {\mathcal J}_t^{(s_{t-}, j)} 
%	- \lambda_{s_{t-}, j} e^{\eta(s_{t-}, j)} dt \big) \, ,\\
%	d X_t &=& d (\Psi_t U_{C,t})
%%d \Pi_t &=& 
%\end{eqnarray} 
%
%sub problem for HH: 
%\begin{eqnarray}
%	\frac{e^{-\rho t} U_C(C_t)}{U_C(C_0)} &=& e^{-\int_0^t r_s d s } \quad for \quad all \quad \omega \\ 	
%	E_{t-1} \frac{e^{-\rho} U_C(C_t)}{U_C(C_{t-1})} &=& E_{t-1}e^{-\int_{t-1}^t r_s d s } =q_{t-1,t}\\ 
%	E_{t-1} MRS_{t-1, t} R_{t-1, t} &=& 1 \quad  \frac{1}{R_{t-1, t}} =q_{t-1,t}\\ 
%	E_{t-1} [M_{t-1,t} R_{t-1, t}]&=& 1 \quad  \frac{1}{R_{t-1, t}} =q_{t-1,t}\\ 
%	E_{t-1} [M_{t-1,t} R^{i}_{t-1, t}]&=& 1 \quad   \quad for \quad all i \quad under \quad complete \quad markets  \\ %  \frac{1}{R_{t-1, t}} =q_{t-1,t}\\ 
%	\frac{e^{-\rho t} U_C(C_t)}{U_C(C_0)} &=& e^{-\int_0^t r_s d s } \\ 
%	U(C_t) = V_W(W_t, t, \cdots) 
%\end{eqnarray}
%
%%\Pi_t %= \int_{0}^t\,   e^{\int_s^{t} r_ud u} \left(b_{0, s} +G_s-\tau_{0, s} N_{0, s} \right)ds  \quad  \Pi_0 = 0\, , 
%
%\newpage 

\noindent Thus,  the time~$t$ tax rate  that applies to taxable income  over a short time interval $(t, t+dt)$ is set by the instantaneous predecessor of the time-$t$ government, also known as the continuation  planner at time $t-$. Via  the household's first-order necessary condition, this timing protocol induces
\begin{equation}\label{eq:Ctt}
C_{t,t} = C_{t-,t},
\end{equation}
where $C_{t-,t}$ is the time-$t$ consumption effectively chosen by the time-$t-$ government.

To achieve this outcome,  the government at time~$t-$ hands over  to the time-$t$ government, a portfolio of financial assets that delivers a balance of  instantaneous debt  $B_t$ for each possible state $s_t$ at $t$  and a stream of term debt services $\{b_{t-,u};\, u \ge t\}$. Importantly,  %with a combination of non-state-contingent debt and state-contingent hedging contracts described in XXX, 
we show that there exists a financial portfolio  with the  following properties for  each  $s_t$ that always implements the Ramsey plan, provided that it exists: %satisfies
\begin{equation}\label{eq:implement_strategy}
B_t = \Pi_t^* \quad \text{and} \quad b_{t-,u} = b_{0-,u};\quad u \ge t , 
\end{equation}
where $ \Pi_t^*$ is the state-contingent cumulative liability whose evolution is given in \eqref{dPiast} under the Ramsey plan. To be precise, using the  dynamic trading strategy that features a combination of state non-contingent debt and a portfolio of Arrow securities described in Subsection \ref{subsec:financing1}, we can show that \eqref{eq:implement_strategy} holds for all paths. %that makes \eqr
As \eqref{eq:implement_strategy} requires that the term debt service flow inherited by the time-$t$ government equals  the tail  $\{b_{0-,u};\, u \ge t\}$ of the initial term debt services stream at all $t$, %as a result  %that  the cumulative liabilities under the Ramsey plan up to time~$t$ $\Pi_t^*$ has to be fully absorbed by the balance of the 
the government's instantaneous debt balance $B_t$ must perfectly track the government's cumulative liability $\Pi_t^\ast$ for all paths. This is achieved by properly choosing the dynamically changing hedging contract holdings $\Delta_{t}$ given in \eqref{Delta}. 

%In our implementation,   the government at time $t-$  leaves a debt structure $\{{B}_t, {b}_{0, u}; u\geq t\}$ to the time $t$ government, where 	\begin{equation} \label{key4} 
%{B}_{t} = \Pi_{t}^* \;\; \text{and} \;\;  {b}_{t-, \, u}=b_{0,\, u}; \;\, u\geq t .
%\end{equation}

%\NW{\paragraph{Debt Management with Arrow Securities.}
%Starting with $B_0=\Pi_0=0$, we can This debt policy can be achieved by dynamic debt management as follows
%\begin{eqnarray} \label{finB2}
%d B_t  %=  -\Pi_t\frac{dq_{0, t}}{q_{0, t}} + \left(b_{0, t} +G_t-\tau_{0, t} N_{0, t}\right)dt 
%& = &   {r_{0,t-}^*} B_{t-} d t  + \left(b_{0, t-} +G_{t-}-\tau_{0, t-} N_{0, t-}\right)dt \notag \\
%&  & \quad\quad + \sum_{j\neq s_{t-}} \alpha(s_{t-}, j)B_{t-} \left(d {\mathcal J}_t^{(s_{t-}, j)} - \NW{ \lambda_{s_{t-}, j}e^{\eta^*(s_{t-}, j)}dt} \right)    
%\end{eqnarray}
%in which $\{ \alpha(s_{t-}, j); j\neq s_{t-}\}$ represents the time $t-$ continuation planner's arrow security holdings such that
%\begin{equation} \label{hedging}
%\alpha(s_{t-}, j) = \left( e^{-\eta^*(s_{t-}, j)}-1 \right) , \forall \, j\neq s_{t-} \,.
%\end{equation}  
%The Arrow securities effectively act as the state-contingent instantaneous debt.  With this strategy, at time $t$, the continuation planer's short-term debt balance is effectively
%\begin{equation}
% B_t = {B}_{t-}+ \sum_{j\neq s_{t-}}\, \left( e^{-\eta^*(s_{t-}, j)}-1\right) {B}_{t-} d {\mathcal J}_t^{(s_{t-}, j)} = \sum_{j\neq s_{t-}}\, e^{-\eta^*(s_{t-}, j)} \Pi^*_{t-} d {\mathcal J}_t^{(s_{t-}, j)}   =\Pi_t^*\,.
%\end{equation}
%}
%

%\begin{assumption} [{\bf local commitment condition}]\label{assump_localcommit}
%	
%	The time $t$ continuation planner must admisister tax rate $\tau_{t-, t}$, which is chosen by the predeceasing government, for time $t$.
%\end{assumption} 

Under Assumption \ref{ass:timing} that specifies  an instantaneous commitment,  the dynamic trading strategy described above allows us to ensure that  $X_t$  defined in \eqref{X} equals $X_t^\ast$  under the Ramsey plan. That is, 
the following identity holds:
%	For any time $t$ government that inherits from previous government an instantaneou debt $\widehat{B}_t$ which dues immediately, when the government has a reoptimization option to choose consumption $C_{t, t}$ and labor supply $N_{t, t}$, it must commits to the purchasing power of the instantaneous debt measured by the marginal unitlity such that 
%$ \widehat{B}_tU_{C, t}(C_{t, t}, N_{t, t})=B_{t-}U_{C, t-}(C_{0, t-}^*, N_{0, t-}^*)$.
\begin{equation} \label{key5} 
{B}_tU_{C, t}(C_{t, t}, N_{t, t})=\Pi_t^* U_{C, t}(C_{t-, t}, N_{t-, t}) = X_t^*,   \end{equation}  
% \begin{equation} \label{key5} 
%\widehat{B}_tU_{C, t}(C_{t, t}, N_{t, t})=B_{t-}U_{C, t-}(C_{0, t-}^*, N_{0, t-}^*)	 ,\end{equation}  
where $C_{t, t}$ and $N_{t, t}   $
are the time-$t$ government's choices.  
%and $B_{t-}$ is the instantaneous debt balance that finances the government's cumulative deficits up to time $t-$. 
%The following assumption summarize this local commitment condition. 
%The purchasing power of the  instantaneous debt inherited by a time $t$ government is preserved, i.e., 
Note that if the Ramsey plan has been  implemented  until time $t-$, then $C_{t-, t}=C_{0, t}^*$ and $N_{t-,t}=N_{0,t}^*$. Now we are ready to state  our  implementation result and then  prove it. 
\begin{proposition}\label{propimplement}
	%	The Ramsey plan can be  implemented by leaving $\{b_{0-,u}\}$ untouched, hedging according to (\ref{hedging}), and adjusting instantaneous debt according to  (\ref{eq:implement_strategy}) under the local commitment condition stated in Assumption \ref{ass:timing}. 
	Under our Assumption~\ref{ass:timing} timing protocol, a Ramsey plan can be implemented by accepting the tail of  initial debt service stream  $\{b_{0-,u}\}$ and adjusting instantaneous debt according to~\eqref{eq:implement_strategy} .
	%Under Assumption \ref{assump_localcommit}, the instantaneou debt management policy is .
\end{proposition}



\begin{proof}
	We  use a  recursion to prove the above proposition. 
	Under the assumption that the Ramsey plan is implemented  up to time $t-$, it is sufficient if we prove that the Ramsey plan is also implemented at time $t$. %\footnote{At $t=0$, the Ramsey planner is also  a time $0$ continuation Ramsey planner.} 
	We can show that  at time $t$ the continuation Ramsey planner, confronting ${B}_t = \Pi_t^*$ and ${b}_{t-, u}=b_{0, u}$ for all $u\geq t$ and the inherited hedging contract holdings $\{ \alpha(s_{t-}, j); j\neq s_{t-}\}$, will confirm the continuation of the Ramsey plan by choosing $C_{t,u}=C_{0,u}^\ast$ and $N_{t,u}=N_{0,u}^\ast$ for all $u\geq t$. We can prove the preceding claim by subdividing time-$t$ continuation Ramsey planner's  problem into two subproblems. 
	
	
	First, the  time-$t$ continuation Ramsey planner chooses  $\{C_{t, u}, u\geq t\}$ to attain  %(\ref{Vcontinue}). 
	\begin{equation}\label{Vimplementation}
	%{V}(X_t, t) = 
	\max_{\{C_{t, u}; u\geq t \}}\; 
	\mathbb{E}_t \left[	\int_t^{\infty} e^{-\rho(u-t)} U({C}_{t, u}, {C}_{t, u}+G_u) d u \right]  ,
	\end{equation} %where maximization 
	subject to $\widehat{X}_t={B}_tU_{C, t}(C_{t, t}, N_{t, t})$ and the following dynamics for $u\geq t$: 
	\begin{eqnarray}\label{dyhatt}  
	d \widehat{X}_u  & = &  \big(\rho \widehat{X}_{u-}-C_{t, u-}U_{C, u-} -N_{t, u-}U_{N, u-}  +{b}_{0, u-}U_{C, u-}\big)du\,. % \notag\\ % \\
	%	& & \quad +\WJ{\sum_{k\neq s_{t-}}\left(\widehat{\alpha}_{t-}{(s_{t-}, k)} + e^{\kappa(s_{t-}, k)}-1 \right) X_{t-}\left(\NW{d{\mathcal J}_t^{(s_{t-}, k)}} -\lambda_{s_{t-}, k}dt \right)} \,.
	%&  & +\, .
	\end{eqnarray} 
	%with $\widehat{X}_u = \widehat{B}_uU_{C, t}$. %Here $U_{C, t}, U_{N, t}$ are short notation for $U_{C, t}(C_{u, t}, N_{u, t}), U_{N, t}(C_{u, t}, N_{u, t})$, respectively. 
	Let $\widehat{V}(\widehat{X}_t, s_t; \Phi_t)$ denote  the  value function for the continuation planner defined in (\ref{Vimplementation}). Similar to our analysis in Section \ref{sec:bellmanramsey}, we can show that $\widehat{V}(\widehat{X}_t, s_t; \Phi_t)$ is additively separable: $\widehat{V}(\widehat{X}_t, s_t; \Phi_t)= -\Phi_t \widehat{X}_t + \widehat{H}(s_t; \Phi_t)$, where $\widehat{H}(s_u; \Phi_t)$ for $u\geq t$ solves the following differential equation:%\footnote{\WJ{See more details in Appendix XX.}}
	\begin{align}\label{HJBHtt} %\nonumber
	\rho \widehat{H}(s_u; \Phi_t)  = &\;
	\max_{C_{t, u}} \, U(C_{t, u},C_{t, u}+G_u) 
	+ \Phi_t \big[C_{t, u} U_{C, u}+(C_{t, u}+G_u) U_{N, u}-b_{0-, u} U_{C, u} \big]  \notag\\
	& \quad\quad \;\,	+\sum_{j\neq s_u}\lambda_{s_u , \, j}\left[ {\widehat{H}(j; \Phi_t)}-\widehat{H}(s_u; \Phi_t) \right]\, , 
	\end{align}
	where $C_{t,u} =C(s_u; \Phi_t)$, and $C(s_u; \Phi_t)$ is  the same as the function defined in (\ref{optC00}).
	
	%	a differential equation (XX)  in Appendix XX, and 
	%	that   $C_{t,u} =C(s_u; \Phi_t)$  where $C(s_u; \Phi_t)$ is  the same as the function defined in (\ref{optC00}). This is because the initial debt term structure $\{b_{0,t}; t\geq 0\}$ is unchanged in our implementation with instantaneous debt only. 
	
	
	Under the local commitment condition given in  Assumption \ref{ass:timing}, the time $t$ continuation Ramsey planner has to choose $C_{t, t}=C_{t-, t}=C_{0, t}^*$. Because 1.) ${B}_t =\Pi_t^*$ under our implementation and 2.) $\widehat{X}_t = {B}_tU_{C, t}(C_{t, t}, N_{t, t})= \Pi_{t}^*U_{C, t}(C_{0, t}^*, N_{0, t}^*) = X_t^*$, we have $\widehat{X}_t= X^*_t$. 
	
	
	
	
	
	
	{We now turn to establishing that $\Phi_t^* = \Phi_0^*$.} We start by  proving that the time-$t$ continuation Ramsey planner  chooses $C_{t,u} = C_{0, u}^*$ for all $u > t$. Because the optimal continuation plan $\{C_{t, u}^*;\, u > t\}$ satisfies $C_{t,u}^* = C(s_u; \Phi_t^*)$ for some $\Phi_t^* > 0$ and the Ramsey plan satisfies $C_{0,u}^* = C(s_u; \Phi_0^*)$, we need to verify that $\Phi_t^* = \Phi_0^*$, which 
	we do via contradiction. Suppose that there exists an $\Phi_t^* \ne \Phi_0^*$ such that the associated  continuation Ramsey plan yields a higher value than the tail of the Ramsey plan, meaning that
	\begin{equation}\label{eq:contradict_ineq}
	\mathbb{E}_t \int_t^\infty e^{-\rho(u-t)} U(C_{t,u}^*, C_{t,u}^*+G_u)\, du > \mathbb{E}_t \int_t^\infty e^{-\rho(u-t)} U(C_{0,u}^*, C_{0,u}^*+G_u)\, du,
	\end{equation}
	where $C_{t,u}^* = C(s_u; \Phi_t^*)$. Recall that the time-$t$ continuation Ramsey planner chooses an  allocation $\{C_{t,u}^*;\, u > t\}$ that satisfies the following continuation implementability condition:
	\begin{equation}\label{eq:cont_implement}
	\mathbb{E}_t	\int_t^\infty e^{-\rho(u-t)}\bigl(C_{t,u}^*\, U_{C,u} + (C_{t,u}^*+G_u)\, U_{N,u}\bigr)\, du \ge \mathbb{E}_t \int_t^\infty e^{-\rho(u-t)} b_{0-,u}\, U_{C,u}\, ds + B_t\, U_{C,t} \, .
	\end{equation}
	Construct a plan $\{\widetilde{C}_{0,s};\, s \ge 0\}$ in which  the allocation from~0 to~$t$ follows the Ramsey plan, $\widetilde{C}_{0,s} = C_{0,s}^*$ for $s \in [0,t]$, but the allocation from time~$t$ onward follows the time-$t$ continuation Ramsey plan: $\widetilde{C}_{0,s} = C_{t,s}^*$ for $s \in (t,\infty)$. Using~\eqref{eq:cont_implement} and $B_t = \Pi_t^*$, we can verify that $\{\widetilde{C}_{0,s};\, s \ge 0\}$ satisfies the implementability condition~\eqref{impletconstraint} for the Ramsey plan. Next, using~\eqref{eq:contradict_ineq}, we can show that this proposed plan $\{\widetilde{C}_{0,s};\, s \ge 0\}$ yields a value for the household that is higher than the value under the Ramsey plan:
	\begin{align*}
	&  \mathbb{E}_0	\int_0^\infty e^{-\rho u} U(\widetilde{C}_{0,u}, \widetilde{C}_{0,u}+G_u)\, du \notag\\
	&= \mathbb{E}_0 \left[ \int_0^t e^{-\rho u} U(\widetilde{C}_{0,u}, \widetilde{C}_{0,u}+G_u)\, du + \mathbb{E}_t\int_t^\infty e^{-\rho u} U(\widetilde{C}_{0,u}, \widetilde{C}_{0,u}+G_u)\, du \right] \\
	&= \mathbb{E}_0 \left[ \int_0^t e^{-\rho u} U(C_{0,u}^*, C_{0,u}^*+G_u)\, du +  \mathbb{E}_t \int_t^\infty e^{-\rho u} U(C_{t,u}^*, C_{t,u}^*+G_u)\, du \right] \\
	&>  \mathbb{E}_0 \left[ \int_0^t e^{-\rho u} U(C_{0,u}^*, C_{0,u}^*+G_u)\, du + \mathbb{E}_t \int_t^\infty e^{-\rho u} U(C_{0,u}^*, C_{0,u}^*+G_u)\, du \right] \\
	&=  \mathbb{E}_0 \int_0^\infty e^{-\rho u} U(C_{0,u}^*, C_{0,u}^*+G_u)\, du.
	\end{align*}
	This contradicts that $\{C_{0,u}^*\}$ is the Ramsey plan. We thus conclude $\Phi_t^* = \Phi_0^*$ and  that $C_{t,u}^* = C(s_u; \Phi_t^*) = C(s_u; \Phi_0^*) = C_{0,u}^*$ for all $u \ge t$.
	
	
\end{proof}



\section{Quantitative Examples}

The two quantitative examples highlight how the Ramsey planner manipulates equilibrium prices and marginal utilities to manage debt costs. A recurring theme is that the state variable $X_t = \Pi_t\, U_{C,t}$ remains a continuous function of time even when the exogenous term debt process or the government expenditure process jumps discontinuously across regimes.

To illustrate how our model works, we  posit  a regime-switching model with two states $s_t=\{H, L\}$.  We endow our representative agent  with the instantaneous  utility function  %\begin{equation}\label{DNYUtility}
$U(C, N) = \log C - \eta \frac{N^\gamma}{\gamma}  \, , $ 
%\end{equation}
where $\eta>0$ and $\gamma\geq 1$, the same  one employed by  \citet*{debortoli2021optimal}.
In subsection \ref{subsec:b},  we study a setting with  a stochastic $b_{0-,t}$, while in subsection \ref{subsec:G}, we study a  setting with a stochastic $G_{t}$. In both subsections,  we set $\rho=0.02$, $\eta=1$, the transition rate $\lambda_{HL}$ from  $H$ to $L$  to $0.2$ and the transition rate $\lambda_{LH}$  from  $L$  to $H$ to $0.3$. 

%\subsection{Markovian $\{b_{0-,t}\}$ }
%
%Figure~1 shows a Ramsey plan when 
%%Consider an exponential debt term structure: 
%$b_{0-,t} = b_0 e^{-\theta t}$, where $b_0 = 0.2$ and $\theta = 0.6$, and that  $G_t = 0.2$ for all~$t$. The Ramsey tax rate $\tau_{0,t}^*$ increases over time (panel~B), which implies that consumption $C_{0,t}^* = 1-\tau_{0,t}^*$ and the  labor supply $N_{0,t}^* = C_{0,t}^*+G_t$ both decrease over time. The equilibrium interest rate is $r_{0,t}^* = \rho + g_{0,t}^*$ where $g_{0,t}^* = \dot{C}_{0,t}^*/C_{0,t}^*$ is the equilibrium consumption growth rate,
%which is % Note that the consumption growth rate $g_{0,t}^*$ is 
%negative:\footnote{ $C_{0,t}^* = 2b_0 e^{-\theta t}\Phi_0^*/(\sqrt{1+4b_0 e^{-\theta t}\Phi_0^*(1+\Phi_0^*)}-1)$ and $\Phi_0^* = 0.2652$. Therefore, $g_{0,t}^* = dC_{0,t}^*/dt / C_{0,t}^* < 0$ because $dC_{0,t}^*/dt = -b_0\Phi_0^*\theta e^{-\theta t}/\sqrt{1+4b_0 e^{-\theta t}\Phi_0^*(1+\Phi_0^*)} < 0$.} {it is less than  $-\rho$ at $t=0$, then gradually } increases over time, eventually {approaching} zero. Coincidentally, the interest rate is negative at $t=0$,  increases over time, and eventually approaches $\rho = 0.02$ (see panel~C).
%
%The Ramsey planner {front-loads} consumption so heavily  that increases in output $N_{0,t}^* = C_{0,t}^*+G_t$ dominate decreases in the tax rate $\tau_{0,t}^* = 1-C_{0,t}^*$, making the primary surplus $S_{0,t}^* = \tau_{0,t}^*\, N_{0,t}^* - G_t = (1-C_{0,t}^*)(C_{0,t}^*+G_t) - G_t$ increase over time. Since  $b_{0-,t}$ decreases over time, the wedge $S_{0,t}^*-b_{0-,t}$ is  negative initially, turns positive at $t = 5.8$, and ultimately converges to $0.7\%$ (see panel~D).
%
%By inducing an increasing equilibrium $r_{0,t}^*$ path (panel~C), the planner lowers the time-0 value of its wedge $S_{0,t}^*-b_{0-,t}$ shown in panel~D. As a result, the cumulative liability $\Pi_t^* = \int_0^t e^{\int_0^u r_{0,v}^*\, dv}(S_{0,u}^*-b_{0-,u})\, du$ increases (panel~E). Since   $\Pi_t^*$ and $U_{C,t}^*$ both increase over time,
%the state variable $X_t^* = \Pi_t^*\, U_{C,t}^*$ also increases over time (panel~F). 
%%% [Figure 1 would be included here]



\begin{figure}[h!]
	\includegraphics[width=\linewidth, height=0.7\textheight]{figure_dynamic/dynamic_b.eps}
	\caption{\label{ramsey_1}{Ramsey plan in a  two-state  setting with high and low levels of initial term debt.} Parameter values: {$b_H=0.2, b_L = 0, G_t=0.2$ for $t\geq 0$.} %$\lambda_{HL}=0.2, \lambda_{LH}=0.3$. 
		% \TS{$C_{0, t}=1-\tau_{0, t}, N_{0, t}=C_{0, t}+G_t$} 
	}%$\{b_{0, t}\}$ with $b_0=0.2$.} \WJ{$
\end{figure}






\subsection{State-Dependent Initial Term Debt $\{b_{0-,t}\}$ } \label{subsec:b}

%This subsection analyzes an example where the initial debt structure $\{b_{0-,t}\}$ is cyclical. 
%$b_{0-,t} = b_L$ if $s_t=L$ setting with two possible states

Consider a setting where the initial  debt is  high or low: $b_{0-,t} = b_H$ if $s_t=H$ and $b_{0-,t} = b_L$ if $s_t=L$ for $t \ge 0$. 
%$b_{0-,t}$, equals $b_H$ and $b_L$ in the two states: in that 
We set   $b_H = 0.2$, $b_L=0$, and  $G_t=0.2$ for all $t$. 
Figure~1 plots outcomes under the   Ramsey plan.  Panel A plots a  path of $b_{0-, t}$. 

%= \tau_L 1_{\{s_t=H\}}+\tau_H 1_{\{s_t=L\}}$, where $
Panel B shows that  the Ramsey tax rate is low (high) when debt is high (low). Specifically, the  tax rate $\tau_{0,t}^* = 0.31$ in state $H$ and  $\tau_{0,t}^* =0.43$ in state $L$.  
This is because the planner wants to reduce the market value of its debt by making marginal utility in the high-debt state low. As $C_{0,t}^* = 1-\tau_{0,t}^*$, low marginal utility, which maps to high consumption, is achieved with a low tax rate. This explains why taxes are low in high-debt states. 

Panel C shows that the  equilibrium interest rate is low  when debt is high. 
Note that  the interest rate $r_{0,t}^*$ is low and   equals $-0.023$ in the high-debt state, but is high and   equals %$%\tau_{0,t}^* =0.43$ in the $s_t=L$ state. 
$0.073$ in the no-debt state. Similarly, this is because the planner wants to make the interest rate low when its debt is high in order to reduce its debt services. %by making marginal utility in the high-debt state low.
%Mathamtically,   
Panel D shows that the primary surplus $S_{0,t}^\ast$ in excess of the initial term debt service $b_{0-,t}$,  $S_{0,t}^\ast-b_{0-,t}$, is low (high) when debt is high (low):   $S_{0,t}^*-b_{0-,t}= 0.131$ in the high-debt state and  $S_{0,t}^*-b_{0-,t}= -0.126$ in the no-debt state.\footnote{
We obtain this result as follows. First,    $C_{0,t}^* = \bigl(\sqrt{1+4b_{0-, t}\Phi_0^*(1+\Phi_0^*)}+1\bigr)/(2(1+\Phi_0^*))$, which increases in $b_{0-,t}$. 
Second, note that $S_{0,t}^*-b_{0-,t}=\tau_{0, t}^*N_{0, t}^*-b_{0-,t} = C_{0, t}^*(1-C_{0, t}^*-G_t)-b_{0-,t}$. Third,  using $\frac{\partial C_{0, t}^*}{\partial b_{0-, t}}=\frac{\Phi_0^*}{2C_{0, t}^*(1+\Phi_0^*)-1}>0$,  we conclude  that $\frac{\partial (S_{0,t}^*-b_{0-,t})}{\partial b_{0-, t}} = \frac{\partial C_{0, t}^*}{\partial b_{0-, t}}\left(2C_{0, t}^*\Phi_0^*-G_t \right)-\Phi_0^*-1 <0 $. 
% is increasing in $b_{0-,t}$
} %Lagrangian multiplier  $\Phi_0^* = 0.7494$.} %decreases in $b_{0, t}$. Note that 
%\WJ{Again, this is because the planner wants to make the interest rate low when its debt is high in order to reduce its debt services.}  



Panel~E shows that the implied cumulative liability~$\Pi_t^*$  (with no rescheduling of term debts) decreases in the no-debt state and increases in the high-debt state. Note that when the state transitions, $\Pi_t^*$  generally jumps. %This is because the government wants to smooth the purchasing power of %This reflects the planner's preferences to smooth  
Panel~F shows that the state variable $X_t = \Pi_t^*\, U_{C,t}^*$, which tracks the purchasing power of $\Pi_t^\ast$, also decreases in the no-debt state and increases in the high-debt state. Notice that,  unlike $\Pi_t^*$, $X_t$ is continuous even when the state transitions. This is due to the household's optimal response to the government's tax plan and the associated equilibrium prices, all of which are determined at $t=0$. %By optimally manipulating the household's marginal utility, the planner ensures that 
The  jump of $\Pi_t^\ast$  is offset  by a jump of $U_{C,t}^\ast$, leaving $X_t = \Pi_t^*\, U_{C,t}^*$ unchanged when the state transitions. 

%stochastic discount factor 


%are countercyclical with respect to the~$b_{0-,t}$ process. Specifically, the government accumulate savings when debt payment is low and accumulate liabilities when debt payment is high.


%$r_{0,t}^* = r_L 1_{\{s_t=H\}}+r_H 1_{\{s_t=L\}}$, where It is also countercyclical and reaches a low level when $\{b_{0-,t}\}$ reaches a high level 

%
%so consumption process is procyclical and reaches a high level when $\{b_{0-,t}\}$ also reaches a high level.\footnote{$C_{0,t}^* = \bigl(\sqrt{1+4b_{0-, t}\Phi_0^*(1+\Phi_0^*)}-1\bigr)/(2(1+\Phi_0^*))$, which is cyclical with the same period as~$b_{0-,t}$. Lagrangian multiplier  $\Phi_0^* = 0.7494$.}
%%
%%They evidently differ from the peaks/troughs in panels~A, B, and~D for~$b_{0-,t}$, $\tau_{0,t}^*$ and with the wedge between primary surplus and~$b_{0-,t}$: $S_{0,t}^*-b_{0-,t}$.\footnote{The wedge shares peaks and troughs with the tax rate $\{\tau_{0,t}^*\}$ process.} 
%%Panel~D also shows that the primary surplus exceeds~$b_{0-,t}$ when $b_{0-,t} = b_L$ and otherwise when $b_{0-,t} = b_H$.
%
%{The tax rate and the wedge $S_{0,t}^*-b_{0-,t}$ processes move in an opposite direction from  the~$b_{0-,t}$ process because the marginal utility (and hence bond price) is low when $\{b_{0-,t}\}$ is high. Such covariation minimizes  the value of the government's cumulative liability.}
%
%\NW{Delete the above if not needed.} 

%% [Figure 2 would be included here]



\begin{figure}[h!]
	\includegraphics[width=\linewidth, height=0.7\textheight]{figure_dynamic/dynamic_g.eps}
	\caption{\label{ramsey_2}{Ramsey plan in a  two-state  setting with high and low levels of government spending.} Parameter values: {$b_{0-, t} = 0.1$ for $t\geq 0$, $G_H=0.2$, and $G_L=0.15$.}		% \TS{$C_{0, t}=1-\tau_{0, t}, N_{0, t}=C_{0, t}+G_t$} 
	}%$\{b_{0, t}\}$ with $b_0=0.2$.} \WJ{$
\end{figure}




%\subsection{Two Regime $\{G_t\}$}

\subsection{State-Dependent Government Spending  $\{G_t\}$} \label{subsec:G}

%the initial  debt is  high or low: $b_{0-,t} = b_H$ if $s_t=H$ and $b_{0-,t} = b_L$ if $s_t=L$ for $t \ge 0$. 
%$b_{0-,t}$, equals $b_H$ and $b_L$ in the two states: in that 
%We set   $b_H = 0.2$, $b_L=0$, and  $G_t=0.2$ for all $t$.  Figure~1 plots the  Ramsey plan solution and implications in this setting.  Panel A plots a  path of $b_{0-, t}$. 

%This subsection considers a setting



Now  government spending switches between  high and low: $G_{t} = G_H$ if $s_t=H$ and $G_{t} = G_L$ if $s_t=L$ for $t \ge 0$. 
%$b_{0-,t}$, equals $b_H$ and $b_L$ in the two states: in that 
We set   $G_H = 0.2, G_L=0.15$  and  $b_{0-,t}=0.1$ for all $t$. 
Figure~2 plots outcomes under the Ramsey plan.  Panel A plots a  path of $G_{t}$. 



Panel B shows that the  tax rate is constant: $\tau_{0,t}^* = 1-C_{0,t}^* = 0.33$ for all~$t$. This is because the maximand in~\eqref{optC00} yields an optimal~$C$ that is independent of~$G_t$ when $b_{0-,t}$ is constant. This tax-smoothing result echoes \cite{barro1979determination}'s finding that  intertemporal properties of an optimal tax rate are disconnected from intertemporal properties of the government expenditure process.
%% CHANGE NOTE: Added economic explanation for why C* is constant; original just
%% stated the result.
Since   consumption~$C_{0,t}^*$ is constant, the equilibrium interest rate  equals the household's discount rate at all $t$: $r_{0,t}^* =  \rho =2\%$ (see panel~C). Panel~D confirms that $S_{0,t}^*-b_{0-,t}$ moves one-for-one in the opposite direction from government spending $G_t$,  as both the debt service $b_{0-,t}$ and the tax rate $\tau_{0,t}^\ast$ are constant over time. 


%reaching a high level when the~$G_t$ process reaches a low level. %Solid dots indicate troughs/peaks in panels~A and~D. Triangles in panel~D indicate points where $S_{0,t}^* = b_{0-,t} = 0.1$.  

%, which  


Unlike  Figure~1, Panel~E shows that the government's cumulative liability is {\it continuous} even as the state transitions. This  follows from the tax (and consumption) smoothing  in this setting. Accruing $-(S_{0,u}^*-b_{0-,u})$ from $0$ to $t$ gives $\Pi_t^* = - \int_0^t e^{\rho (t-u)}(S_{0,u}^*-b_{0-,u})\, du$. Thus, $\Pi_t^*$ increases (decreases) when  $S_{0,u}^*-b_{0-,u}$ is  negative  (positive). As $U_{C,t}^*=1/C_{0,t}^*=1/0.67$ for all $t$, $X_t= \Pi_t^* U_{C,t}^\ast= 1.48  \Pi_t^*$, thus making Panel~F  proportionally track Panel~E at all $t$. %in this case.  %= 

%decreases when $S_{0,t}^*-b_{0-, t}= 0.019 >0$ and $G_t=G_L$ and increases when $S_{0,t}^*-b_{0-, t}= -0.015 <0$ and $G_t=G_H$.\footnote{This follows from $d\Pi_t^*/dt = e^{\rho t}(S_{0,t}^*-b_{0-,t})$.} The state variable $X_t^* = \Pi_t^*\, U_{C,t}^*$ moves synchronously as~$\Pi_t^*$  (see the panels~E and~F) because constant consumption implies constant~$U_{C,t}^*$,.






\section{Alternative  Recursive Formulation}\label{sec:indeterminacy} %Via Dynamic 

This section replaces the backward-looking cumulative liability $\Pi_t$ with a forward-looking promised liability $\Omega_t$ and constructs an alternative state variable $Y_t = \Omega_t\, U_{C,t}$. The law of motion for $Y_t$ contains a coefficient on a martingale increment that is indeterminate, reflecting the government's freedom in choosing Arrow securities.  Despite this indeterminacy, the value function retains the same additively separable structure and the Ramsey plan is unique: the optimal consumption and tax paths are the same as those found in section~\ref{sec:bellmanramsey}.

We proceed in three steps.
{ 
	1. For a given consumption process, we define  a process $\{\Omega_t;\, t \ge 0\}$ that tracks the government's promised (present discounted) value of all future liabilities in the absence of any debt restructuring. 2. From the  $\{\Omega_t;\, t \ge 0\}$ process, we  construct an alternative state variable process $\{Y_t;\, t \ge 0\}$. 3. We frame  a dynamic program with $\{Y_t;\, t \ge 0\}$ as a state variable and a state-$s_t$-dependent value function, $K(Y_t,s_t)$, that is additively separable: $K(Y_t,s_t) = -\Psi_0\, Y_t + H(s_t;\, \Psi_0),$ where  $\Psi_0$ is a scalar to be determined. 
} 

In addition to providing an alternative approach for solving the Ramsey plan, we also discuss difficulties in implementing the Ramsey plan if Arrow securities are not chosen properly. 





\subsection{{Cumulative Liabilities with no Debt Restructuring}}



For a given competitive equilibrium (see Definition \ref{defn1}), 
we can construct the following  $\{\Omega_t\}$ process: 
\begin{equation} \label{Omega}
\Omega_t %= \int_{0}^t\,   e^{\int_s^{t} r_ud u} \left(b_{0, s} +G_s-\tau_{0, s} N_{0, s} \right)ds 
= \mathbb{E}_t \int_{t}^\infty\, \frac{q_{0, u}}{q_{0, t}} \left(b_{0-, u} +G_u-\tau_{0, u} N_{0, u} \right)du \, , \quad  \Omega_0 = 0\, , 
\end{equation}   
where 
the ratio $q_{0,u}/q_{0,t}$ discounts the flow $(b_{0-,u}+G_u-\tau_{0,u}\, N_{0,u})\, du$ back to the government's time-$t$ promised liability. 
This $\{\Omega_t\}$ process is the government's time-$t$ discounted value of  its liabilities in the absence of  debt restructuring from $t$ to $\infty$.  


Note that %\{t\geq 0\} 
\begin{equation}
q_{0, t}\Pi_t+q_{0, t}\Omega_t  =  \mathbb{E}_t \bigg[ \int_{0}^\infty\, {q_{0, u}}\left(b_{0-, u} +G_u-\tau_{0, u} N_{0, u} \right)du\bigg]
\end{equation}
is a martingale. Using the martingale representation theorem, we obtain:
\begin{lemma}\label{lemmaOmega0} %\neq s_{t-} _{t-}(s_{t-}, s_{t})
	For a given  $\{\alpha_t; t\geq 0\}$ process set at $t=0$, 	we can express  the dynamics for 
	the  $\{\Omega_t; t\geq 0\}$ process defined in \eqref{Omega}  as follows:
	\begin{align}\label{dOmega0}
	d\Omega_t  
	\quad =& \quad r^{}_{0, t-}\Omega_{t-}dt 
	+ \left(b_{0-, t-} + G_{t-} - \tau_{0, t-} N_{0, t-} \right)dt   \notag \\
	&\quad + \sum_{j\neq s_{t-}} \alpha_{t-} \Omega_{t-} 
	\big[\underbrace{d {\mathcal J}_t^{(s_{t-}, j)} 
		- \lambda_{s_{t-}, j} e^{\eta(s_{t-}, j)} dt}_{\text{increment of a  martingale}} \big] \, ,
	\end{align}
	where $\alpha_{t-}=  \alpha(s_{t-}, j)$ for $j\neq s_{t-}$ and $\eta(s_{t-}, s_t) = \ln\frac{U_{C, t}}{U_{C, t-}}$ .  %\NW{Make it general: }
\end{lemma}

It is important to note that the $\{\Omega_t\}$ process defined in \eqref{Omega} is conceptually different from  the $\{\Pi_t\}$ process defined in \eqref{Pi}, which tracks cumulative liabilities accrued from $0$ to $t$ in the complete absence of debt restructuring.  Comparing their respective dynamics, \eqref{dOmega0} and  \eqref{dPi0},  %We make two observations regarding \eqref{dOmega0}. 
we see that the only difference is the last  term that involves the increment of a martingale, which captures the state transition. For the cumulative liability $\Pi_t$ process,  the coefficient multiplying the martingale increment in   \eqref{dPi0} is unique and equals $\left(e^{-\eta(s_{t-}, j)} - 1 \right)\Pi_{t-} $. By contrast, for the promised liability $\Omega_t$ process, the coefficient multiplying the martingale increment in  \eqref{dOmega0}, $\alpha_{t-} \Omega_{t-} $, is indeterminate. This is 
because the cumulative liability $\{\Pi_t\}$ process is ``backward-looking,'' while the promised liability $\{\Omega_t\}$ process is ``forward-looking.'' That is, there  exists a unique $ \alpha(s_{t}, j) $ process, namely, $ \alpha(s_{t}, j) =\left(e^{-\eta(s_{t}, j)} - 1 \right)$, that  makes $\Pi_t = \Omega_t$ 
for all $t\geq 0$. 



Next, we use the $\{\Omega_t\}$ process to introduce a state variable $\{Y_t; t\geq 0\}$.


\subsection{A State Variable $Y_t$} \label{sec:Y}


Inspired by the martingale property of the $\{q_{0, t}\Pi_t+q_{0, t}\Omega_t\}$ process, 
we define %Next, we introduce $X_t$  which is $\Pi_t$ multiplied by the household's marginal utility of consumption $U_{C, t}$: %multiplied by the cumulative primary deficits 
\begin{align}\label{Y}
Y_t = \Omega_t U_{C, t}\, ,
\end{align}
%In effect, 
%We later show that $X_t$ is the 
which will serve as a state variable for a recursive formulation of the Ramsey problem. Evidently,  $Y_t$ incorporates information about both  $\Omega_t$ and the marginal utility of consumption  $U_{C, t}$. 
%The reason that $X_t$ allows us to recursify our problem is can be viewed as the new worth of cumulative primary deficit measured by the marginal utility up to time $t$.  
%To highlight the mechanism, 

Combining \eqref{Y} and \eqref{foch2} yields
\begin{equation}
Y_t =\Omega_t e^{\rho t}q_{0, t}U_{C, 0}\,. %= e^{\rho t} U_{C, 0}\, \mathbb{E}_t \int_{t}^\infty\, {q_{0, u}} \left(b_{0-, u} +G_u-\tau_{0, u} N_{0, u} \right)du \,.
\end{equation}
Differentiating both sides of the above equation and using $C_t = N_t - G_t$, $e^{\rho t} U_{C, 0} q_{0, t}=U_{C, t}$, and $(1-\tau_{0, t}) U_{C, t}  = -  U_{N, t}$, as implied by \eqref{resourceconstraint} and \eqref{foch1}-\eqref{foch2}, we obtain
\begin{eqnarray}\label{dY} 
dY_t  %&= & \rho  Y_{t-} dt +  e^{\rho t} U_{C, 0} q_{0, t} \left(b_{0-, t} +G_t-\tau_{0, t} N_{0, t} \right)dt   \notag\\
%&= & \rho  X_t dt +  U_{C, t} \left(b_{0-, t} +G_t - N_{0,t} + (1-\tau_{0, t}) N_{0, t} \right)dt   \notag\\
& = & \left(\rho Y_{t-}+b_{0-, t-}U_{C, t-}-C_{0, t-}U_{C, t-} -N_{0, t-}U_{N, t-}  \right)dt\,  \notag \\
& & \quad + \sum_{j\neq s_{t-}} \beta_{t-}Y_{t-} 
\big(d {\mathcal J}_t^{(s_{t-}, j)} 
- \lambda_{s_{t-}, j} e^{\eta(s_{t-}, j)} dt \big) \,,
\end{eqnarray}
where 
\begin{equation} \label{beta}
\beta_{t-}=\beta(s_{t-}, j) = \big[1+\alpha(s_{t-}, j)\big]e^{\eta(s_{t-}, j)}-1\,. 
\end{equation}


  


We have established: 
\begin{proposition} 
	%	Regardless whether a jump arrives or not at $t$, 
	For a Ramsey plan that commits to $\{\alpha_t; t\geq 0 \}$, 
	$Y_t$ defined in (\ref{Y})  evolves according to  (\ref{dY}).
\end{proposition}
%Because  welfare maximization calls for  continuity of the household's net worth $\Pi_t$ multiplied by the marginal utility $U_{C, t}$, 
%That is, the Ramsey planner wants to  make $X_t$ continuous,  even where  the initial debt structure $\{b_{0-,t}\}$ or  the government spending process $\{G_t\}$ jumps. 

Note that, as we discussed earlier, if and only if $\alpha(s_{t-}, j)= e^{-\eta(s_{t-}, j)}-1 $ for all $t$, the promised liability process $\{\Omega_t\}$ coincides with the cumulative liability process $\{\Pi_t\}$ at all $t$, which in turn implies that the $\{Y_t;\, t\geq 0\}$ process defined in \eqref{Y} coincides with the $\{X_t; t\geq 0\}$ process defined in \eqref{X}. 




 
In section \ref{sec:DPdefine1}, we cast  the Ramsey problem  in terms of  the $\{Y_t; t\geq 0\}$ process in   two steps: 1) for a given $Y_0$, we define and solve  a dynamic programming (DP) problem with $Y_t$ and $s_t$ as state variables; 2) we characterize the Ramsey plan. %solve for $X_0$ and $C_0$ 






\subsection{Dynamic Program (DP) and Ramsey Plan} \label{sec:DPdefine1}


\begin{definition} [{\bf  DP Problem}] \label{definDP1}
	Confronting  $Y_0$ at $t=0$, 
	%\paragraph{Dynamic programming for the continuation planner}  at time $0$ confronts $X_0$ as a state variable and
	the government  solves  \eqref{continuation_obj}
	%. Then the continuation Ramsey planer chooses $\{C_{0, t}; t\geq \varepsilon\}$ to optimize
	%	\begin{equation}\label{continuation_obj1}
	%	\max_{\{C_{0, t}; t\geq 0 \}}\; 
	%	\mathbb{E}_0 \,\bigg[\int_{0}^{\infty} e^{-\rho t} U({C}_{0, t}, {C}_{0, t}+G_t) d t  \bigg],
	%	\end{equation}
	subject to the law of motion (\ref{dY}). % of   state variable $X_t$. 
\end{definition} 

%As we have effectively coded the history-relevant information via $Y_t$ given in \eqref{Y}, we next use dynamic programming to formulate and solve this problem. %(see e.g., Fleming and Soner, 2006).
% casts in terms of  a family of optimization problems. 
%Starting from   states vector $(X_t, t)$ at $t\in[0, \infty)$, 
Let $K(Y_0, s_0)$ denote the (optimal) value function for this problem. To solve $K(Y_0, s_0)$, we first construct the following time$-t$ optimization problem: %value function 
%\NW{\begin{equation}
%	e^{-\rho t}  {V}(X_t, t) = \max_{\{C_{0, s}; s\geq t \}}\; 
%	\int_t^{\infty} e^{-\rho s} U({C}_{0, s}, {C}_{0, s}+G_s) d s .
%\end{equation}}
\begin{equation}\label{Vcontinue1}
%{V}(X_t, t) = 
\max_{\{C_{0, u}, \, \beta_u; u\geq t \}}\; 
\mathbb{E}_t	\, \bigg[\int_t^{\infty} e^{-\rho(u-t)} U({C}_{0, u}, {C}_{0, u}+G_u) d u \bigg] \,  
\end{equation}
subject to (\ref{dY}).
Let $K(Y_t, s_t)$ denote the value function where $(Y_t, s_t)$ are the state variables. 
%
%Next, we construct a martingale that will help us to derive the process defined by
%\begin{equation*}
%	\int_{0}^t e^{-\rho t} U(C_{0, t}, C_{0, t}+G_t)dt + e^{\rho t}V(X_t, s_t) 
%\end{equation*}
%is a martingale under physical measure so its drift is zero
%\begin{equation} \label{martingaleV}
%	\mathbb{E}_t \left[d \left(e^{\rho t}V(X_t, s_t)  \right) \right] + e^{-\rho t} U(C_{0, t}, C_{0, t}+G_t)dt =0
%\end{equation}
%Simplifying \eqref{martingaleV} for all $t\geq 0$, 
{In Appendix \ref{appendix:A}}, we show that $K(Y_t, s_t)$ satisfies the following HJB equation: 
%\begin{eqnarray}\label{HJBV} \nonumber
%\rho V (X, s) & = &  \max_{N\in[G, 1], \beta} \quad U(N-G, 1-N) \\
%& & \quad  + \frac{\partial V}{\partial X} \left(\rho X+N U_{L}(N-G, 1-N)-C U_{C}(N-G, 1-N) \WJ{- \sum_{k\neq s}\lambda_{s k}\beta(s, k)}\right)  \nonumber \\
%&& \quad +\sum_{k\neq s}\lambda_{sk}\left[ \WJ{V(X+\beta(s, k), k)}-V(X, s) \right],
%\end{eqnarray}
%\begin{eqnarray}\label{HJBV} \nonumber
%\rho V (X, s) & = &  \max_{N\in[G, 1], \beta} \quad U(N-G, 1-N) \\
%& & \quad\quad \quad \quad+ \frac{\partial V}{\partial X} \left(\rho X+N U_{L}(N-G, 1-N)-C U_{C}(N-G, 1-N) \right)  \nonumber \\
%&& \quad \quad \WJ{- \frac{\partial V}{\partial X} \left(\sum_{k\neq s}\lambda_{s k}\beta(s, k) \right) }+\sum_{k\neq s}\lambda_{sk}\left[ \WJ{V(X+\beta(s, k), k)}-V(X, s) \right],
%\end{eqnarray}
\begin{align}\label{HJBV01} 
\rho K(Y_t, s_t) = &   \max_{C_{0, t} ,  \, \beta_t} \;\; U(C_{0, t}, C_{0, t}+G_t) 
+ \frac{\partial K }{\partial Y_t} \left[\rho Y_t- C_{0, t} U_{C, t}-(C_{0, t}+G_t) U_{N, t}+b_{0-, t} U_{C, t} \right]   \notag \\
& \quad\quad \;\;\,\, +\sum_{j\neq s_t} \lambda_{s_t, j}\left[K(Y+\beta_t Y, j)-K(Y, s_t) -\beta_t\frac{\partial K }{\partial Y_t} \right]	 \, . 
%	&& \quad   \;\;\; %\WJ{- \frac{\partial V}{\partial X} \left(\sum_{k\neq s}\lambda_{s k}\beta(s, k) \right) }
%& \quad\quad \,\,	+\WJ{\sum_{k\neq s}\lambda_{s_tk}\left[ V(X+\beta{(s_t, k)}X, k)-V(X, s_t) - \frac{\partial V(X, s_t)}{\partial X} \beta{(s_t, k)}X   \right] }.
\end{align} 
%
%
We guess and  verify that $K(Y_t, s_t)$ takes the following  form:
\begin{eqnarray}\label{Vadditive01} 
K(Y_t, s_t)=-\Psi_0 Y_t+H(s_t; \Psi_0)\, ,
\end{eqnarray}
where $\Psi_0$ is an object to be chosen by the Ramsey planner and $H(s_t; \Psi_0)$ is given by \eqref{HJBH00}. %Similar to $V(X_t, s_t)$ given in \eqref{Vadditive0}, $K(Y_t, s_t)$ also takes an additively separable form.
Note that the second term in \eqref{Vadditive01} has the same functional form as the second term in \eqref{Vadditive0}: 
$H(s_t; \Psi_0)$ in \eqref{Vadditive01} is also characterized by Proposition \ref{proH}.


The $\{\beta_t; t\geq 0\}$ process  associated with the $\{Y_t; t\geq 0\}$ process given in \eqref{Y} is indeterminate, a consequence of  the additively separable form \eqref{Vadditive01} of the value function $K(Y_t, s_t)$. 

Replacing $\Phi_0$ in the optimal consumption and labor supply expressions \eqref{optC00} and \eqref{optN00} with $\Psi_0$, we obtain the optimal consumption $C(s_t; \Psi_0)$ and labor supply $N(s_t; \Psi_0)$ for this alternative recursive formulation, paralleling the analysis in subsection \ref{sec:DPdefine}. 

In general, for a competitive equilibrium, $\Psi_0$ in \eqref{Vadditive01} does not equal $\Phi_0$ in \eqref{Vadditive0}. But under the Ramsey plan, $\Psi_0^\ast=\Phi_0^\ast$, which we  show below. 

The implementability constraint \eqref{impletconstraint} requires that $\Psi_0$ satisfy $I_0(s_0; \Psi_{0}) =0$, where $I_0( \cdot\,; \,\cdot )$ is defined in \eqref{eq:Psi0}. Thus $\Psi_0$ has the same admissible set as $\Phi_0$, i.e., $\Psi_0\in\mathcal A$, where $\mathcal A$ is defined in subsection \ref{sec:DPRamsey}.  Using (\ref{Vadditive01}) and $Y_0=0$, the Ramsey planner chooses $\Psi_{0}$ by solving:
\begin{eqnarray}\label{MMM} 
M(s_0)  & = & \max_{\Psi_0\in \mathcal A}\, K(0, s_0; \Psi_0) \;\;
= \;\; \max_{\Psi_0\in \mathcal A} \, H(s_0; \Psi_0)  \,.
\end{eqnarray}

% \begin{eqnarray}\label{M}
% M(s_0) %(\{b_{0, t}; t\geq 0\})
% \; = \; \max_{\Psi_{0}} \; K(Y_{0}, s_0; \Psi_0)  \, ,
% \end{eqnarray}
% subject to the implementability constraint (\ref{impletconstraint}).
% 
%
% Using (\ref{Vadditive01}) and $Y_0=0$, the  Ramsey problem  is %equivalent to
% \begin{eqnarray}\label{MMM} 
% M(s_0)  & = & \max_{\Psi_0\in \mathcal A}\, K(0, s_0) \;\;
% = \;\; \max_{\Psi_0\in \mathcal A} \, H(s_0; \Phi_0)  \,.
% \end{eqnarray}

Because both problems reduce to maximizing $H(s_0;\cdot)$ over the same admissible set $\mathcal{A}$, problem \eqref{MMM} coincides with problem \eqref{WWW}. It follows that $\Psi_0^*=\Phi_0^*$, so the optimal consumption policy $C(s_t; \Psi_0^*)$ coincides with $C(s_t; \Phi_0^*)$ for all $s_t$ and $\{C(s_t; \Psi_0^*); t\geq 0\} = \{C(s_t; \Phi_0^*); t\geq 0\}$. Therefore, we have established: 
\begin{proposition} %$\Psi_0^*$ for  equals $\Phi_0^*$ for
	When a Ramsey plan exists, the problem \eqref{MMM} is equivalent to the  problem \eqref{WWW}, and $\Psi_0^* = \Phi_0^*$ and $M(s_0)=W(s_0)$ hold. 
\end{proposition}








\section{Concluding Remarks}\label{sec:conclusion}

This paper extended the deterministic continuous-time formulation of the Lucas-Stokey optimal tax problem in \cite{JiangSargentWang2024} to a stochastic environment driven by a Markov chain.  Several lessons emerged.

\paragraph{Instantaneous debt as a sufficient instrument.}
A central message of both the deterministic predecessor and this stochastic extension is that the Ramsey plan can be implemented without rescheduling the initial term structure of government debt.  It suffices for the government to trade locally risk-free instantaneous debt --- together with, in the stochastic setting, competitively priced hedging contracts --- while leaving the inherited coupon stream $\{b_{0-,t}\}$ untouched.  This contrasts with \cite{LucasStokey1983}, who implemented the Ramsey plan by restructuring the entire term structure period by period.

\paragraph{Complete markets in continuous time.}
Translating \citeauthor{LucasStokey1983}'s (1983) complete-markets assumption to continuous time required care. The dynamic spanning constructions pioneered by \cite{BlackScholes1973}, \cite{merton1973theory}, and \cite{HarrisonandKreps1979} show that instantaneous debt together with $n-1$ hedging contracts that pay off at Markov state transitions suffice to span all contingencies over each infinitesimal interval.  These hedging contracts allow the government to offset jumps in bond prices and in its cumulative liability at state transitions, maintaining dynamic completeness without access to a  full set of dated Arrow securities. \NW{$n-1$ is not defined. to discuss.} 

\paragraph{A natural state variable.}
The cumulative government liability $\Pi_t$ --- the running present-value sum of primary deficits and inherited debt service --- can jump at state transitions because bond prices and tax revenues jump.  Yet its purchasing-power-adjusted counterpart $X_t = \Pi_t\, U_{C,t}$ evolves continuously in time.  This continuity makes $X_t$ a tractable state variable and underlies the additively separable structure $V(X_t,s_t) = -\Phi_0\, X_t + H(s_t; \Phi_0)$ of the Ramsey planner's value function.  The separability in turn enables us to reduce the planner's infinite-dimensional problem to a one-dimensional search  over the scalar $\Phi_0$.

\paragraph{Local commitment and sequential implementation.}
Our timing protocol grants each government only ``local commitment'': its current-period tax rate is set by its immediate predecessor.  Under this protocol, the continuation government at date~$t$ inherits instantaneous debt $B_t = \Pi_t^\ast$ and the tail of the original term debt, and faces the same purchasing-power-adjusted state $X_t^\ast$ as the date-0 Ramsey planner.  Because the value function is linear in $X_t$, the continuation government's incentives are aligned with those of its predecessor, and it voluntarily confirms the Ramsey plan.  This structure helps us understand and address the \cite*{debortoli2021optimal} counterexample.

\paragraph{Tax smoothing and hedging.}
The quantitative illustrations in section~7 revealed a sharp dichotomy.  When only the initial term debt  switches between regimes, the Ramsey planner manipulates interest rates and bond prices to reduce the  market value of outstanding obligations,  allowing taxes  and consumption to vary across states.  When instead only government spending switches while the initial term debt is constant, the planner achieves perfect tax smoothing: consumption, tax rates, and interest rates all remain constant across states, funded entirely by adjusting instantaneous debt and hedging positions.  The latter result extends to continuous time \citeauthor{barro1979determination}'s (1979) insight that optimal tax rates need not inherit the serial correlation properties of the government expenditure process.

\paragraph{Uniqueness despite indeterminacy.}
Section~\ref{sec:indeterminacy} showed that an alternative recursive formulation built on a forward-looking  promised liability $\Omega_t$ leads to indeterminacy in the portfolio process --- the government has freedom  in how it distributes  obligations  across hedging contracts.  Nevertheless, the Ramsey plan itself is unique: the optimal multiplier $\Psi_0^\ast$ equals $\Phi_0^\ast$, and consumption-tax paths coincide with those from the backward-looking formulation in section~\ref{sec:bellmanramsey}.  Indeterminacy is  confined to the  portfolio  that finances a unique allocation. 

%\newpage
\bigskip
%\footnotesize
%\begin{center}
%\textbf{References}
%\end{center}
%	\singlespacing
%\setlength{\bibsep}{0.3ex}
\bibliographystyle{chicago}
\bibliography{LSreference}




\clearpage




\pagebreak
\section*{Appendices}
%\appendix{\bf \centering \LARGE \noindent Appendices}
\appendix{}
%\setcounter{page}{1} 
%\setcounter{Theorem}{0} 
%\renewcommand{\thepage}{I-\arabic{page}}
%\setcounter{theorem}{0}
%\renewcommand{\thetheorem}{\Alph{section}.\arabic{theorem}}
\setcounter{equation}{0}
\renewcommand{\theequation}{A-\arabic{equation}}
\setcounter{section}{0}
\renewcommand{\thesection}{\Alph{section}}
\renewcommand{\thesubsection}{\Alph{section}.\arabic{subsection}}
\setcounter{figure}{0}
\renewcommand{\thefigure}{A-\arabic{figure}}
\setcounter{footnote}{0}
\renewcommand{\thefootnote}{A-\arabic{footnote}}
\bigskip%


\section{Technical Details} \label{appendix:A}

\medskip
\noindent\textbf{Regularity Conditions for Proposition~\ref{proplag}.}\quad We provide two regularity conditions that guarantee that a Ramsey plan exists. First, for a given initial debt structure, the admissible allocation set is not empty, so there exist allocations $\{C_{0,t};\, t \ge 0\}$ that satisfy the implementability condition~\eqref{impletconstraint}. Second, $(CU_C + (C+G)U_N - bU_C)$ is concave for $G > 0$, $b \ge 0$.

%For economies analyzed in Debortoli, Nunes, and Yared (2021) and this paper, these two regularity conditions are satisfied when~$b_0$ is not too high.

\bigskip

\noindent\textbf{Proof of Lemma \ref{lemmaPi0}: Dynamics of the $\{\Pi_t\}$ Process.}\quad 
Below, we prove Lemma \ref{lemmaPi0} by deriving the dynamics of $\Pi_t$ defined in \eqref{Pi}. 
Using the property that $\{X_t\}$ given in \eqref{dX} is continuous in $t$ and combining \eqref{X} with \eqref{foch2}, we obtain
\begin{equation}\label{Pi2}
\Pi_t=q_{0, t}^{-1} \int_{0}^t\, {q_{0, u}} \left(b_{0-, u} +G_u-\tau_{0, u} N_{0, u} \right)du\,.
\end{equation}
Given a consumption process $\{C_{0, t}; t\geq 0\}$, the SDF can be written  as  (see, e.g., \cite{duffie2010dynamic}): %requires that  
\begin{eqnarray}\label{sdfq2}
d q_{0, t} = - r^{}_{0, t-} q_{0, t-}dt  -  q_{0, t-}  \sum_{j\neq s_{t-}} \zeta(s_{t-}, j)\left(d {\mathcal J}_t^{(s_{t-}, j)} - \lambda_{s_{t-}, j} d t \right)\,,
\end{eqnarray}
where $ \zeta(s_{t-}, j)$ is given by
\begin{equation}
\zeta(i, j) = -\big(e^{\eta(i, j)}-1 \big)\,,
\end{equation}
and {$\eta(i,j)$ for  $i\neq j$} is given by: %equals logarithmic %given by %the relative jump size which satisfies 
\begin{equation}\label{kappa} 
\eta(i, j)= \ln\frac{U_C(C(j), N(j))}{U_C(C(i), N(i))} \end{equation}

Applying Ito's lemma to \eqref{Pi2} and using \eqref{dX} and \eqref{sdfq2}, we obtain
\begin{align}
d\Pi_t  = &  q_{0, t-}^{-1} d \left[ \int_{0}^t\, {q_{0, u}} \left(b_{0-, u} +G_u-\tau_{0, u} N_{0, u} \right)du \right] +   \left[\int_{0}^{t-}\, {q_{0, u}} \left(b_{0-, u} +G_u-\tau_{0, u} N_{0, u} \right)du \right] d\left( q_{0, t}^{-1}\right)   \notag\\
 & \quad  + (\Pi_t - \Pi_{t-}) \sum_{j\neq s_{t-}} d {\mathcal J}_t^{(s_{t-}, j)}  \notag\\
 = & \left(b_{0-, t-} +G_{t-}-\tau_{0, t-} N_{0, t-} \right)dt    \notag \\
& -  \left[\int_{0}^{t-}\, {q_{0, u}} \left(b_{0-, u} +G_u-\tau_{0, u} N_{0, u} \right)du \right] q_{0, t-}^{-2} \big(- r^{}_{0, t-} q_{0, t-}dt  +  q_{0, t-}  \sum_{j\neq s_{t-}} \zeta(s_{t-}, j) \lambda_{s_{t-}, j} d t \big) \notag\\
&   \quad  + (\Pi_t - \Pi_{t-}) \sum_{j\neq s_{t-}} d {\mathcal J}_t^{(s_{t-}, j)}  \notag\\
=&   \left(b_{0-, t-} +G_{t-}-\tau_{0, t-} N_{0, t-} \right)dt  - \Pi_{t-} \left(- r^{}_{0, t-} dt  +    \sum_{j\neq s_{t-}} \zeta(s_{t-}, j) \lambda_{s_{t-}, j} d t  \right)  \notag \\
 & \quad  + (\Pi_t - \Pi_{t-}) \sum_{j\neq s_{t-}} d {\mathcal J}_t^{(s_{t-}, j)}  \notag\\
 = & \left(b_{0-, t-} +G_{t-}-\tau_{0, t-} N_{0, t-} \right)dt + r^{}_{0, t-}\Pi_{t-}dt \notag\\
 & \quad + \Pi_{t-} \sum_{j\neq s_{t-}} \left[ \left( \frac{\Pi_{t}}{\Pi_{t-}}-1  \right) d {\mathcal J}_t^{(s_{t-}, j)} -\zeta(s_{t-}, j) \lambda_{s_{t-}, j} d t   \right]  \notag \\
 =&\left(b_{0-, t-} +G_{t-}-\tau_{0, t-} N_{0, t-} \right)dt + r^{}_{0, t-}\Pi_{t-}dt  \notag \\
 &\quad   +    \sum_{j\neq s_{t-}} \left(e^{-\eta(s_{t-}, j)}-1 \right)\Pi_{t-} \left(d {\mathcal J}_t^{(s_{t-}, j)} - \lambda_{s_{t-}, j}e^{\eta(s_{t-}, j)} d t \right)
\end{align}
where the last equality holds because that $\left( \frac{\Pi_{t}}{\Pi_{t-}}-1  \right) = e^{-\eta(s_{t-}, s_t)}-1$ and $\zeta(s_{t-}, s_t)= -\left(e^{\eta(s_{t-}, s_t)}-1\right) =  e^{\eta(s_{t-}, s_t)} \left(e^{-\eta(s_{t-}, s_t)}-1\right) $.

\bigskip
\noindent\textbf{Derivation of HJB Equation \eqref{HJBV0}.}\quad 
Next, we construct a martingale that helps characterize the process
\begin{equation*}
    \int_0^t e^{-\rho u} U(C_{0,u}, C_{0,u}+G_u)\, du
    + e^{-\rho t} V(X_t, s_t).
\end{equation*}
Under the physical measure, this process is a martingale, and hence its drift must be zero. Therefore,
\begin{equation} \label{martingaleV}
    \mathbb{E}_t\left[d\left(e^{-\rho t}V(X_t, s_t)\right)\right]
    + e^{-\rho t} U(C_{0,t}, C_{0,t}+G_t)\,dt = 0.
\end{equation}
Applying Ito's lemma to the left hand side of \eqref{martingaleV}, we obtain the HJB equation \eqref{HJBV0}.

\bigskip
\noindent\textbf{Proof of Proposition \ref{proH}} \quad
Substituting optimal consumption defined in \eqref{optC00} into HJB equation \eqref{HJBH00}, we oabtain for $s\in (1, 2, \cdots, n)$
\begin{equation}
\rho H(s; \Phi_0) = A(s; \Phi_0) + \sum_{k\neq s} \lambda_{s, k}\left[H(k; \Phi_0) - H(s; \Phi_0) \right] \,.
\end{equation}
We then obtain
\begin{equation}
\left ( \rho \mathbb{I} -{\Lambda} \right) \mathbb{H} =\mathbb{A} \,,
\end{equation}
which gives \eqref{HH}. 



%\medskip
%\paragraph{Technical Details for Proposition \ref{propimplement}} \quad




\end{document}

